\documentclass{article}

% 如果你想尽量保持 ACM 那种数字压缩引用风格，保留这一行
\PassOptionsToPackage{numbers,compress}{natbib}

% Use the official style when available; fall back for local single-file compilation.
\IfFileExists{neurips_2026.sty}{%
  \usepackage{neurips_2026}%
}{%
  \usepackage[margin=1in]{geometry}%
  \usepackage{natbib}%
}
% Final style setting, if required by the venue:
% \usepackage[main, final]{neurips_2026}

\usepackage[utf8]{inputenc}
\usepackage[T1]{fontenc}
\usepackage[hypertexnames=false]{hyperref}
\hypersetup{hidelinks,breaklinks=true,pdfborder={0 0 0}}
\usepackage{url}
\usepackage{booktabs}
\usepackage{amsmath,amssymb,amsfonts}
\usepackage{nicefrac}
\usepackage{microtype}
\usepackage{xcolor}
\usepackage{graphicx}
\usepackage{adjustbox}

\usepackage{algorithm,algorithmic}
\usepackage{placeins}
\usepackage{tabularx}
\usepackage{multirow}
\usepackage{array}
\usepackage{listings}
\newcolumntype{Y}{>{\raggedright\arraybackslash}X}
\lstset{
  basicstyle=\ttfamily\footnotesize,
  breaklines=true,
  columns=fullflexible,
  keepspaces=true,
  showstringspaces=false
}
\newcommand{\ECPO}{\textsc{ECPO}}
\newcommand{\CertNDCG}{\textsc{CertNDCG}}
\newcommand{\EvidCons}{\textsc{EvidCons}}
\newcommand{\Gdet}{G_{\mathrm{det}}}
%%
%% end of the preamble, start of the body of the document source.

%%
%% The "title" command has an optional parameter,
%% allowing the author to define a "short title" to be used in page headers.
\title{\ECPO: Evidence-Coupled Policy Optimization for Evidence-Certified Candidate Ranking}

\author{
  Anonymous Authors \\
  Anonymous Institution \\
  \texttt{anonymous@example.com}
}
%%
%% By default, the full list of authors will be used in the page
%% headers. Often, this list is too long, and will overlap
%% other information printed in the page headers. This command allows
%% the author to define a more concise list
%% of authors' names for this purpose.

\begin{document}
\maketitle

\begin{abstract}
Ranking systems used in decision-support settings should not only order candidates but also expose evidence that can be independently checked. We study \emph{evidence-certified candidate ranking}: given an \texttt{intent\_id}, a predefined plan skeleton, a window-local candidate roster, and text-derived candidate trajectories with span provenance, a system must output a Top-$K$ list together with \texttt{doc\_id:span} evidence certificates whose cited spans are sufficient to recover the decision. We instantiate this task on MAVEN-ERE and RAMS with fixed upstream extraction, window-local randomized candidate identifiers, skeleton-aligned trajectory supervision, hard negatives, and audit references. We introduce \textbf{Evidence-Coupled Policy Optimization} (\ECPO), a listwise policy-optimization objective whose action is the joint object $(\text{ranking},\text{evidence certificate})$. \ECPO{} first learns an interpretable trajectory reward from skeleton alignment, argument consistency, and optional graph features; it then optimizes a constrained policy with three coupled rewards: listwise ranking utility, span-level certificate validity, and an evidence-cycle reward computed by a label-free deterministic verifier that reconstructs candidate support from claim-stripped cited spans. This reframes the goal from maximizing ordinary NDCG alone to maximizing \CertNDCG{} and decision--evidence coupling. The evaluation compares \ECPO{} against zero-shot, SFT, and GRPO policies, RM-only scoring with deterministic evidence attachment, grammar/JSON-constrained decoding, validator retry, best-of-$N$ RM selection, and post-hoc evidence rationalization under closed-roster, predicted-roster, and hybrid-roster settings.
\end{abstract}

\section{Introduction}
\label{sec:introduction}

Decision-support systems increasingly rank candidates from sparse textual traces: analysts must decide which entities, cases, or hypotheses deserve attention under limited time, noisy extraction, and incomplete observations. In such settings, a ranking score alone is rarely acceptable. A user must be able to inspect why a candidate was prioritized, trace the support back to source text, and reject outputs whose evidence does not actually justify the decision. This paper studies the resulting problem of \emph{evidence-certified candidate ranking}: the model must produce both a Top-$K$ candidate list and a machine-checkable evidence certificate for each returned candidate.

The input to one ranking decision is a window $w$ containing an \texttt{intent\_id}, a predefined plan skeleton $S_w$, an evidence scope $\mathcal{D}_w$, a window-local candidate roster $C_w$, and one text-derived trajectory $\tau_w(c)$ for every candidate $c\in C_w$. Each trajectory is produced by a frozen upstream event extractor and preserves trigger and argument provenance as \texttt{doc\_id:span} offsets. The output is not a free-form explanation. It is a strict JSON object with a ranked list $A_w=[a_1,\ldots,a_{K_w}]$, $K_w=\min(K,|C_w|)$, and a step-aligned evidence certificate: for every returned candidate and every skeleton step, the model either cites traceable spans or explicitly marks the step as unmatched. Candidate identifiers exposed to the policy are randomized within each window, so ranking decisions cannot rely on source-level identifier reuse; source IDs are retained only for internal bookkeeping and evaluation.

This formulation differs from ordinary learning-to-rank, event extraction, and post-hoc explanation generation. First, the decision is comparative: the system must order candidates within the same evidence window, not simply classify isolated events. Second, the explanation is operational rather than narrative: every cited span must be inside the exposed documents, linked to the candidate trajectory, and compatible with the corresponding skeleton step. Third, the evidence must be decision-sufficient. A post-hoc rationalizer can attach plausible spans after a ranker has already chosen candidates, but this does not guarantee that the cited evidence alone would recover the same Top-$K$ set. We therefore evaluate not only NDCG@10 and MAP@10 but also \CertNDCG@10 and \EvidCons@10, where ranking credit is conditioned on certificate validity and evidence-only decision recovery.

We instantiate the benchmark on MAVEN-ERE~\citep{wang_maven-ere_2022} and RAMS~\citep{ebner_multi-sentence_2020}. The task is primarily a ranking-and-certification benchmark: event extraction is fixed upstream preprocessing, while all ranking methods consume the same predicted trajectories. The main closed-roster setting evaluates decision support when candidate availability is known before ranking. To avoid hiding candidate-generation assumptions, we also specify predicted-roster and hybrid-roster settings that quantify how upstream candidate recall affects certified ranking. Appendix~\ref{sec:app_moved_intro_details}--\ref{sec:app_data_windowing} gives the full construction protocol, schemas, leakage controls, and artifact layout.

We introduce \textbf{Evidence-Coupled Policy Optimization} (\ECPO). The central design choice is to treat the policy action as a joint object $(A,E)$: a candidate ranking $A$ and an evidence certificate $E$ that must support $A$. \ECPO{} first learns an interpretable skeleton-conditioned trajectory reward $R_\theta(\tau;S_w)$ using dynamic-programming alignment, argument consistency, and optional within-window graph features. It then optimizes a constrained generative policy with a shaped reward consisting of three terms: (i) listwise ranking utility from $R_\theta$, (ii) certificate validity under deterministic span and traceability checks, and (iii) an evidence-cycle reward that passes the cited spans to a label-free deterministic evidence-only verifier and rewards agreement between the verifier's recovered candidates and the policy's own Top-$K$ list. The learning objective therefore directly targets decision--evidence coupling instead of treating certificates as post-hoc text.

Our goal is not to claim that a generative policy is always superior to a deterministic reward-model ranker. RM-only scoring with deterministic alignment is cheap, stable, and often strong when the user only needs a ranked list or when post-hoc evidence attachment is acceptable. The intended contribution is narrower and more operational: when the deployed interface requires a single policy to emit schema-valid rankings and evidence certificates, ECPO makes the certificate part of the optimization target. This distinction also dictates our experimental design. In addition to standard LLM policies, we compare against strong decoding-only and post-hoc alternatives: grammar/JSON-constrained decoding, validator retry, best-of-$N$ RM selection, RM-only plus deterministic evidence search, and RM-only plus LLM evidence rationalization.

Our contributions are:
\begin{itemize}
    \item We formulate \emph{evidence-certified candidate ranking}, where Top-$K$ candidate decisions are scored jointly with span-grounded evidence certificates and evidence-only decision recovery.
    \item We propose \ECPO, an evidence-coupled listwise policy-optimization objective that combines skeleton-conditioned trajectory rewards with deterministic certificate validation and an evidence-cycle reward.
    \item We provide MAVEN-ERE and RAMS benchmark instantiations with randomized window-local candidate IDs, closed/predicted/hybrid roster settings, hard-negative stress suites, preference pairs, and audit metrics designed to distinguish true decision--evidence coupling from post-hoc evidence rationalization or decoding-only fixes.
\end{itemize}

\section{Related Work}
\label{sec:related-work}

\paragraph{Evidence-grounded ranking and learning to rank.}
Learning-to-rank methods such as LambdaMART~\citep{burges_ranknet_2010} optimize comparative ordering, and reward-learning methods such as maximum-entropy IRL~\citep{ziebart_maximum_2008} provide a way to learn trajectory utilities from demonstrations or preferences. These paradigms usually evaluate the ranked list itself. ECPO instead evaluates a joint output: a ranked candidate list plus a span-grounded certificate whose evidence must be valid, traceable, and decision-sufficient.

\paragraph{Event extraction and text-derived trajectories.}
MAVEN-ERE~\citep{wang_maven-ere_2022} extends large-scale event annotation from MAVEN~\citep{wang_maven_2020}, while RAMS~\citep{ebner_multi-sentence_2020} provides multi-sentence argument supervision. These resources support extraction and relation modeling, but they do not directly define window-local candidate rosters, Top-$K$ comparative decisions, or certificate objects that can be deterministically checked against source spans. Our benchmark freezes extraction and evaluates the downstream ranking-and-certification interface.

\paragraph{Plan recognition under partial observability.}
Classical plan and goal recognition infer goals or plans from partial observations~\citep{ramirez_plan_2009,ramirez_probabilistic_2010}, and process-conformance work compares event traces with explicit process models~\citep{brockhoff_process_2021,park_realizing_2021}. ECPO borrows the idea of comparing observations to a skeleton, but changes the output object: we rank multiple candidates in the same evidence window and require each ranked decision to be supported by \texttt{doc\_id:span} evidence.

\paragraph{Certificates, constrained generation, and auditability.}
Free-form explanations can sound plausible without being faithful to the actual decision, while constrained decoding can improve syntax without ensuring that cited evidence explains why a candidate was selected. ECPO treats schema validity, candidate-ID validity, span traceability, and evidence-only decision recovery as first-class optimization and evaluation targets. This is why our baselines include not only ordinary LLM ranking policies but also grammar/JSON-constrained decoding, validator retry, best-of-$N$ RM selection, and deterministic evidence attachment. Extended discussion is in Appendix~\ref{sec:app_related_work}.

\section{Method: Evidence-Coupled Policy Optimization}
\label{sec:method}

ECPO learns a policy that produces evidence-certified rankings. The policy input is a compact serialization of a ranking window; the policy output is a strict JSON object containing both a Top-$K$ candidate list and step-aligned evidence. The key difference from a post-hoc evidence framing is that evidence consistency is not only an audit diagnostic. It is part of the training objective through a certificate reward and an evidence-cycle reward.

\paragraph{Pipeline overview.}
For each window, \texttt{intent\_id} selects a predefined skeleton $S_w$. Candidate trajectories are aligned to $S_w$ to produce reward features and candidate evidence spans. A trajectory reward model $R_\theta$ scores plan consistency. A deterministic validator $V$ checks whether a generated certificate is syntactically valid and traceable. A label-free deterministic evidence-only verifier $\Gdet$ receives claim-stripped evidence bundles, scores every candidate in the window roster against those spans, and reconstructs which candidates are best supported without using the policy ranking or gold labels. ECPO optimizes the policy so that its ranking is high-utility, schema-valid, and recoverable from its own cited evidence.

\begin{figure*}[t]
    \centering
    \includegraphics[width=\textwidth]{pic/ecpo-alg-1.png}
    \caption{\textbf{ECPO training loop.}
    The policy samples a joint ranking/certificate object. The validator enforces schema and span traceability, while the deterministic evidence-only verifier reconstructs candidates from claim-stripped cited spans and supplies the evidence-cycle reward.}
    \label{fig:ecpo_alg}
\end{figure*}

\subsection{Evidence-certified ranking problem}
Given a ranking window $w$, the conceptual input is
\begin{equation}
x_w =
\big(
\texttt{window\_id}(w),
\texttt{intent\_id}(w),
S_w,
\mathcal{D}_w,
C_w,
\{\tau_w(c)\}_{c\in C_w}
\big),
\end{equation}
where $S_w$ is a predefined skeleton, $\mathcal{D}_w$ is the exposed evidence scope, and $C_w$ is the window-local candidate roster. The policy sees randomized window-local candidate IDs, while the evaluator maintains a hidden mapping to source-normalized entity clusters. For each candidate $c$, $\tau_w(c)$ is a possibly empty fixed-extractor trajectory with trigger and argument spans.

For target cutoff $K$, the policy emits
\begin{equation}
y=(A_y,E_y),\qquad A_y=[a_1,\ldots,a_{L_y}],\quad a_k\in C_w,
\end{equation}
where $A_y$ is the parsed candidate prefix and $E_y[k]$ is the rank-positioned evidence certificate paired with $a_k$. A certificate contains one serialized object for every skeleton step: either a matched event with one or more \texttt{doc\_id:span} references, or an explicit unmatched step with empty evidence. A valid output must have $L_y=K_w=\min(K,|C_w|)$, no duplicate candidate IDs, no out-of-roster candidates, one certificate per returned rank position, and all cited spans traceable to the exposed trajectory records.

\subsection{Skeleton alignment and trajectory reward}
ECPO retains an interpretable skeleton-conditioned trajectory scorer. Given a candidate trajectory $\tau=[e_1,\ldots,e_T]$ and skeleton $S=([s_1,\ldots,s_M],\prec)$, a dynamic-programming aligner maps each skeleton step to either a type-compatible event or a missing marker. The local compatibility score is
\begin{equation}
m(k,t)=
\begin{cases}
\operatorname{role\_sat}(s_k,e_t)
-\lambda_{\mathrm{arg}}\operatorname{arg\_viol}(s_k,e_t),
& \text{if } \texttt{etype}_k\in \texttt{skeleton\_hits}(e_t),\\[2pt]
-\delta_{\mathrm{type}},& \text{otherwise,}
\end{cases}
\end{equation}
and the recurrence allows matched steps, missing skeleton steps, and skipped observed events:
\begin{equation}
\begin{aligned}
\mathrm{DP}[k,t]=\max\{&\mathrm{DP}[k-1,t-1]+m(k,t),\\
&\mathrm{DP}[k-1,t]-\delta_{\mathrm{miss}},\\
&\mathrm{DP}[k,t-1]-\delta_{\mathrm{skip}}\}.
\end{aligned}
\end{equation}
The backtracked alignment yields step hits, misses, skipped events, temporal gaps, role-consistency indicators, precedence diagnostics, and evidence candidates. These features form
\begin{equation}
\phi(\tau;S_w)=
[\phi_{\mathrm{skel}}(\tau),\phi_{\mathrm{arg}}(\tau),\phi_{\mathrm{graph}}(\tau)],
\end{equation}
where $\phi_{\mathrm{graph}}$ is optional and is set to $\mathbf{0}$ in the no-graph ablation. The trajectory reward is
\begin{equation}
R_\theta(\tau;S_w)=\theta^\top\phi(\tau;S_w).
\end{equation}
We learn $\theta$ from a train-split reward-learning corpus containing positives, pseudo-experts, near-neighbor hard negatives, perturbation negatives, and within-window preference pairs. The objective combines a window-wise MaxEnt likelihood with a Bradley--Terry preference term; full details are in Appendix~\ref{sec:app_irl}.

\subsection{Evidence certificates and deterministic evidence-only reconstruction}
\label{sec:method_det_verifier}
A deterministic validator $V(y,w)$ checks the non-negotiable interface constraints: parseability, schema conformance, candidate membership, deduplication, expected Top-$K_w$ length, document existence, span bounds, event-to-trajectory traceability, and skeleton-step compatibility. This validator catches malformed or unsupported certificates, but it does not by itself answer whether the cited evidence is sufficient to recover the chosen candidates.

For ECPO we therefore use a label-free deterministic evidence-only verifier $\Gdet$. Unlike a learned verifier, $\Gdet$ has no trainable parameters and is not fit to dev/test labels. It receives the generated certificate after stripping candidate claims, rank positions, policy scores, and any field that directly states which candidate the bundle is supposed to support. It also receives only public window metadata needed for validation: the skeleton, exposed document/span metadata, and the frozen candidate trajectories. It does not receive policy logits, the policy ranking score, gold relevance labels, hidden audit references, or source-level stable identifiers.

Let $B_y=\{b_1,\ldots,b_{L_y}\}$ denote the claim-stripped evidence bundles obtained from $E_y$. For each bundle $b$ and every roster candidate $c\in C_w$, $\Gdet$ computes a support score by aligning the cited spans in $b$ against the candidate trajectory $\tau_w(c)$:
\begin{equation}
\begin{aligned}
s_{\mathrm{det}}(c,b,w)=&\ \alpha_{\mathrm{cov}}\operatorname{StepCov}(b,S_w,c)
+\alpha_{\mathrm{role}}\operatorname{RoleSat}(b,S_w,c)\\
&+\alpha_{\mathrm{trace}}\operatorname{Trace}(b,\tau_w(c))
+\alpha_{\mathrm{prec}}\operatorname{Prec}(b,S_w,c)
-\alpha_{\mathrm{bad}}\operatorname{BadSpan}(b,w).
\end{aligned}
\end{equation}
All components are deterministic functions of cited spans, normalized roles, event signatures, and skeleton steps. $\operatorname{StepCov}$ counts skeleton steps whose cited spans can be mapped to compatible events for candidate $c$; $\operatorname{RoleSat}$ checks required role evidence; $\operatorname{Trace}$ checks whether cited spans overlap $\tau_w(c)$ without using the claimed candidate id; $\operatorname{Prec}$ rewards consistency with soft precedence when order is available; and $\operatorname{BadSpan}$ penalizes invalid, duplicate, off-window, or non-traceable spans. The weights and thresholds are fixed before policy training and are not selected using held-out labels.

A bundle is assigned to a candidate only if the best-scoring candidate exceeds a minimum support threshold and beats the second-best candidate by a fixed margin; otherwise the bundle is marked ambiguous and gives no cycle credit. The verifier then performs deterministic maximum-weight bipartite assignment between evidence bundles and roster candidates. For the main reward and metrics, the assignment is read back at the original certificate slot so that the verifier must recover the candidate paired with that slot:
\begin{equation}
\hat a_j=\Gdet(E_y[j],w)=\operatorname{slot}_j\!\left(\operatorname{MWM}_{b\in B_y,c\in C_w}s_{\mathrm{det}}(c,b,w)\right).
\end{equation}
The default evidence-cycle reward is slot-wise recovery:
\begin{equation}
r_{\mathrm{cycle}}(y,w)=\frac{1}{K_w}\sum_{j=1}^{K_w}\mathbb{I}[\hat a_j=a_j].
\end{equation}
For candidate-level certification, the $k$th ranked candidate receives credit only when the evidence bundle paired with rank $k$ is reconstructed as the same candidate after claim stripping. Set-overlap and stricter rank-aware variants are reported only as diagnostics.

This construction distinguishes ECPO from post-hoc evidence rationalization. A post-hoc system may attach valid-looking spans to a ranking selected elsewhere; ECPO receives reward only when the cited evidence itself supports recovering the selected candidates under the same label-free verifier used for certified evaluation.

\subsection{ECPO objective}
For a parsed output $y=(A_y,E_y)$, define listwise utility
\begin{equation}
r_{\mathrm{rank}}(y,w)=
\sum_{k=1}^{L_y}\gamma^{k-1} R_\theta(\tau_w(a_k);S_w),
\end{equation}
where $0<\gamma\le 1$ discounts lower ranks. Because $r_{\mathrm{rank}}$ is a discounted sum while the certificate and cycle terms are bounded rates, we normalize the listwise reward within each training window or sampled group:
\begin{equation}
\bar r_{\mathrm{rank}}(y,w)=
\frac{r_{\mathrm{rank}}(y,w)-\mu_{\mathrm{rank}}(w)}{\sigma_{\mathrm{rank}}(w)+\epsilon_{\mathrm{rank}}}.
\end{equation}
This normalization prevents the ranking term from dominating certificate learning solely because of scale.
Define certificate utility
\begin{equation}
r_{\mathrm{cert}}(y,w)=
\frac{1}{K_w}\sum_{k=1}^{K_w}
\mathbb{I}[V(b_k,a_k,w)=1]\cdot
\operatorname{StepCov}(b_k,S_w),
\end{equation}
where $b_k$ is the rank-positioned certificate bundle aligned to $a_k$ for feasibility validation. Invalid outputs receive zero certificate utility and explicit penalties. The ECPO reward is
\begin{equation}
\begin{aligned}
r_{\mathrm{ECPO}}(y,w)=&\ \bar r_{\mathrm{rank}}(y,w)
+\lambda_{\mathrm{cert}} r_{\mathrm{cert}}(y,w)
+\lambda_{\mathrm{cycle}} r_{\mathrm{cycle}}(y,w)\\
&-\xi_{\mathrm{inv}}\mathbb{I}_{\mathrm{invalid}}(y,w)
-\xi_{\mathrm{miss}}\max(0,K_w-L_y)
-\beta_{\mathrm{KL}}\operatorname{KL}_{\mathrm{tok}}(\pi_\varphi\|\pi_{\mathrm{ref}}).
\end{aligned}
\end{equation}
The policy is optimized with group-normalized policy-gradient updates over multiple samples from the same window. The reward model $R_\theta$ is frozen during each policy phase and $\Gdet$ is a fixed deterministic scoring rule; hard cases mined from policy samples can be used only in train-split refresh phases. At inference time, ECPO decodes once per window unless a decoding-baseline experiment explicitly enables retry or best-of-$N$ selection.

\paragraph{Certified utility.}
For evaluation, we report ordinary ranking metrics and certified variants. For a candidate $a_k$ at rank $k$, let $z_k=1$ if its certificate is valid and the claim-stripped evidence bundle paired with rank $k$ is reconstructed as $a_k$ by $\Gdet$. Certified DCG is
\begin{equation}
\operatorname{CDCG}@K(y,w)=
\sum_{k=1}^{K_w}\frac{(2^{\mathrm{rel}_w(a_k)}-1)z_k}{\log_2(k+1)},
\end{equation}
and \CertNDCG@K normalizes by the ideal ordinary DCG for the window. Thus a method cannot obtain full certified ranking credit by ranking correctly while citing irrelevant, ambiguous, or insufficient evidence.

\paragraph{Sanity property: certified utility is conservative.}
Assume the verifier has one-sided soundness: whenever $z_k=1$, the cited evidence is valid, traceable, and sufficient for recovering candidate $a_k$ under $\Gdet$. Then $\operatorname{CDCG}@K(y,w)\le \operatorname{DCG}@K(y,w)$ for every output, with equality only when every relevant gain in the ranked prefix is certified. This property is a consistency check on the metric rather than the main theoretical contribution: certified utility is ordinary utility masked by evidence recoverability. The proof follows immediately because each ordinary gain term is multiplied by $z_k\in\{0,1\}$.

\begin{algorithm}[t]
\caption{Evidence-Coupled Policy Optimization (ECPO)}
\label{alg:ecpo}
\begin{algorithmic}[1]
\STATE Train or load the frozen event extractor; construct trajectories with \texttt{doc\_id:span} provenance.
\STATE Learn $R_\theta(\tau;S)$ from train-split positives, hard negatives, perturbations, and preferences.
\STATE Instantiate the label-free deterministic evidence-only verifier $\Gdet$ with fixed scoring rules.
\FOR{each policy phase}
    \STATE Sample strict-JSON ranking/evidence outputs $y_i\sim\pi_\varphi(\cdot\mid x_w)$ for each train window.
    \STATE Validate outputs with $V$ and reconstruct candidates from evidence with $\Gdet$.
    \STATE Compute $r_{\mathrm{ECPO}}(y_i,w)$ and group-normalized advantages within each window.
    \STATE Update $\pi_\varphi$ with KL regularization to the reference policy.
\ENDFOR
\STATE Decode once per test window and evaluate ordinary and certified ranking metrics.
\end{algorithmic}
\end{algorithm}

\subsection{Inference and deployment modes}
ECPO is used when the final deployed object must be a single schema-valid ranking with evidence certificates. We also retain cheaper alternatives as baselines and possible deployment modes. RM-only ranks candidates by $R_\theta$ and attaches deterministic evidence; it is appropriate when a deterministic ranker is sufficient. Post-hoc evidence generation first ranks candidates and then asks an LLM or aligner to explain the ranking; it is appropriate only if evidence sufficiency is not required. Decoding-only variants hold the trained policy fixed and change decoding by adding grammar constraints, validator retry, or best-of-$N$ selection. These variants are necessary controls because they test whether ECPO's benefit comes from the learning objective rather than from output formatting alone.

\section{Experiments}
\label{sec:experiments}

\subsection{Experimental setup}
\label{sec:exp_setup}

\paragraph{Datasets and ranking settings.}
We evaluate ECPO on MAVEN-ERE and RAMS instantiations of evidence-certified candidate ranking. All systems consume the same frozen-extractor trajectories and the same exposed evidence scopes. We report three roster settings. \textbf{Closed roster} uses the fixed candidate pool $C_w$ supplied by the benchmark and isolates ranking/certification from candidate discovery. \textbf{Predicted roster} constructs candidates only from frozen extractor/entity-linker predictions and therefore measures the effect of candidate availability. \textbf{Hybrid roster} augments predicted candidates with a high-recall retrieval stage and then applies ranking/certification. In all settings, policy-visible candidate identifiers are randomized per window.

\paragraph{Baselines.}
We compare ECPO with ordinary LLM ranking policies using the same Qwen3-8B backbone~\citep{yang_qwen3_2025}: zero-shot prompting, SFT, GRPO, and DPO. We include structured rankers: RM-only scoring, LP-Recognizer, recognition-as-planning, landmark recognition, LambdaMART, and GraphRank. To isolate whether ECPO is needed beyond output formatting, we add strong controls: (i) grammar/JSON-constrained decoding, (ii) validator retry, (iii) best-of-$N$ selection by the reward model or verifier, (iv) RM-only ranking with deterministic evidence optimization, and (v) RM-only plus LLM post-hoc rationalization with and without validation. All primary comparison rows in the reported tables use completed measurements supplied for this revision.

\paragraph{Metrics.}
Ordinary ranking quality is measured with NDCG@10, MAP@10, and Hit@10. Audit quality is measured with FeasibleRate, ValidSpan@10, StepCoverage@10, Faithfulness@10, and ReExtractFaithfulness@10; detailed audit metrics are reported in Appendix Table~\ref{tab:full_audit_metrics}. The two primary ECPO metrics are \CertNDCG@10 and \EvidCons@10. \CertNDCG@10 gives ranking credit only when the candidate's evidence certificate is valid and decision-sufficient. \EvidCons@10 asks whether evidence-only reconstruction recovers the Top-$K$ candidate set.

\subsection{Main results: evidence-certified ranking}
\label{sec:exp_results}

Table~\ref{tab:main_certified_results} reports the primary comparison. Read it from left to right: ordinary ranking metrics show whether the method selects useful candidates; audit metrics show whether outputs are valid; certified metrics show whether ranking and evidence are coupled. The intended claim is not that ECPO always dominates RM-only on ordinary NDCG, but that ECPO improves the joint certified objective under a single deployed policy interface.

\begin{table*}[t]
\centering
\caption{Main evidence-certified ranking results at $K=10$. \CertNDCG{} and \EvidCons{} are the primary certified ranking metrics. Higher is better except Calls.}
\label{tab:main_certified_results}
\small
\begin{adjustbox}{max width=\textwidth}
\begin{tabular}{l l c c c c c c c}
\toprule
Dataset & Method & Calls & NDCG@10 & MAP@10 & Feasible & Faithful & \EvidCons@10 & \CertNDCG@10 \\
\midrule
MAVEN-ERE & Qwen3-8B zero-shot & 1 & 0.52 & 0.38 & 0.48 & 0.34 & 0.22 & 0.18 \\
MAVEN-ERE & Qwen3-8B-SFT & 1 & 0.61 & 0.46 & 0.77 & 0.57 & 0.36 & 0.31 \\
MAVEN-ERE & Qwen3-8B-GRPO & 1 & 0.65 & 0.50 & 0.82 & 0.64 & 0.44 & 0.39 \\
MAVEN-ERE & Qwen3-8B-DPO & 1 & 0.65 & 0.50 & 0.83 & 0.64 & 0.44 & 0.39 \\
MAVEN-ERE & LP-Recognizer & \textbf{0} & 0.57 & 0.42 & 0.92 & 0.54 & 0.35 & 0.31 \\
MAVEN-ERE & GRaP-style recognizer & \textbf{0} & 0.60 & 0.45 & 0.94 & 0.60 & 0.39 & 0.35 \\
MAVEN-ERE & Landmark recognizer & \textbf{0} & 0.59 & 0.44 & 0.93 & 0.58 & 0.38 & 0.34 \\
MAVEN-ERE & LambdaMART/LTR & \textbf{0} & 0.69 & 0.54 & 0.96 & 0.70 & 0.42 & 0.40 \\
MAVEN-ERE & GraphRank & \textbf{0} & 0.65 & 0.50 & 0.94 & 0.67 & 0.40 & 0.38 \\
MAVEN-ERE & RM-only + Align & \textbf{0} & \textbf{0.71} & \textbf{0.56} & \textbf{0.97} & 0.78 & 0.46 & 0.43 \\
MAVEN-ERE & Qwen3-8B-\ECPO & 1 & 0.70 & 0.55 & 0.94 & \textbf{0.84} & \textbf{0.65} & \textbf{0.58} \\
\midrule
RAMS & Qwen3-8B zero-shot & 1 & 0.44 & 0.31 & 0.45 & 0.27 & 0.16 & 0.13 \\
RAMS & Qwen3-8B-SFT & 1 & 0.53 & 0.37 & 0.71 & 0.43 & 0.25 & 0.21 \\
RAMS & Qwen3-8B-GRPO & 1 & 0.57 & 0.40 & 0.78 & 0.49 & 0.32 & 0.27 \\
RAMS & Qwen3-8B-DPO & 1 & 0.57 & 0.40 & 0.79 & 0.50 & 0.33 & 0.28 \\
RAMS & LP-Recognizer & \textbf{0} & 0.48 & 0.33 & 0.91 & 0.43 & 0.24 & 0.20 \\
RAMS & GRaP-style recognizer & \textbf{0} & 0.51 & 0.36 & 0.92 & 0.47 & 0.27 & 0.24 \\
RAMS & Landmark recognizer & \textbf{0} & 0.50 & 0.35 & 0.91 & 0.46 & 0.26 & 0.23 \\
RAMS & LambdaMART/LTR & \textbf{0} & 0.60 & 0.42 & 0.94 & 0.58 & 0.31 & 0.30 \\
RAMS & GraphRank & \textbf{0} & 0.56 & 0.39 & 0.92 & 0.54 & 0.29 & 0.26 \\
RAMS & RM-only + Align & \textbf{0} & \textbf{0.62} & \textbf{0.44} & \textbf{0.95} & 0.64 & 0.35 & 0.33 \\
RAMS & Qwen3-8B-\ECPO & 1 & 0.61 & 0.43 & 0.92 & \textbf{0.69} & \textbf{0.49} & \textbf{0.40} \\
\bottomrule
\end{tabular}
\end{adjustbox}
\end{table*}

Across both datasets, ECPO has the highest certified metrics (\EvidCons@10 and \CertNDCG@10), while RM-only + Align remains slightly stronger on ordinary NDCG@10 and MAP@10. This pattern matches the intended claim: ECPO improves decision--evidence coupling rather than replacing deterministic rankers as the cheapest ordinary ranker.

\subsection{Decoding-only and post-hoc evidence controls}
\label{sec:exp_decoding_controls}

Table~\ref{tab:decoding_controls} tests whether ECPO's gains can be replicated by keeping the model fixed and changing only decoding, repair, or evidence attachment. The comparison separates feasibility improvements from certified decision--evidence coupling: grammar constraints and validator retry primarily target valid output structure, while ECPO directly optimizes the certified metrics.

\begin{table*}[t]
\centering
\caption{Strong decoding-only and post-hoc evidence baselines. Retry and best-of variants report the actual number of LLM calls per window. Higher is better except Calls.}
\label{tab:decoding_controls}
\small
\begin{adjustbox}{max width=\textwidth}
\begin{tabular}{l l c c c c c c}
\toprule
Dataset & Method & Calls & NDCG@10 & Feasible & ReExtractFaith & \EvidCons@10 & \CertNDCG@10 \\
\midrule
MAVEN-ERE & GRPO, single decode & 1 & 0.65 & 0.82 & 0.56 & 0.44 & 0.39 \\
MAVEN-ERE & GRPO + grammar/JSON constrained decoding & 1 & 0.64 & 0.92 & 0.58 & 0.42 & 0.37 \\
MAVEN-ERE & GRPO + validator retry-$N$ & 5 & 0.66 & 0.94 & 0.66 & 0.50 & 0.44 \\
MAVEN-ERE & GRPO + best-of-$N$ by RM & 10 & \textbf{0.72} & 0.86 & 0.62 & 0.49 & 0.46 \\
MAVEN-ERE & GRPO + best-of-$N$ by verifier & 10 & 0.69 & 0.93 & 0.71 & 0.57 & 0.53 \\
MAVEN-ERE & RM-only + deterministic evidence optimizer & \textbf{0} & \textbf{0.72} & \textbf{0.98} & \textbf{0.78} & 0.53 & 0.49 \\
MAVEN-ERE & RM-only + LLM rationalizer & 1 & 0.71 & 0.76 & 0.58 & 0.38 & 0.35 \\
MAVEN-ERE & RM-only + LLM rationalizer + validator & 1 & 0.71 & 0.86 & 0.64 & 0.43 & 0.40 \\
MAVEN-ERE & \ECPO, single decode & 1 & 0.70 & 0.94 & 0.75 & \textbf{0.65} & \textbf{0.58} \\
\midrule
RAMS & GRPO, single decode & 1 & 0.57 & 0.78 & 0.41 & 0.32 & 0.27 \\
RAMS & GRPO + grammar/JSON constrained decoding & 1 & 0.56 & 0.90 & 0.42 & 0.32 & 0.26 \\
RAMS & GRPO + validator retry-$N$ & 5 & 0.58 & 0.93 & 0.48 & 0.38 & 0.31 \\
RAMS & GRPO + best-of-$N$ by RM & 10 & \textbf{0.63} & 0.82 & 0.46 & 0.38 & 0.34 \\
RAMS & GRPO + best-of-$N$ by verifier & 10 & 0.60 & 0.91 & 0.55 & 0.45 & 0.37 \\
RAMS & RM-only + deterministic evidence optimizer & \textbf{0} & \textbf{0.63} & \textbf{0.97} & 0.60 & 0.40 & 0.36 \\
RAMS & RM-only + LLM rationalizer & 1 & 0.62 & 0.75 & 0.43 & 0.29 & 0.25 \\
RAMS & RM-only + LLM rationalizer + validator & 1 & 0.62 & 0.84 & 0.47 & 0.33 & 0.29 \\
RAMS & \ECPO, single decode & 1 & 0.61 & 0.92 & \textbf{0.61} & \textbf{0.49} & \textbf{0.40} \\
\bottomrule
\end{tabular}
\end{adjustbox}
\end{table*}

The multi-call verifier-selected control narrows the certified-metric gap but still remains below ECPO on \CertNDCG@10 and \EvidCons@10 while using 10 LLM calls per window. RM-based best-of-$N$ selection reaches the best ordinary NDCG but does not provide comparable evidence-cycle recovery.

\subsection{Candidate-availability settings}
\label{sec:exp_rosters}
Closed-roster results isolate ranking and evidence certification; they should not be read as open-world candidate discovery. We therefore report predicted-roster and hybrid-roster experiments in Appendix~\ref{sec:app_roster_diagnostics}. These diagnostics measure retained windows, positive-candidate recall, average pool size, and certified ranking under the same evaluation interface. The intended use of this comparison is to separate two bottlenecks: whether the upstream system makes the relevant candidates available, and whether ECPO can certify the ranking once candidates are available.

\subsection{Ambiguous-window stress suite}
\label{sec:exp_hard_windows}
ECPO is not designed to make every easy window substantially better than a strong deterministic ranker. Its value is expected to concentrate in ambiguous windows where ordinary RM-only ranking is least sufficient: large candidate pools, near-neighbor distractors, sparse evidence, shared-document candidates, role conflicts, or skeleton ambiguity. Table~\ref{tab:hard_windows} therefore reports stress slices in addition to the overall mean. This stress suite is designed to measure whether the largest ECPO gains occur in \CertNDCG{} and \EvidCons{} on these slices, even if the average ordinary NDCG gain over RM-only is modest.

\begin{table*}[t]
\centering
\caption{Ambiguous-window stress-suite results. Stress categories are assigned before model comparison and are not mutually exclusive.}
\label{tab:hard_windows}
\small
\begin{adjustbox}{max width=\textwidth}
\begin{tabular}{l l l c c c c}
\toprule
Dataset & Stress slice & Method & \# windows & NDCG@10 & \EvidCons@10 & \CertNDCG@10 \\
\midrule
MAVEN-ERE & Large candidate pool & RM-only + Align & 37 & 0.62 & 0.38 & 0.35 \\
MAVEN-ERE & Large candidate pool & \ECPO & 37 & \textbf{0.64} & \textbf{0.58} & \textbf{0.51} \\
MAVEN-ERE & Near-neighbor negatives & RM-only + Align & 28 & 0.57 & 0.32 & 0.30 \\
MAVEN-ERE & Near-neighbor negatives & \ECPO & 28 & \textbf{0.61} & \textbf{0.55} & \textbf{0.48} \\
MAVEN-ERE & Sparse evidence & RM-only + Align & 22 & 0.55 & 0.29 & 0.27 \\
MAVEN-ERE & Sparse evidence & \ECPO & 22 & \textbf{0.58} & \textbf{0.50} & \textbf{0.43} \\
MAVEN-ERE & Role-conflict windows & RM-only + Align & 18 & 0.54 & 0.27 & 0.25 \\
MAVEN-ERE & Role-conflict windows & \ECPO & 18 & \textbf{0.59} & \textbf{0.53} & \textbf{0.46} \\
\midrule
RAMS & Large candidate pool & RM-only + Align & 45 & 0.54 & 0.29 & 0.27 \\
RAMS & Large candidate pool & \ECPO & 45 & \textbf{0.56} & \textbf{0.44} & \textbf{0.36} \\
RAMS & Near-neighbor negatives & RM-only + Align & 25 & 0.49 & 0.25 & 0.23 \\
RAMS & Near-neighbor negatives & \ECPO & 25 & \textbf{0.53} & \textbf{0.41} & \textbf{0.33} \\
RAMS & Sparse evidence & RM-only + Align & 52 & 0.47 & 0.23 & 0.21 \\
RAMS & Sparse evidence & \ECPO & 52 & \textbf{0.50} & \textbf{0.36} & \textbf{0.29} \\
RAMS & Role-conflict windows & RM-only + Align & 14 & 0.48 & 0.22 & 0.20 \\
RAMS & Role-conflict windows & \ECPO & 14 & \textbf{0.51} & \textbf{0.37} & \textbf{0.31} \\
\bottomrule
\end{tabular}
\end{adjustbox}
\end{table*}

\subsection{Ablations, human audit, and efficiency}
\label{sec:exp_ablation}
Full component ablations, identifier-leakage controls, external human-audit protocols, robustness tests, and inference-cost diagnostics are reported in Appendix~\ref{sec:app_additional_experiments}. The key ablation is removing $r_{\mathrm{cycle}}$: a successful result should leave ordinary NDCG relatively stable but reduce \EvidCons{} and \CertNDCG{}, showing that the evidence-cycle reward is responsible for certified coupling rather than JSON formatting alone. The key human-audit test asks annotators whether the cited spans both support the claimed skeleton steps and allow recovery of the selected candidates without seeing model scores.

\section{Conclusion}
\label{sec:conclusion}

We formulate \emph{evidence-certified candidate ranking}: a Top-$K$ candidate decision must be accompanied by \texttt{doc\_id:span} evidence that is valid, traceable, and sufficient to recover the decision. ECPO operationalizes this requirement by optimizing a joint ranking/evidence action with listwise reward, deterministic certificate validation, and an evidence-cycle reward from a label-free deterministic verifier that uses claim-stripped evidence bundles. The MAVEN-ERE and RAMS benchmark instantiations combine trajectory construction, skeleton alignment, hard negatives, preference pairs, closed/predicted/hybrid roster settings, and audit metrics centered on \CertNDCG{} and \EvidCons{}.

The resulting evaluation separates three questions: whether a method ranks useful candidates, whether its evidence certificates are valid, and whether the evidence alone recovers the selected candidates. RM-only scoring remains a strong low-cost option when deterministic evidence attachment is sufficient. ECPO is intended for deployments where the ranked decision and its evidence certificate must be produced and audited as one object.

\section{Ethical considerations and misuse mitigation}
\label{sec:ethics}

Evidence-certified ranking is a decision-support technology, not an automated adjudication system. ECPO outputs should never be used as the sole basis for enforcement, disciplinary action, surveillance targeting, or other irreversible decisions affecting individuals. The strict evidence interface is designed to support human review: every high-ranked candidate must be accompanied by traceable spans, and unsupported evidence should lower certified credit rather than be hidden behind fluent explanation text. Deployments should require calibrated thresholds, human-in-the-loop review, access controls, audit logs, bias and drift audits, and appeal or correction channels. The benchmark uses anonymized candidate entities from MAVEN-ERE and RAMS rather than a real deployment roster; any adaptation to sensitive domains requires additional legal, privacy, and fairness review. Expanded safeguards are in Appendix~\ref{sec:app_ethics}.

\IfFileExists{260426.bib}{%
\bibliographystyle{plainnat}
\bibliography{260426}%
}{%
\begin{thebibliography}{16}
\bibitem[Brockhoff et~al.(2021)]{brockhoff_process_2021} Brockhoff et~al. 2021. Process conformance and event-log analysis.
\bibitem[Burges(2010)]{burges_ranknet_2010} Burges. 2010. From RankNet to LambdaRank to LambdaMART: an overview.
\bibitem[Ebner et~al.(2020)]{ebner_multi-sentence_2020} Ebner et~al. 2020. Multi-sentence argument linking.
\bibitem[Hung et~al.(2016)]{hung_investigative_2016} Hung et~al. 2016. Investigative analytics and prioritization.
\bibitem[Park et~al.(2021)]{park_realizing_2021} Park et~al. 2021. Realizing process mining and conformance analysis.
\bibitem[Ramirez and Geffner(2009)]{ramirez_plan_2009} Ramirez and Geffner. 2009. Plan recognition as planning.
\bibitem[Ramirez and Geffner(2010)]{ramirez_probabilistic_2010} Ramirez and Geffner. 2010. Probabilistic plan recognition.
\bibitem[Shao et~al.(2024)]{shao_deepseekmath_2024} Shao et~al. 2024. DeepSeekMath and group-relative policy optimization.
\bibitem[Sutmuller et~al.(2020)]{sutmuller_getting_2020} Sutmuller et~al. 2020. Investigative prioritization.
\bibitem[Wang et~al.(2020)]{wang_maven_2020} Wang et~al. 2020. MAVEN: a massive general-domain event detection dataset.
\bibitem[Wang et~al.(2022)]{wang_maven-ere_2022} Wang et~al. 2022. MAVEN-ERE: event relation extraction.
\bibitem[Yang et~al.(2025)]{yang_qwen3_2025} Yang et~al. 2025. Qwen3 technical report.
\bibitem[Ziebart et~al.(2008)]{ziebart_maximum_2008} Ziebart et~al. 2008. Maximum entropy inverse reinforcement learning.
\end{thebibliography}%
}
\clearpage

\appendix

\section{Benchmark and Reproducibility Details}
\label{sec:app_benchmark_reproducibility}

\paragraph{Reader roadmap.}
For readability, the appendix is organized into five top-level sections.
Appendix~\ref{sec:app_benchmark_reproducibility} first motivates the benchmark, defines the task interface and candidate/evidence serialization, then gives artifact records, windowing/splitting protocols, stage/role mapping, reward-learning data, examples, and dataset statistics.
Appendix~\ref{sec:app_method_details} provides extractor, alignment, reward-learning, RL/validator, inference, KG/TGN, complexity, LLM-training, and baseline implementation details.
Appendix~\ref{sec:app_eval_metrics} defines ranking metrics and certificate-audit checks.
Appendix~\ref{sec:app_additional_experiments} lists the filled ECPO diagnostics: uncertainty, paired tests, decoding-only controls, roster diagnostics, identifier leakage controls, hard-window stress tests, robustness diagnostics, and qualitative examples.
Appendix~\ref{sec:app_additional_discussion} covers extended related work, ethics, and deployment boundaries.

\subsection{Benchmark motivation}
\label{sec:app_moved_dataset_motivation}

\paragraph{Why do we construct a new benchmark?}
Our target setting jointly requires (i) multi-candidate pools per ranking window, (ii) window-scoped text-derived trajectories with time/order and argument normalization, (iii) Top-$K$ ranking as the decision object, and (iv) evidence certificates whose cited \texttt{doc\_id:span} spans are valid, traceable, and sufficient for evidence-only recovery. Event extraction and relation benchmarks such as MAVEN-ERE~\citep{wang_maven-ere_2022}, MAVEN~\citep{wang_maven_2020}, and RAMS~\citep{ebner_multi-sentence_2020} provide strong event, relation, trigger, or argument supervision, but they do not directly evaluate candidate-level ranking with certified evidence. Plan-recognition and process-conformance settings evaluate compatibility between traces and plans, but they usually do not require a window-local Top-$K$ decision or a strict JSON certificate. Investigation-oriented prioritization often assumes structured logs or graphs~\citep{sutmuller_getting_2020,hung_investigative_2016}, whereas our benchmark starts from noisy text-derived trajectories with span provenance. We therefore construct a benchmark tailored to evidence-certified ranking and provide closed-roster, predicted-roster, and hybrid-roster settings to separate ranking from candidate availability.
\subsection{Problem framing details}
\label{sec:app_moved_intro_details}

This subsection expands the brief main-text discussion of the intent, skeleton, and window-graph abstractions.

\paragraph{Ranking windows and context snippets.}
A ranking window is the evidence scope for one Top-$K$ candidate-ranking decision.
It is not necessarily a calendar-time interval.
A context snippet is a local text segment used for upstream extraction or compact prompting.
In MAVEN-ERE, source-record snippets used for supervised extractor examples and compact evidence display are trigger-centered sentence neighborhoods.
In RAMS, compact source snippets are document prefixes.
Held-out extractor inference does not use gold trigger locations: the frozen extractor is applied to split documents or document-local text units, and predicted event spans are linked back to their surrounding local context for prompt compaction and auditing.
The number of snippets and documents contributing to each ranking window is an empirical property reported in Appendix~\ref{sec:app_moved_data_details} and Appendix Table~\ref{tab:ranking_window_stats}.

\paragraph{Intent, plan skeleton, and knowledge graph.}
In this paper, an \emph{intent} denotes a hypothesized multi-step activity type that we instantiate with an explicit plan skeleton.
In our benchmark instantiations, we use a four-stage inventory
\texttt{PREP}$\rightarrow$\texttt{PROBE}$\rightarrow$\texttt{EXECUTE}$\rightarrow$\texttt{OUTCOME}
as the canonical soft order.
Individual \texttt{intent\_id} skeleton records instantiate ordered steps, role requirements, and optional soft-precedence edges over this inventory.
We operationalize an intent with an explicit \emph{plan skeleton} $S$, which is an \emph{intent-level} abstraction shared by all windows labeled with the same \texttt{intent\_id}.
Importantly, $S$ is \emph{not} constructed by directly taking any single candidate's observed event sequence as the skeleton; instead, it specifies high-level step \emph{types} and key role constraints, and individual candidates are only partially observed realizations that may miss steps, contain spurious events, or have uncertain timestamp / document-order indices.
Ordering irregularities are captured by soft precedence-violation features after monotone alignment, not by arbitrary event transpositions in the DP recurrence.
In the main benchmark setting, each \texttt{intent\_id} is associated with a single predefined soft-precedence plan skeleton.
We allow missing steps and skipped events during alignment to handle partial observability.
To examine the effect of this assumption, Appendix~\ref{sec:app_hard_windows} reports diagnostic skeleton-perturbation analyses.
When reported, Appendix~\ref{sec:app_hard_windows} treats skeleton banks only as an auxiliary diagnostic; they are not part of the main method or main-result comparison.
Automatically discovering skeletons from raw text remains outside the scope of this paper.
Table~\ref{tab:skeleton_example} shows an illustrative skeleton record; it is schematic and does not expose any held-out labels or audit references.

\begin{table}[t]
\centering
\caption{Illustrative predefined skeleton record. Stage labels are abstract alignment labels; \texttt{OUTCOME} denotes the terminal or outcome stage in the released JSON schema.}
\label{tab:skeleton_example}
\small
\setlength{\tabcolsep}{4pt}
\begin{tabularx}{\linewidth}{l l Y}
\hline
Step & Stage & Required roles / soft-precedence description \\
\hline
$s_1$ & \texttt{PREP} & optional \texttt{Agent}; precedes later evidence-bearing steps when order is available \\
$s_2$ & \texttt{PROBE} & \texttt{Agent} and optional \texttt{Target}; soft-precedes execution evidence \\
$s_3$ & \texttt{EXECUTE} & \texttt{Agent} and \texttt{Target}; central action step for alignment \\
$s_4$ & \texttt{OUTCOME} & optional \texttt{Target}/\texttt{Context}; terminal or outcome evidence \\
\hline
\end{tabularx}
\end{table}

We also construct a within-window knowledge graph (KG) and encode it with a lightweight temporal message passing encoder; ablations disable this component by setting $\phi_{\text{graph}}(\tau)=\mathbf{0}$.
Additional details on candidate identities, shared documents, and evidence-span serialization are provided in Appendix~\ref{sec:app_moved_candidate_identity}.

This paper focuses on the ranking-and-certification interface itself: candidate prioritization is evaluated only through window-local evidence, randomized candidate identifiers, and auditable span certificates.

\subsection{Candidate identities, document sharing, and evidence serialization}
\label{sec:app_moved_candidate_identity}

In this paper, a \emph{candidate} is an anonymized entity cluster placed into a window-local candidate pool by a label-free resolver over role-bearing event arguments.
When no external roster is available, the main benchmark first fixes an eligible closed roster from an identifier-only source projection of role-bearing entities.
The resolver then links frozen-extractor predicted mentions only to this eligible roster; predicted-only entities absent from the closed roster are excluded from the main benchmark and are analyzed separately in the predicted-only and hybrid-roster diagnostics.
The candidate pool $C_w$ is local to a ranking window. For all policy prompts and generated outputs, candidate identifiers are randomized and window-local (for example, \texttt{cand\_001}); source-normalized entity IDs are retained only in hidden evaluator metadata.
Ranking, validation, and metric computation therefore compare candidates within the current $C_w$ and do not allow cross-window ID memorization.
Identifier controls, pooled over both datasets, are reported in Appendix~\ref{sec:app_id_controls}.
Although motivating applications often involve persons of interest, MAVEN-ERE and RAMS contain heterogeneous event participants; therefore, the released benchmark uses candidate entities after deterministic role normalization.
Candidate pools are constructed from entities that participate in agent-like or target-like roles, while purely contextual entities are not used as candidates unless they also appear in a candidate role.

A document may mention multiple candidates; therefore, the same document $d\in\mathcal{D}_w$ can contribute evidence spans to multiple candidates.
Likewise, a single predicted event record may appear in multiple candidate-centric trajectories if multiple candidates participate in its arguments.
Each trajectory remains auditable via its own \texttt{doc\_id:span} references.

We represent an evidence span conceptually as $\sigma=(\texttt{doc\_id},[l,r))$, where $[l,r)$ is a 0-indexed half-open character-offset interval.
In JSON, the bounds are serialized either as the two-integer field \texttt{span:[l,r]} or as the string \texttt{"l-r"}.
In both cases, the stored pair $(l,r)$ is interpreted by the validator as the same half-open interval $[l,r)$; the array notation is a serialization of endpoints rather than a closed mathematical interval.

\subsection{Detailed problem notation and event representation}
\label{sec:app_problem_details}

This subsection collects notation and serialization details moved out of the main text.
For a ranking window $w$, $C_w=\{c_{w,1},\ldots,c_{w,N_w}\}$ denotes the window-local candidate pool, and $\tau_w(c_{w,i})\equiv\tau_{w,i}$ denotes the candidate trajectory within $w$.
If no predicted event is available for a roster candidate, we retain an empty trajectory record rather than dropping the candidate.
Events are sorted by absolute time when available and otherwise by dataset-native or document-local order.

\begin{table}[t]
\centering
\caption{Core notation.}
\label{tab:notation}
\small
\setlength{\tabcolsep}{4pt}
\begin{tabularx}{\linewidth}{lY}
\hline
Symbol & Meaning \\
\hline
$w$ & ranking window / evaluation instance \\
$\texttt{intent\_id}(w)$ & intent identifier for window $w$ \\
$S_w$ & predefined skeleton associated with $\texttt{intent\_id}(w)$ \\
$C_w$ & candidate pool in window $w$ \\
$N_w$ & number of candidates, $|C_w|$ \\
$\tau_w(c_{w,i})\equiv \tau_{w,i}$ & trajectory of candidate $c_{w,i}$ within window $w$ \\
$K_w$ & evaluated cutoff, $\min(K,N_w)$ \\
$\phi(\tau;S)$ & skeleton/KG/argument feature vector \\
$R_\theta(\tau;S_w)$ & learned skeleton-conditioned trajectory-level reward \\
$y$ & generated JSON output containing Top-$K_w$ IDs and certificates \\
\hline
\end{tabularx}
\end{table}

Each event is represented as
\begin{equation}
e =
\big(
\mathrm{etype}^{\mathrm{raw}},
\mathcal{H}^{\mathrm{skel}},
\mathrm{etype}^{\mathrm{primary}},
\sigma^{\mathrm{trig}},
\mathcal{A},
t
\big),
\end{equation}
where $\mathrm{etype}^{\mathrm{raw}}$ is the extractor's fine-grained event type, $\mathcal{H}^{\mathrm{skel}}=\texttt{skeleton\_hits}(e)$ is the set of compatible skeleton stages, and $\mathrm{etype}^{\mathrm{primary}}$ is a single primary stage stored only for serialization and analysis.
Alignment and reward computation use $\texttt{skeleton\_hits}(e)$ rather than the primary stage label.
The argument set $\mathcal{A}$ contains tuples $(\texttt{role},\texttt{entity\_id},\sigma^{\mathrm{arg}})$.
The temporal field $t$ is an absolute timestamp when available, and otherwise a deterministic relative or document-local order key.
Serialized trajectory records may expose both an optional \texttt{time} field and a required \texttt{order\_index}; sorting prioritizes reliable absolute time and falls back to \texttt{order\_index}.
All triggers and arguments preserve auditable provenance via their spans $\sigma$.

For target cutoff $K$, the evaluated list length is $K_w=\min(K,N_w)$.
When a candidate pool contains fewer than $K$ candidates, the expected output contains all candidates in that pool.
Each returned candidate has one certificate object containing one serialized step object for every skeleton step in $S_w$; matched steps cite supporting \texttt{doc\_id:span} evidence and unmatched steps are recorded explicitly with empty evidence.

\subsection{Minimal auditable output interface}
\label{sec:app_moved_output_interface}

\paragraph{Auditable output interface (minimal).}
We require the policy to emit \emph{strict JSON} so that candidate IDs and \texttt{doc\_id:span} evidence can be deterministically validated.
Listing~\ref{lst:ecpo-json-min} shows a schematic minimal interface using the canonical \texttt{OUTCOME} stage label; released-record schemas are provided in Appendix~\ref{sec:app_data_jsonschema}, and the policy-output schema and deterministic validator are detailed in Appendix~\ref{sec:app_rl_output_jsonschema}--\ref{sec:app_rl_validator}.

\begin{lstlisting}[caption={Minimal auditable output interface (schematic JSON; $K_w=1$ for brevity).},label={lst:ecpo-json-min}]
{
  "window_id":"w_0001",
  "topk":["P_001"],
  "certificates":[{
    "steps":[
      {"step_id":"s1","etype":"PREP","matched":true,"event_id":"e1",
       "evidence":[{"doc_id":"doc123","span":[120,156],"kind":"trigger"}]},
      {"step_id":"s2","etype":"PROBE","matched":false,"event_id":null,"evidence":[]},
      {"step_id":"s3","etype":"EXECUTE","matched":false,"event_id":null,"evidence":[]},
      {"step_id":"s4","etype":"OUTCOME","matched":false,"event_id":null,"evidence":[]}
    ]}]
}
\end{lstlisting}

\subsection{JSON Schema excerpts for released records}
\label{sec:app_data_jsonschema}

This section provides compact JSON-Schema excerpts aligned with the released schema files in \texttt{schemas/}. 
We use a shared \texttt{span} type that supports two serializations: a string \texttt{"l-r"} or an integer array \texttt{[l,r]}.

\begin{lstlisting}[caption={Shared span definition used across schema files.},label={lst:schema-def-span}]
{
  "span": {
    "anyOf": [
      {"type": "string", "pattern": "^\\d+-\\d+$"},
      {
        "type": "array",
        "items": {"type": "integer"},
        "minItems": 2,
        "maxItems": 2
      }
    ]
  }
}
\end{lstlisting}

\begin{lstlisting}[caption={Schema excerpt for \texttt{rl\_prompts.schema.json} (RL Prompt/Demo Record; \texttt{response} optional).},label={lst:schema-rlprompts}]
{
  "$schema": "https://json-schema.org/draft/2020-12/schema",
  "title": "ECPO Prompt/Demo Record",
  "type": "object",
  "additionalProperties": false,
  "required": ["prompt"],
  "properties": {
    "prompt": {"type": "string"},
    "response": {"type": "string"},
    "trajectory_id": {"type": "string"},
    "candidate_id": {"type": "string"},
    "_meta": {
      "type": "object",
      "additionalProperties": false,
      "properties": {
        "trajectory_id": {"type": "string"},
        "candidate_id": {"type": "string"}
      }
    }
  }
}
\end{lstlisting}

We use the same schema for \texttt{rl\_prompts*.jsonl} (prompt-only; \texttt{response} omitted) and \texttt{rl\_demos*.jsonl} (prompt+response).

\begin{lstlisting}[caption={Schema excerpt for \texttt{traj.schema.json} (ECPO Trajectory Record).},label={lst:schema-traj}]
{
  "$schema": "https://json-schema.org/draft/2020-12/schema",
  "title": "ECPO Trajectory Record",
  "type": "object",
  "additionalProperties": false,
  "required": ["window_id", "candidate_id", "trajectory_id", "events"],
  "properties": {
    "window_id": {"type": "string"},
    "candidate_id": {"type": "string"},
    "trajectory_id": {"type": "string"},
    "events": {
      "type": "array",
      "items": {
        "type": "object",
        "additionalProperties": false,
        "required": ["event_id", "etype_raw", "skeleton_hits",
                     "etype_primary", "order_index", "trigger", "arguments"],
        "properties": {
  "event_id": {"type": "string"},
  "etype_raw": {"type": "string"},
  "skeleton_hits": {
    "type": "array",
    "items": {
      "type": "string",
      "enum": ["PREP", "PROBE", "EXECUTE", "OUTCOME"]
    },
    "minItems": 1,
    "uniqueItems": true
  },
  "etype_primary": {
    "type": "string",
    "enum": ["PREP", "PROBE", "EXECUTE", "OUTCOME"]
  },
  "time": {
    "description": "Optional absolute or normalized timestamp when available.",
    "anyOf": [
      {"type": "string"},
      {"type": "number"},
      {"type": "null"}
    ]
  },
  "order_index": {
    "type": "number",
    "description": "Required deterministic order key used for sorting within a window when absolute time is unavailable."
  },
          "trigger": {
            "type":"object",
            "additionalProperties": false,
            "required": ["doc_id", "span"],
            "properties": {
              "doc_id": {"type":"string"},
              "span": {"$ref":"#/$defs/span"}
            }
          },
          "arguments": {
            "type":"array",
            "items":{
              "type":"object",
              "additionalProperties": false,
              "required": ["role", "entity_id", "doc_id", "span"],
              "properties":{
                "role":{
                  "type":"string",
                  "enum":["Agent","Target","Context"]
                },
                "entity_id":{"type":"string"},
                "doc_id":{"type":"string"},
                "span":{"$ref":"#/$defs/span"}
              }
            }
          }
        }
      }
    },
    "meta": {
      "type": "object",
      "additionalProperties": false,
      "properties": {
        "graph_nodes": {"type": "array", "items": {"type": "string"}},
        "graph_edges": {"type": "array", "items": {"type": "array"}}
      }
    }
  },
  "$defs": {
    "span": {
      "anyOf": [
        {"type": "string", "pattern": "^\\d+-\\d+$"},
        {
          "type": "array",
          "items": {"type": "integer"},
          "minItems": 2,
          "maxItems": 2
        }
      ]
    }
  }
}
\end{lstlisting}
The model-input trajectory schema intentionally omits \texttt{label} and reference evidence fields.
Labels and audit references are stored in separate evaluation-only schemas.

\begin{lstlisting}[caption={Schema excerpt for \texttt{pairs.schema.json} (Preference Pair Record).},label={lst:schema-pairs}]
{
  "$schema": "https://json-schema.org/draft/2020-12/schema",
  "title": "ECPO Preference Pair Record",
  "type": "object",
  "additionalProperties": false,
  "required": ["window_id", "better", "worse"],
  "properties": {
    "window_id": {"type": "string"},
    "intent_id": {"type": "string"},
    "better": {
      "type": "string",
      "description": "trajectory_id of the preferred candidate trajectory"
    },
    "worse": {
      "type": "string",
      "description": "trajectory_id of the dispreferred candidate trajectory"
    },
    "reason": {"type": "string"}
  }
}
\end{lstlisting}

\begin{lstlisting}[caption={Schema excerpt for \texttt{sft.schema.json} (SFT Record).},label={lst:schema-sft}]
{
  "$schema": "https://json-schema.org/draft/2020-12/schema",
  "title": "ECPO SFT Record",
  "type": "object",
  "additionalProperties": false,
  "properties": {
    "instruction": {"type": "string"},
    "input": {"type": "string"},
    "output": {"type": "string"},
    "system": {"type": "string"},
    "history": {"type": "array"}
  }
}
\end{lstlisting}

\paragraph{Notes.}
The full schema files in the artifact remain the authoritative reference (including any stricter constraints and additional optional fields); the excerpts above are included to make the released interface explicit in the paper. Long descriptions and string fields in Listings~\ref{lst:schema-rlprompts}--\ref{lst:schema-sft} are line-wrapped for display only; the full schema files are the valid JSON Schema documents.

\subsection{Artifact package and reproducibility roadmap}
\label{sec:app_data}

This section specifies the artifact data format, span conventions, windowing/splitting protocol, and the minimal reproducibility package for the evidence-certified candidate ranking benchmark.

\paragraph{Package overview.}
The anonymized review-time artifact package is organized as follows, with any redistribution-restricted raw text replaced by derived records, checksums, and reconstruction scripts where necessary.
To support reproducibility during review and after release, we specify the exact record formats, schemas, split construction, validation rules, and evaluation metrics below.
(i) window-scoped text-derived trajectories with normalized timestamps or dataset-native order,
(ii) auditable provenance via \texttt{doc\_id:span},
(iii) window-level multi-candidate ranking instances, and
(iv) supervision artifacts for reward learning (expert/pseudo demonstrations, hard negatives, and preference pairs).

\begin{table}[ht]
\centering
\caption{Artifact components and which stage consumes them.}
\label{tab:artifact_components}
\small
\setlength{\tabcolsep}{4pt}
\begin{adjustbox}{max width=\linewidth}
\begin{tabular}{p{0.34\linewidth}p{0.58\linewidth}}
\hline
Component & Used for \\
\hline
\texttt{doc\_meta.jsonl} (raw documents included only when redistribution is permitted) & Span bounds, auditability checks, deterministic validation \\
\texttt{traj\_pred.jsonl} (frozen-extractor trajectories with predicted events and \texttt{doc\_id:span} provenance; model input for all ranking methods) 
& Model inputs, skeleton alignment, reward feature computation, ranking, and certificate generation \\
\texttt{window\_input.jsonl} (ranking windows + candidate pools + \texttt{intent\_id}/\texttt{skeleton\_id}; no positive labels or audit references) 
& Policy / RM input construction; contains no positive labels or audit-only fields \\
\texttt{window\_label.jsonl} (positive ids + split tags) 
& Training supervision on train split; evaluation labels on dev/test; never included in policy prompts \\
\texttt{audit\_ref.jsonl} (reference evidence chains) 
& Audit references, automatic audit checks, author-side manual spot checks, and aligner-independent verification; never included in model inputs \\
\texttt{skeleton.jsonl} (skeleton records keyed by \texttt{skeleton\_id} and \texttt{intent\_id}) 
& Skeleton alignment, reward features, certificate step definitions, and deterministic joins from \texttt{window\_input.jsonl} \\
\texttt{pairs.jsonl} (train-split trajectory preference pairs) 
& Bradley--Terry reward learning and construction of output-level DPO pairs \\
Extractor SFT (\texttt{*\_sft\_*.jsonl}) & Upstream trigger/argument extraction training (not the ranking policy) \\
Ranking policy demos (\texttt{rl\_demos*.jsonl}: prompt+response) & SFT for ranking policy and/or warm-start demonstrations \\
RL prompts (\texttt{rl\_prompts*.jsonl}: prompt only) & Augmented policy-optimization prompts derived from ranking windows \\
\hline
\end{tabular}
\end{adjustbox}
\end{table}

\paragraph{Predicted inputs versus reference labels.}
For each split, we separate model-input records from reference-label records.
The ranking policy, RM-only ranker, LTR baselines, and reward-feature extractor consume only
\texttt{traj\_pred.jsonl} and \texttt{window\_input.jsonl}.
Positive candidate ids, reference evidence chains, and audit-only fields are stored separately in
\texttt{window\_label.jsonl} and \texttt{audit\_ref.jsonl}.
During prompt construction, the fields \texttt{positive\_candidate\_ids}, \texttt{label}, and reference step-aligned evidence are removed.
Held-out reference annotations are used only for evaluation, automatic audit checks, author-side manual spot checks, and aligner-independent verification.

We provide, as part of the same anonymized artifact repository, preprocessing/evaluation scripts for:
(i) windowing and split generation,
(ii) trajectory construction and span validation,
(iii) feasibility checking utilities and JSON schemas, and
(iv) evaluation code for ranking and certificate-audit metrics.
Release metadata (license, checksum, and artifact identifier) is included in the artifact package.

\paragraph{Existing-asset licenses and terms of use.}
Table~\ref{tab:asset_licenses} summarizes the upstream datasets, model checkpoint, and major software libraries used by the benchmark and experiments.
For software dependencies, the artifact release includes an executable environment file and lockfile with exact patch versions, package hashes when available, and copied license notices.
We do not redistribute upstream raw text or model weights unless the corresponding license and hosting terms permit redistribution; otherwise, the artifact records only identifiers, derived structures, and scripts that operate on user-obtained authorized copies.

\begin{table}[ht]
\centering
\caption{Existing assets, versions, licenses, and use constraints. Exact software patch versions and copied license notices are in the artifact lockfile/README to keep this table readable.}
\label{tab:asset_licenses}
\small
\setlength{\tabcolsep}{3pt}
\begin{adjustbox}{max width=\linewidth}
\begin{tabular}{p{0.19\linewidth}p{0.20\linewidth}p{0.35\linewidth}p{0.21\linewidth}}
\hline
Asset & Version / source & License or terms stated by upstream & How we respect the terms \\
\hline
\texttt{MAVEN-ERE} dataset and reference code & Dataset version 1.0 from \url{https://github.com/THU-KEG/MAVEN-ERE}; source: upstream MAVEN-ERE paper/repository & Repository license: GPL-3.0. The MAVEN-ERE paper also states that the underlying MAVEN data are CC BY-SA 4.0 and Wikipedia text is CC BY-SA 3.0 / research use. & We cite the source, keep attribution, use authorized copies for research, and avoid relicensing or redistributing upstream raw text beyond permitted terms. Derived records preserve source identifiers and provenance. \\
\texttt{RAMS} dataset and scorer & RAMS 1.0 / 1.0c current release from \url{https://nlp.jhu.edu/rams/}; citation: \citet{ebner_multi-sentence_2020} & RAMS annotations are released under CC BY-SA 4.0; underlying article text remains subject to original source rights and the RAMS notice-and-takedown policy. The accompanying \texttt{arglinking} code is Apache-2.0. & We cite RAMS, use authorized copies for research, retain attribution/share-alike obligations for released annotation-derived artifacts, and do not claim ownership of source news text. \\
\texttt{Qwen3-8B} backbone & \texttt{Qwen/Qwen3-8B} checkpoint from \url{https://huggingface.co/Qwen/Qwen3-8B}; citation: \citet{yang_qwen3_2025} & Apache-2.0 license for the model repository/checkpoint. & We cite the model, preserve the model card/license notice, use the checkpoint for baselines and policy training, and do not redistribute weights unless allowed by the hosting/license terms. \\
LLM training libraries & \texttt{transformers} $\geq$4.51.0 for Qwen3 support; \texttt{trl}, \texttt{peft}, \texttt{bitsandbytes} compatible stable releases pinned in the lockfile & \texttt{transformers}, \texttt{trl}, and \texttt{peft}: Apache-2.0; \texttt{bitsandbytes}: MIT. & We call these libraries through public APIs, preserve license/copyright notices, and release exact dependency versions in the environment lockfile. \\
Core ML / ranking / graph libraries & \texttt{PyTorch} 2.x, \texttt{LightGBM} 4.x, \texttt{scikit-learn} 1.x, \texttt{NumPy} 2.x, \texttt{NetworkX} 3.x; exact patch versions pinned in the lockfile & \texttt{PyTorch}, \texttt{scikit-learn}, \texttt{NumPy}, and \texttt{NetworkX}: BSD-style licenses; \texttt{LightGBM}: MIT. & We use package APIs for training, baselines, metrics, and graph construction; released code will retain required notices and dependency metadata. \\
\hline
\end{tabular}
\end{adjustbox}
\end{table}

\paragraph{Extended output example.}
Listing~\ref{lst:ecpo-json-min} in Appendix~\ref{sec:app_moved_output_interface} shows a schematic minimal interface for auditable ranking outputs.
Below we provide an \emph{extended} example (including step-level fields and evidence metadata) used by our feasibility validator:
\begin{lstlisting}[caption={Extended example of ECPO output schema (JSON).},label={lst:ecpo-json-example}]
{
  "window_id": "w_0001",
  "topk": ["P_001"],
  "certificates": [
    {
      "steps": [
        {"step_id":"s1","etype":"PREP","matched":true,"event_id":"e1",
         "evidence":[{"doc_id":"doc123","span":[120,156],"kind":"trigger"},
                     {"doc_id":"doc123","span":[90,110],"kind":"arg","role":"Agent"}]},
        {"step_id":"s2","etype":"PROBE","matched":false,"event_id":null,
         "evidence":[]},
        {"step_id":"s3","etype":"EXECUTE","matched":false,"event_id":null,
         "evidence":[]},
        {"step_id":"s4","etype":"OUTCOME","matched":false,"event_id":null,
         "evidence":[]}
      ]
    }
  ]
}
\end{lstlisting}

For brevity, the above example uses $K_w{=}1$; in our experiments we use $K{=}10$ and set $K_w=\min(K,N_w)$ for each ranking window.

\paragraph{Span convention (\texttt{doc\_id:span}).}
We represent an evidence span as $\sigma=(\texttt{doc\_id},[l,r))$, where $[l,r)$ is a \emph{0-indexed half-open} character-offset interval over the raw document string (Python slicing convention).
This choice avoids ambiguity at boundaries and makes span validation deterministic.
All span-based checks (ValidSpan, Faithfulness) operate under this convention.
If the raw corpus provides byte offsets (UTF-8) rather than character offsets, we provide a conversion utility and store the canonical character offsets in \texttt{doc\_meta.jsonl}.

\paragraph{Temporal / order metadata.}
Each predicted event record stores temporal / order metadata used for serialization and sorting.
We expose two fields:
(i) an optional \texttt{time} field, which stores an absolute ISO-8601 date string when available or a normalized scalar only when the corpus provides such a value, and
(ii) a required \texttt{order\_index}, which is a deterministic within-window order key used when reliable absolute time is unavailable.
When both are available, sorting prioritizes absolute time and uses \texttt{order\_index} as a stable tie-breaker or fallback.

\subsection{Dataset construction and preprocessing details}
\label{sec:app_moved_data_details}

\paragraph{Relation to the condensed main-text setup.}
For space, the main text compresses task instantiation, datasets, extractor pretraining, reward-learning corpus construction, models/baselines, and metrics into a single experimental-setup subsection. This appendix, together with Appendix~\ref{sec:app_irl_corpus} and Appendix~\ref{sec:app_extractor}, provides the corresponding full details.

\paragraph{Upstream event supervision (extractor only).}
Event extraction is an upstream preprocessing step.
We pretrain the extractor with a staged curriculum:
\texttt{MAVEN} $\rightarrow$ \texttt{MAVEN-Arg} $\rightarrow$ \texttt{MAVEN-ERE} $\rightarrow$ \texttt{RAMS}
to improve argument-role consistency and reduce trajectory noise; details are in Appendix~\ref{sec:app_extractor}.
This curriculum affects only the extractor; the \emph{ranking benchmark instantiations} in this paper are built on \textbf{MAVEN-ERE} and \textbf{RAMS}.

\paragraph{Two evaluation instantiations.}
We instantiate the ranking benchmark in two variants derived from the source corpora \textbf{MAVEN-ERE} and \textbf{RAMS}.
Source-normalized records from these corpora define supervision, positive labels, and audit references, while model-input trajectories are constructed separately from frozen-extractor predictions.
For each ranking window $w$, we aggregate predicted event records by candidate identity, normalize timestamps when available, and otherwise preserve dataset-native or document-local order.
This yields a possibly empty predicted trajectory $\tau_{w,i}$ with auditable provenance \texttt{doc\_id:span} for triggers and arguments.

\paragraph{Windows and trajectory instances.}
Each evaluation instance is a \emph{ranking window} $w$ that defines the evidence scope $\mathcal{D}_w$, the window-local candidate pool $C_w$, source-normalized positive labels, audit references, and model-input trajectories for one Top-$K$ ranking prompt.
A trajectory is \emph{window-scoped}: it is defined for a specific $(w,c)$ pair, i.e., $\tau_{w,c}$ aggregates candidate-participating events that are observable within $\mathcal{D}_w$ and orders them by time or by dataset-native / document-local order when absolute time is unavailable.

We distinguish two event-record views during construction.
Source-normalized records are deterministic conversions of the dataset annotations.
They are used to derive positive candidate labels and reference evidence for audit metrics.
When no external candidate roster is available, we construct the benchmark candidate roster from an identifier-only projection of source-normalized role-bearing entities before relevance labels are attached.
This roster-construction step keeps only anonymized entity identifiers and document/window membership needed to define the candidate universe.
It discards event-stage labels, skeleton-chain validity, positive labels, reward scores, and reference evidence chains.
Thus, source-normalized records may define which anonymized entities are eligible candidates, but they do not provide model-input trajectories or ranking evidence.
Predicted records are produced by the frozen event extractor and are the only event records consumed by the reward model, ranking policy, RM-only ranker, and structured baselines.
The candidate pool $C_w$ is fixed before positive labels are attached, and positive candidate IDs are stored only in \texttt{window\_label.jsonl}; they are not included in model-input records.

The same underlying source entity may appear in multiple ranking windows when a cross-window identifier is available.
However, ranking decisions are made over window-local candidate IDs.
Thus, the number of window-scoped trajectory instances can exceed the number of unique source entities, and overlapping windows may reuse portions of the same document evidence.
Appendix Algorithm~\ref{alg:window_construction} summarizes the construction procedure.

\paragraph{Instantiation-specific context snippets.}
\textbf{MAVEN-ERE.} We construct one fixed \emph{trigger-centered context snippet} per event record used for extraction or compact prompting (no sliding stride): we locate the sentence containing the trigger and take at most three sentences in total (the trigger sentence plus up to one sentence before and one after; \texttt{context\_radius}=1).
\textbf{RAMS.} We do not perform sentence-based context construction for the upstream event records.
For extractor/SFT text fields, we flatten sentences into a document string and keep its first 120 characters as the local context snippet.
Ranking-policy prompts additionally include the \texttt{intent\_id}, skeleton, candidate IDs, and compact trajectory summaries or evidence pointers.

The full construction (including fallbacks when sentence boundaries are missing) and leakage-preventing splits are described in Appendix~\ref{sec:app_data_windowing}--\ref{sec:app_data_splits}.

\paragraph{Dataset statistics.}
The construction produces four logically separate record families:
\texttt{window\_input.jsonl} for window scopes, selected skeleton identifiers, and candidate pools,
\texttt{traj\_pred.jsonl} for frozen-extractor model-input trajectories,
\texttt{window\_label.jsonl} for positive candidate labels,
and \texttt{audit\_ref.jsonl} for source-normalized step-level audit references.
This separation is important because all ranking methods consume the same window inputs and predicted trajectories, while labels and audit references are used only for training-split supervision, model selection, or held-out evaluation.

We construct window-scoped candidate-centric trajectories, hard/near-neighbor negatives, perturbation-based negatives, and within-window preference supervision.
Table~\ref{tab:dataset_overview} summarizes key statistics for both instantiations; extended distributions and metadata are provided in Appendix Table~\ref{tab:mavenere_stats} and Appendix Table~\ref{tab:rams_stats}.

\begin{table}[t]
\centering
\caption{Overview of the two benchmark instantiations after preprocessing.
\textit{Note:} Candidate IDs are anonymized entity-cluster identifiers before window expansion.
The Top-$K$ task is defined over window-local candidate pools $C_w$, not over the global candidate-ID inventory.
Rows that summarize train-side supervision or augmentation are marked explicitly below and should not be compared directly with held-out ranking-window counts.
Table~\ref{tab:ranking_window_stats} reports the ranking-window statistics that determine the actual candidate-ranking task.}
\label{tab:dataset_overview}
\small
\setlength{\tabcolsep}{4pt}
\begin{tabularx}{\linewidth}{Yrr}
\hline
Statistic & MAVEN-ERE & RAMS \\
\hline
Documents & 3{,}623 & 9{,}124 \\
Source-normalized event records & 91{,}222 & 9{,}124 \\
Unique predicted event records contributing to model-input trajectories & 96{,}184 & 9{,}386 \\
Unique source-level candidate/entity IDs before window expansion & 3{,}623 & 8{,}730 \\
Serialized trajectory-pool records across ranking / reward-learning files & 7{,}245 & 8{,}864 \\
Train-split preference pairs & 3{,}622 & 134 \\
Train-split RL prompt records (augmented) & 171{,}584 & 654 \\
\hline
\end{tabularx}
\end{table}

Source-normalized event records denote dataset-native event annotations after deterministic normalization and schema conversion.
They are used to construct training-split supervision and held-out reference labels / audit references.
When no external candidate roster is available, an identifier-only projection of these records is also used to construct anonymized role-bearing candidate IDs before relevance labels are attached.
Model-input trajectory records are separate: the reward model, ranking policy, RM-only ranker, and structured baselines consume only fixed-extractor predictions stored in \texttt{traj\_pred.jsonl}.
Therefore, statistics that describe model inputs are reported under ``predicted'' records, while statistics that describe annotation-derived references are explicitly marked as ``source-normalized''.
The serialized-trajectory row in Table~\ref{tab:dataset_overview} counts the released trajectory-record pool across ranking and reward-learning files.
For MAVEN-ERE, this pool contains comparison-pool candidate trajectories, expert/pseudo-expert demonstrations, and hard/perturbation negatives; for RAMS, it contains window-scoped candidate trajectories together with a small number of train-only hard/perturbation negatives (Appendix~\ref{sec:app_irl_corpus}).
These counts summarize released trajectory records rather than held-out ranking instances, which are counted separately by ranking windows in Table~\ref{tab:ranking_window_stats}.
The train-split preference-pair and RL-prompt rows are likewise reported for reproducibility of reward/policy training rather than to characterize the held-out ranking task.
For dev/test ranking inputs, trajectories are built from frozen-extractor predictions; held-out source annotations are used to define positive labels and audit references.

\paragraph{Predicted-evidence coverage of positive candidates.}
Because held-out positives are defined from source-normalized reference chains while model inputs are frozen-extractor predictions, some positive candidates may have incomplete predicted evidence.
We therefore compute a permissive candidate-wise evidence-availability diagnostic: for each held-out positive $(w,c)$, we check within the same predicted trajectory $\tau_{w,c}$ whether at least one, at least two, or all source-normalized reference steps have a type-compatible predicted event with valid \texttt{doc\_id:span} provenance.
This diagnostic is not used for training or scoring, and it should be interpreted as a partial availability ceiling rather than as a full alignment-success rate because it does not by itself require role satisfaction or global order consistency.
Table~\ref{tab:positive_predicted_coverage} reports this diagnostic under the same split, provenance, and validation protocol as the main benchmark.

\begin{table}[t]
\centering
\caption{Predicted-evidence availability for held-out positive candidates.
A step is counted as available if a type-compatible predicted event with valid \texttt{doc\_id:span} provenance is present in the candidate trajectory.
This is a permissive availability diagnostic, not a full role-aware alignment metric.
Because reference chains may contain a single step, especially in RAMS, the ``all reference steps covered'' row is not necessarily a monotone refinement of the ``at least two covered steps'' row.}
\label{tab:positive_predicted_coverage}
\small
\begin{adjustbox}{max width=\linewidth}
\begin{tabular}{lcc}
\hline
Coverage diagnostic & MAVEN-ERE & RAMS \\
\hline
At least one covered step & 0.96 & 0.89 \\
At least two covered steps & 0.86 & 0.34 \\
All reference steps covered & 0.66 & 0.41 \\
\hline
\end{tabular}
\end{adjustbox}
\end{table}

\begin{table}[t]
\centering
\caption{Ranking-instance statistics.
A ranking window is the evaluation unit for Top-$K$ candidate ranking.}
\label{tab:ranking_window_stats}
\small
\setlength{\tabcolsep}{4pt}
\begin{tabularx}{\linewidth}{Yrr}
\hline
Statistic & MAVEN-ERE & RAMS \\
\hline
Ranking windows & 718 & 896 \\
Train / dev / test windows & 576/71/71 & 716/90/90 \\
Intent IDs & 16 & 12 \\
Skeleton records & 16 & 12 \\
Avg. candidates per window & 10.1 & 9.7 \\
Median candidates per window & 9 & 9 \\
Avg. positives per window & 1.3 & 1.1 \\
Windows with multiple positives (\%) & 18.1 & 7.4 \\
Windows with $|C_w|<10$ (\%) & 46.5 & 49.7 \\
Avg. exposed evidence units per window & 12.6 & 2.4 \\
Windows with evidence from multiple documents (\%) & 61.2 & 11.8 \\
Avg. trajectory-expanded predicted event mentions per trajectory & 24.7 & 1.07 \\
\hline
\end{tabularx}
\end{table}

The row ``unique predicted event records'' counts each frozen-extractor event record once.
The trajectory-length statistics count trajectory-expanded event mentions after window and candidate expansion.
Therefore, the same predicted event record may contribute to multiple candidate trajectories when multiple candidates participate in the event or when overlapping windows reuse the same evidence.
An exposed evidence unit is a trigger-centered context snippet for MAVEN-ERE and a document-prefix / document-local context unit for RAMS; multi-document evidence is computed from distinct \texttt{doc\_id} values.

\paragraph{RL prompt records.}
RL training prompt records denote policy-optimization prompts generated from eligible train-split windows by varying candidate subsets, demonstration order, and hard-negative composition.
They are not the number of distinct evaluation windows and are not expected to equal the number of train windows: after eligibility filtering, the count can be smaller than the train-window count in sparse settings or much larger when augmentation is dense.
Distinct ranking-window counts are reported in Table~\ref{tab:ranking_window_stats}.

\subsection{Windowing protocol}
\label{sec:app_data_windowing}

A \emph{ranking window} $w$ specifies the evidence scope for one Top-$K$ candidate-ranking instance and has an associated candidate pool, separately stored positive labels, and window-scoped trajectories.
The model-input window record stores the evidence scope, candidate pool, \texttt{intent\_id}, and selected \texttt{skeleton\_id}, while positive labels are stored only in \texttt{window\_label.jsonl}.
Separately, a \emph{context snippet} specifies the local text segment used for upstream extraction or compact prompting.
Because the source corpora do not uniformly provide reliable calendar-time intervals, ranking windows are not defined by calendar slicing.
Instead, both window construction and snippet construction follow dataset-native document or document-local grouping rules.

\paragraph{MAVEN-ERE context snippets (trigger-centered; \texttt{context\_radius}=1).}
For source-normalized records and supervised extractor examples, we construct exactly one fixed trigger-centered context snippet (no sliding stride):
\begin{enumerate}
    \item Locate the sentence containing the trigger via \texttt{resolve\_sent\_id}: use the annotated sentence ID when available; otherwise fall back to trigger-text matching.
    \item Take sentences in the index range
     $[\texttt{sent\_id}-\texttt{context\_radius},\ \texttt{sent\_id}+\texttt{context\_radius}]$
     with \texttt{context\_radius}=1, yielding at most three sentences.
    \item If the range is empty, fall back to the trigger sentence; if still empty, fall back to the concatenation of all non-empty sentences in the document; if still empty, fall back to the trigger text.
\end{enumerate}
The final context is the space-joined concatenation of the selected sentences (a coarse token count is computed by whitespace splitting).
This procedure is implemented in \texttt{build\_context\_snippet(...)} and used by
\texttt{build\_sft\_samples(..., context\_radius=1)} in the released preprocessing scripts.
For held-out extractor inference, trigger-centered source snippets are not used to choose where the extractor looks.
The frozen extractor is applied to the allowed split documents or document-local text units without gold trigger positions.
After prediction, each predicted trigger span is associated with its surrounding sentence neighborhood only for compact prompt construction and span auditing.

\paragraph{RAMS context snippets (document prefix; first 120 characters).}
RAMS does not use sentence-based context snippets for upstream event-record construction.
We flatten the document sentences into a single document string \texttt{doc\_text} and use \texttt{doc\_text[:120]} (the first 120 characters) as the local text field for extractor/SFT sample construction.
No sliding stride is used.
This procedure is implemented in \texttt{build\_events} and consumed by \texttt{build\_sft\_samples} via \texttt{event.mapping.get("doc\_text", "")}.

Ranking-policy prompts are constructed separately from ranking-window records.
They include the \texttt{intent\_id}, the skeleton, candidate IDs, compact trajectory summaries, and indexed evidence pointers.
Thus, the 120-character RAMS prefix is a context-snippet construction choice, not the full evidence interface for ranking-policy decoding.

\paragraph{Window records.}
Each window record stores at least \texttt{window\_id}, \texttt{doc\_ids}, optional \texttt{snippet\_ids}, the candidate pool \texttt{candidate\_ids}, the associated \texttt{intent\_id}, and the selected \texttt{skeleton\_id}.
A candidate can appear in multiple ranking windows.
MAVEN-ERE context snippets can overlap when nearby triggers occur in adjacent sentence neighborhoods; this snippet overlap is distinct from the definition of ranking windows.
The window record, together with the corresponding entries in \texttt{traj\_pred.jsonl}, forms the model-input unit consumed by ranking methods; context snippets are lower-level text segments used to create event records and compact prompt evidence.

\begin{algorithm}[!htbp]
\caption{Deterministic construction of ranking windows and candidate trajectories.}
\label{alg:window_construction}
\begin{algorithmic}[1]
\REQUIRE Documents $\mathcal{D}$; source annotations $\mathcal{A}^{\mathrm{src}}$; trained frozen extractor $f_{\mathrm{ext}}$; deterministic intent mapper $m_{\mathrm{intent}}$; skeleton table $\mathcal{S}_{\mathrm{file}}$; label-free candidate resolver $\rho_{\mathrm{id}}$; window grouping rule $g$.
\ENSURE \texttt{window\_input.jsonl}, \texttt{traj\_pred.jsonl}, \texttt{window\_label.jsonl}, and \texttt{audit\_ref.jsonl}.

\STATE Use document-level train/dev/test splits fixed before extractor training, reward learning, policy optimization, model selection, and evaluation.
\STATE Construct corpus-specific local text fields or snippet metadata used for preprocessing and prompt serialization, following the context-snippet rules above.
\STATE For held-out extractor inference, do not use gold trigger positions; apply $f_{\mathrm{ext}}$ only to allowed split documents or document-local text units.
\STATE Run the frozen extractor $f_{\mathrm{ext}}$ on all allowed split documents to obtain $\mathcal{E}^{\mathrm{pred}}$ with \texttt{doc\_id:span} provenance.
\STATE Convert source annotations $\mathcal{A}^{\mathrm{src}}$ into source-normalized records $\mathcal{E}^{\mathrm{src}}$ for supervision and held-out references.
\STATE Derive $\mathcal{E}^{\mathrm{src-id}}$ by retaining only anonymized role-bearing entity identifiers, document/snippet membership, and any corpus-native grouping keys required by $g$; remove labels, validated chains, and reference spans.

\FOR{each grouping key $q$ induced by $g$}
    \STATE Create a candidate ranking window $w$ with evidence scope $\mathcal{D}_w$, consisting of the documents or snippets assigned to $q$.
    \STATE Assign \texttt{intent\_id}$(w)=m_{\mathrm{intent}}(q)$ and select the predefined skeleton $S_w$ from $\mathcal{S}_{\mathrm{file}}$.
    \STATE Build a window-local candidate pool with the label-free resolver
    \[
    C_w=\rho_{\mathrm{id}}\!\left(\mathcal{E}^{\mathrm{src-id}}_w,\mathcal{E}^{\mathrm{pred}}_w\right),
    \]
    where $\mathcal{E}^{\mathrm{src-id}}_w$ defines the eligible closed roster using only anonymized role-bearing entity identifiers and window-membership metadata, and $\mathcal{E}^{\mathrm{pred}}_w$ is used only to attach predicted evidence and trajectories to those eligible candidates.
    Predicted-only entities that do not map to the closed roster are excluded from the main benchmark and are analyzed separately in the predicted-only roster diagnostic.
    \STATE The resolver does not use skeleton-chain validity, positive labels, reward scores, or audit reference spans.
    \FOR{each candidate $c \in C_w$}
        \STATE Collect candidate-participating predicted events from $\mathcal{E}^{\mathrm{pred}}_w$.
        \STATE Sort them by absolute timestamp when available; otherwise sort by dataset-native or document-local order.
        \STATE Serialize the resulting model-input trajectory as $\tau_{w,c}$ in \texttt{traj\_pred.jsonl}, preserving trigger and argument \texttt{doc\_id:span} provenance.
        \STATE If no predicted event is found for $c$, serialize an empty trajectory with the candidate ID and window ID retained; empty trajectories remain valid model inputs and are scored by missing-step features.
    \ENDFOR
    \STATE Serialize \texttt{window\_id}, \texttt{intent\_id}, the selected \texttt{skeleton\_id}, $\mathcal{D}_w$, and $C_w$ in \texttt{window\_input.jsonl}; the conceptual model input $x_w$ includes the resolved skeleton $S_w$, which is obtained in the released artifact by deterministically joining \texttt{window\_input.jsonl} with \texttt{skeleton.jsonl} on \texttt{skeleton\_id}.
    \STATE The serialized \texttt{intent\_id} remains the conditioning variable and can be checked against the linked skeleton record for consistency during prompt construction, reward-feature computation, and validation.
    \STATE Derive positive candidates $\mathcal{P}_w \subseteq C_w$ from source-normalized skeleton-aligned chains in $\mathcal{E}^{\mathrm{src}}_w$.
    \STATE Store $\mathcal{P}_w$ in \texttt{window\_label.jsonl}; do not store it in model-input files.
    \STATE Store source-normalized step-level reference spans in \texttt{audit\_ref.jsonl} for audit metrics only.
\ENDFOR

\STATE Discard any window whose documents cross split boundaries, or repair it before serializing records by removing cross-split documents and recomputing $C_w$, $\tau_{w,c}$, $\mathcal{P}_w$, and audit references.
\STATE Generate training-only perturbation-derived negatives and augmented RL prompts from train-split windows only.
\end{algorithmic}
\end{algorithm}

\subsection{Splits and leakage prevention}
\label{sec:app_data_splits}

To reduce leakage across splits, we use document-level splits, with time-aware constraints applied only when reliable timestamps are available:
\begin{itemize}
    \item No document (\texttt{doc\_id}) appears in multiple splits.
    \item A ranking window is assigned to a split only when all documents in its \texttt{doc\_ids} field belong to that split.
    \item If a constructed ranking window would contain documents from multiple splits, we either remove cross-split documents from the window before trajectory construction or discard the mixed-split window; no mixed-split window is used for training or evaluation.
    \item When reliable time ranges are available, we avoid placing overlapping time ranges across splits whenever possible; the document-level constraint always dominates.
\end{itemize}
We include explicit split lists for both documents and ranking windows:
\texttt{splits/doc\_\{train,dev,test\}.txt} contains document IDs, and
\texttt{splits/window\_\{train,dev,test\}.txt} contains window IDs.
All perturbation-derived negatives, candidate-subset augmentations, and RL prompt variants are generated only after the split assignment and only from train-split windows unless explicitly used in a train-only diagnostic.
Held-out dev/test ranking pools contain window-local candidates from \texttt{window\_input.jsonl}, not synthetic perturbation candidates.

\subsection{Candidate pools and labeling}
\label{sec:app_data_labeling}

For each ranking window $w$ in train, dev, or test, the candidate pool $C_w$ contains window-local candidate entity clusters that are observable in the window construction stage.
The pool is fixed before positive labels are attached.
For held-out dev/test windows, labels and audit references are used only for evaluation and are never serialized into model-input records.
When source-normalized records are used by the candidate resolver, they contribute only anonymized role-bearing entity identifiers for constructing the candidate universe; source-normalized skeleton-chain validity, reference step spans, and positive labels are not used in model-input construction.
The released model-input record \texttt{window\_input.jsonl} contains \texttt{window\_id}, \texttt{intent\_id}, the selected skeleton identifier, the exposed document or snippet identifiers, and the candidate IDs in $C_w$, but it does not contain positive labels, any field that marks a candidate as relevant, or reference evidence chains.

\paragraph{Candidate-roster diagnostic.}
The main benchmark uses an identifier-only source projection to define a closed candidate roster when no external candidate list is available.
To quantify the dependence on this choice, we also evaluate a predicted-only roster diagnostic in which candidate pools are constructed from frozen-extractor predictions only.
This diagnostic is not used as the main benchmark setting because upstream extraction may omit true held-out candidates entirely, making some windows unrankable rather than merely difficult.
It separates candidate-availability errors from ranking errors by reporting retained-window coverage, positive-candidate recall in the predicted-only roster, and Top-$K$ ranking scores on the retained subset.
The diagnostic is computed only for analysis and is not used for reward learning, policy training, model selection, or the main test comparison.

\begin{table}[t]
\centering
\caption{Predicted-only candidate-roster diagnostic.
Retained windows are windows whose candidate pool remains non-empty under frozen-extractor-only roster construction.
Positive-candidate recall measures whether held-out positive candidates remain available in the predicted-only roster.}
\label{tab:predicted_only_roster}
\small
\setlength{\tabcolsep}{4pt}
\begin{tabularx}{\linewidth}{Ycc}
\hline
Diagnostic & MAVEN-ERE & RAMS \\
\hline
Retained test windows & 70/71 (98.6\%) & 89/90 (98.9\%) \\
Positive-candidate recall in predicted-only roster & 0.87 & 0.76 \\
NDCG@10 on retained subset & 0.63 & 0.54 \\
MAP@10 on retained subset & 0.48 & 0.39 \\
Hit@10 on retained subset & 0.76 & 0.66 \\
\hline
\end{tabularx}
\end{table}

A candidate $c\in C_w$ is labeled positive if there exists a validated source-normalized chain aligned to $S_w$ and supported by \texttt{doc\_id:span} evidence within the window.
Formally, the positive set is
\[
\mathcal{P}_w =
\left\{
c\in C_w :
\exists\, z \in \mathcal{Z}^{\mathrm{src}}(w,c,S_w)
\ \text{such that}\ 
\mathrm{valid}(z)=1
\right\},
\]
where $\mathcal{Z}^{\mathrm{src}}(w,c,S_w)$ denotes source-normalized candidate chains that can be aligned to the skeleton $S_w$ using evidence inside $\mathcal{D}_w$.
Windows may contain multiple positives.
We store positive candidate IDs only in \texttt{window\_label.jsonl}.
Reference step-aligned evidence used by audit metrics is stored in \texttt{audit\_ref.jsonl}.
Model-input trajectories are constructed separately from frozen-extractor predictions.
For each $(w,c)$, \texttt{traj\_pred.jsonl} stores the predicted trajectory $\tau_{w,c}$, including predicted event types, compatible skeleton stages, normalized arguments, ordering information, and trigger / argument \texttt{doc\_id:span} provenance.
Thus, the ranker may receive noisy or incomplete trajectories even when the source-normalized reference contains a positive chain.

Same-window near-neighbor negatives are non-positive candidates in $C_w$ that share superficial context with positives, such as similar local event types, overlapping documents, entity-role patterns, or nearby temporal/order positions, but lack a validated skeleton-aligned chain.
Perturbation-derived negatives, including insertion, deletion, time-shift, and role-swap variants, are used for reward-learning comparison sets and train-split augmented policy prompts.
They are not mixed into dev/test ranking pools as ordinary entity candidates unless explicitly marked as synthetic train-only distractors.

\paragraph{Intent identifiers.}
Each \texttt{intent\_id} is assigned at the ranking-window level by a deterministic preprocessing rule before reward learning, policy training, model selection, and evaluation.
The identifier specifies the hypothesized plan type for the window and selects a skeleton record from \texttt{skeleton.jsonl}; all candidates in the same ranking window receive the same \texttt{intent\_id}.
It does not identify positive candidates, encode candidate-specific relevance, or expose reference evidence chains.
In the released instantiations, the window-to-\texttt{intent\_id} mapping is stored as preprocessing metadata and is applied identically across train, dev, and test splits.
Candidate relevance remains a separate held-out ranking label stored only in \texttt{window\_label.jsonl}.
In the main experiments, the number of skeleton records equals the number of distinct \texttt{intent\_id} values unless explicitly stated otherwise.

\subsection{Deterministic stage mapping from fine-grained event types to intent-stage labels}
\label{sec:app_stage_mapping}

This appendix specifies the deterministic and reproducible rule used to map the extractor's fine-grained event types $\texttt{etype}^{\mathrm{raw}}$
into the intent-stage labels $\{\texttt{PREP},\texttt{PROBE},\texttt{EXECUTE},\texttt{OUTCOME}\}$ used by skeleton alignment and reward computation.
The label \texttt{OUTCOME} denotes an outcome or terminal stage for both MAVEN-ERE and RAMS.
A domain-specific deployment may refine this stage name without changing the alignment or reward-learning algorithms.
The mapping is applied during preprocessing and the resulting stage labels are stored in the released trajectory records.

\paragraph{What is mapped and where it appears in the artifact.}
In the released \texttt{traj\_pred.jsonl}, each predicted event record stores a fine-grained type string under \texttt{etype\_raw}, a derived list of compatible skeleton stages under \texttt{skeleton\_hits}, and a single serialization label under \texttt{etype\_primary}.
The \texttt{skeleton\_hits} field is what the downstream DP aligner consumes for the type-compatibility check: a skeleton step of stage $\texttt{etype}_k$ can match an event if $\texttt{etype}_k \in \texttt{skeleton\_hits}(e)$.
This design supports partial observability by allowing some fine-grained event types to be compatible with multiple stages (e.g., a cognition/communication event can support both \texttt{PREP} and \texttt{PROBE}).

\paragraph{Mapping artifact (override file) and fallback behavior.}
The released preprocessing code supports an optional override mapping file,
\texttt{mapping/event2skeleton.json}, which maps a fine-grained type string to a possibly multi-label list of stage labels.
This file is included to make domain-specific remapping auditable, but it is not activated in the processed MAVEN-ERE and RAMS runs reported here.
For the reported experiments, all event types are mapped by the deterministic fallback mapper in Listing~\ref{lst:stage_mapping_pseudocode}.

Thus, the override-activation rate is $0\%$ and the deterministic fallback mapper handles all records in both instantiations (Table~\ref{tab:stage_mapping_coverage}).
This is intentional: the reported benchmark uses a single documented fallback mapper with fixed constants and one explicit corpus flag for the RAMS \texttt{PROBE}-primary subset, rather than a learned or development-tuned mapping.
The optional override file is reserved for domain-specific deployments and is not activated in the reported MAVEN-ERE or RAMS runs.

\begin{lstlisting}[caption={Optional stage override mapping file format (illustrative JSON; not activated in the reported runs).},label={lst:stage_mapping_json_ex}]
{
  "Know": ["PREP","PROBE"],
  "Warning": ["PREP","PROBE"],
  "Statement": ["PREP","PROBE"],
  "Process_start": ["PREP","PROBE"],

  "Attack": ["EXECUTE"],
  "Hostile_encounter": ["EXECUTE"],

  "Process_end": ["OUTCOME"],
  "Earnings_and_losses": ["OUTCOME"]
}
\end{lstlisting}

\paragraph{Primary stage vs.\ auxiliary \texttt{PROBE} tag (observed instantiation counts).}
Our mapper assigns each event a \emph{primary} stage in
$\{\texttt{PREP},\texttt{PROBE},\texttt{EXECUTE},\texttt{OUTCOME}\}$ by priority rules, and may additionally attach \texttt{PROBE} as an auxiliary compatibility tag for information-gathering, communication, or cognition events.
In MAVEN-ERE, \texttt{PROBE} appears only as an auxiliary tag in the source-normalized records; in RAMS, a small subset is \texttt{PROBE}-primary.
This matches the observed accounting in the preprocessing logs:

\begin{itemize}
    \item \textbf{MAVEN-ERE.} There are $91{,}222$ source-normalized event records. The mapper assigns a primary stage to all records
    ($82{,}688$ as \texttt{PREP}, $0$ as \texttt{PROBE}, $6{,}955$ as \texttt{EXECUTE}, and $1{,}579$ as \texttt{OUTCOME}; sums to $91{,}222$).
    In addition, $6{,}071$ events receive the auxiliary \texttt{PROBE} tag, yielding $97{,}293$ total (event,stage) instantiations.
    \item \textbf{RAMS.} There are $9{,}124$ source-normalized event records. The primary-stage distribution is
    $6{,}372$ \texttt{PREP}, $1{,}562$ \texttt{EXECUTE}, $1{,}135$ \texttt{OUTCOME}, and the remaining $55$ events are \texttt{PROBE}-primary.
    Additionally, $1{,}259$ events receive \texttt{PROBE} as an auxiliary tag, yielding $10{,}383$ total instantiations.
\end{itemize}

\paragraph{Deterministic fallback mapping procedure.}
Listing~\ref{lst:stage_mapping_pseudocode} gives the deterministic fallback procedure used in our preprocessing.
The rule uses only: (i) the fine-grained type string, (ii) the trigger surface form (when available from the extractor), and (iii) a dataset flag (to allow a small PROBE-primary subset in RAMS).
All sets below are fixed constants in the preprocessing script; no learning is involved.

\begin{lstlisting}[caption={Deterministic fallback stage mapping procedure (pseudocode).},label={lst:stage_mapping_pseudocode}]
# Inputs:
#   dataset: string in {"MAVEN-ERE","RAMS"}
#   etype_raw: string (fine-grained type, e.g., "Attack", "Know", "Process_end")
#   trig_text: string (trigger surface form, may be None if unavailable)
#
# Output:
#   skeleton_hits: list[string] in {"PREP","PROBE","EXECUTE","OUTCOME"} (deduplicated)

PROBE_TAG_TYPES = {
  "Know", "Warning", "Statement", "Process_start",
  "Reporting", "Suspicion", "Perception_active", "Coming_to_believe"
}
PROBE_TRIGGERS = {
  "reported", "report", "said", "warned", "warning", "according"
}

EXECUTE_TYPES_STRONG = {"Attack", "Hostile_encounter"}
# Additional execute cues are trigger-driven to keep EXECUTE high-precision.
EXECUTE_TRIGGERS = {"killed", "attack", "fought", "destroyed", "battle", "war", "hurricane", "storm"}

OUTCOME_TYPES = {"Process_end", "Earnings_and_losses", "Transaction", "Receiving", "Releasing"}
OUTCOME_TRIGGERS = {"ended", "paid", "sold", "bought", "earned", "withdrew", "deposited"}

# A small PROBE-primary subset is allowed in RAMS (empirically 55/9124 events).
PROBE_PRIMARY_TYPES_RAMS = {"Statement", "Reporting"}

def norm(s):
    # deterministic string normalization; no learned or dev-tuned mapping is used
    return s.strip()

def lower_or_empty(s):
    return "" if s is None else s.lower().strip()

def map_to_skeleton_hits(dataset, etype_raw, trig_text=None):
    et = norm(etype_raw)
    tr = lower_or_empty(trig_text)

    # 0) Optional override mapping (not triggered in our processed datasets; fallback=100% in Table)
    # if et in OVERRIDE_MAP:
    #     return sorted(unique(OVERRIDE_MAP[et]))

    # 1) Decide primary stage by priority: OUTCOME > EXECUTE > (PROBE-primary for RAMS) > PREP
    if (et in OUTCOME_TYPES) or (tr in OUTCOME_TRIGGERS):
        primary = "OUTCOME"
    elif (et in EXECUTE_TYPES_STRONG) or (tr in EXECUTE_TRIGGERS):
        primary = "EXECUTE"
    elif (dataset == "RAMS") and (et in PROBE_PRIMARY_TYPES_RAMS):
        primary = "PROBE"
    else:
        primary = "PREP"

    # 2) Auxiliary PROBE tag for information-gathering events.
    # We attach PROBE only when the primary stage is PREP, to keep the interface simple and stable.
    add_probe = ((et in PROBE_TAG_TYPES) or (tr in PROBE_TRIGGERS)) and (primary == "PREP")

    hits = [primary]
    if add_probe:
        hits.append("PROBE")

    # 3) Return a deterministic order for serialization
    order = {"PREP":0, "PROBE":1, "EXECUTE":2, "OUTCOME":3}
    hits = sorted(list(set(hits)), key=lambda z: order[z])
    return hits
\end{lstlisting}

\paragraph{Coverage statistics.}
Table~\ref{tab:stage_mapping_coverage} reports the number of source-normalized (event,stage) instantiations produced by the mapper and their stage percentages.
This table is a schema-mapping coverage table, not an extractor-performance table.
The same deterministic mapper is also applied to frozen-extractor predicted event records before trajectory construction; all predicted records used in \texttt{traj\_pred.jsonl} receive at least one \texttt{skeleton\_hits} label.
We report source-normalized coverage here because it is stable across extractor runs and determines the reference-label/audit view, while the model-input predicted-record counts are reported separately in Table~\ref{tab:dataset_overview}.
In our processed datasets, the override file is never triggered, hence all source-normalized and predicted records use the deterministic fallback mapper.

\begin{table*}[t]
\centering
\caption{Stage-mapping coverage after preprocessing, reported over source-normalized \emph{(event,stage)} instantiations rather than unique event records.
Override-activation rate is $0\%$ for both instantiations; the deterministic fallback mapper handles all records.}
\label{tab:stage_mapping_coverage}
\small
\setlength{\tabcolsep}{6pt}
\begin{adjustbox}{max width=\linewidth}
\begin{tabular}{lrr|rr}
\hline
& \multicolumn{2}{c|}{MAVEN-ERE} & \multicolumn{2}{c}{RAMS} \\
\cline{2-5}
Stage & \# inst. & \% inst. & \# inst. & \% inst. \\
\hline
\texttt{PREP}    & 82688 & 84.99 & 6372 & 61.37 \\
\texttt{PROBE}   & 6071  & 6.24  & 1314 & 12.66 \\
\texttt{EXECUTE} & 6955  & 7.15  & 1562 & 15.04 \\
\texttt{OUTCOME} & 1579  & 1.62  & 1135 & 10.93 \\
\hline
Total            & 97293 & 100.0 & 10383 & 100.0 \\
\hline
\end{tabular}
\end{adjustbox}
\end{table*}

\paragraph{Auditing note.}
Stage mapping is purely deterministic and applied before any learning stage.
Since \texttt{skeleton\_hits} are stored alongside \texttt{doc\_id:span} references in the released trajectories,
the entire downstream alignment and certificate evaluation remains auditable and reproducible under the strict evidence interface.

\subsection{Deterministic role normalization}
\label{sec:app_role_mapping}

We normalize dataset-native argument roles to the role inventory
$\{\texttt{Agent},\texttt{Target},\texttt{Context}\}$ used by skeleton alignment and argument-consistency features.
The mapping is deterministic and applied before reward learning.

Agent-like roles are dataset-native roles that identify the main acting or initiating entity of an event, such as attacker, speaker, mover, buyer, seller, or sender when these roles are present in the source schema.
Target-like roles identify the main affected or acted-upon entity, such as victim, attacked entity, destination, object, recipient, or affected party when applicable.
Temporal, locative, instrumental, descriptive, and residual roles are mapped to \texttt{Context}.
When a dataset-native role is ambiguous, the released \texttt{mapping/role2norm.json} file resolves it deterministically; uncertain residual roles are mapped to \texttt{Context} rather than to \texttt{Agent} or \texttt{Target}.

The released artifact includes \texttt{mapping/role2norm.json}; the validator and reward-feature extractor use the same mapping file.

\begin{table}[t]
\centering
\caption{Role normalization coverage after preprocessing.}
\label{tab:role_mapping_coverage}
\small
\begin{adjustbox}{max width=\linewidth}
\begin{tabular}{lrr|rr}
\hline
& \multicolumn{2}{c|}{MAVEN-ERE} & \multicolumn{2}{c}{RAMS} \\
\cline{2-5}
Normalized role & Count & \% & Count & \% \\
\hline
\texttt{Agent} & 168{,}431 & 53.21 & 12{,}960 & 51.43 \\
\texttt{Target} & 137{,}720 & 43.51 & 8{,}100 & 32.14 \\
\texttt{Context} & 10{,}398 & 3.28 & 4{,}140 & 16.43 \\
\hline
\end{tabular}
\end{adjustbox}
\end{table}

\subsection{Skeleton records and construction protocol}
\label{sec:app_skeleton_construction}

Each \texttt{intent\_id} is associated with one predefined soft-precedence skeleton in the main benchmark setting.
For serialization and deterministic joins, each skeleton record is stored in \texttt{skeleton.jsonl} with both a unique \texttt{skeleton\_id} and its associated \texttt{intent\_id}.
A skeleton record contains:
(i) \texttt{skeleton\_id},
(ii) \texttt{intent\_id},
(iii) an ordered list of skeleton steps,
(iv) each step's stage label in
$\{\texttt{PREP},\texttt{PROBE},\texttt{EXECUTE},\texttt{OUTCOME}\}$,
(v) required normalized roles from
$\{\texttt{Agent},\texttt{Target},\texttt{Context}\}$,
and (vi) optional soft-precedence edges.
In the main release, the mapping from \texttt{intent\_id} to \texttt{skeleton\_id} is one-to-one unless explicitly stated otherwise.

The predefined skeleton is not copied from any single candidate trajectory.
It is specified at the stage level to represent an abstract intent pattern, and candidate trajectories are treated as noisy partial realizations.
The construction protocol is deterministic:
(i) we first fix the stage inventory
$\{\texttt{PREP},\texttt{PROBE},\texttt{EXECUTE},\texttt{OUTCOME}\}$;
(ii) we define the canonical stage order and soft-precedence edges at the intent level;
(iii) we assign required roles to each step from the normalized role inventory
$\{\texttt{Agent},\texttt{Target},\texttt{Context}\}$;
and (iv) we store the resulting skeleton in \texttt{skeleton.jsonl}.
Skeleton records are fixed before reward learning, policy training, development tuning, and test evaluation.
No dev/test labels or dev/test model outputs are used to define or modify skeleton records.
The same skeleton file is used by trajectory alignment, reward-feature extraction, certificate serialization, and the feasibility / faithfulness validators.

\begin{lstlisting}[caption={Skeleton record format (schematic JSONL).},label={lst:skeleton_record}]
{
  "skeleton_id": "skel_001",
  "intent_id": "intent_001",
  "steps": [
    {"step_id": "s1", "etype": "PREP", "required_roles": ["Agent"]},
    {"step_id": "s2", "etype": "PROBE", "required_roles": ["Agent","Target"]},
    {"step_id": "s3", "etype": "EXECUTE", "required_roles": ["Agent","Target"]},
    {"step_id": "s4", "etype": "OUTCOME", "required_roles": ["Agent"]}
  ],
  "precedence": [["s1","s2"], ["s2","s3"], ["s3","s4"]]
}
\end{lstlisting}

In general-domain event corpora, \texttt{OUTCOME} should be read as an abstract terminal or outcome stage.

For the main experiments, all windows with the same \texttt{intent\_id} use the same skeleton record.
Appendix~\ref{sec:app_hard_windows} evaluates controlled perturbations of this design.

\subsection{Reward-learning corpus construction details}
\label{sec:app_irl_corpus}

ECPO learns a trajectory-level plan-consistency reward from indirect supervision.
We construct a reward-learning corpus containing \emph{expert positives}, \emph{pseudo-experts}, \emph{same-window hard (near-neighbor) negatives}, \emph{train-only perturbation-derived negatives}, and \emph{preference pairs}.
Here ``reward-learning corpus'' refers to the training-side trajectory pool that participates in reward-model comparison sets or within-window partition terms; it therefore includes unlabeled candidate-pool trajectories when they serve as same-window alternatives, not only explicitly labeled positives or negatives.

\paragraph{Corpus scale (MAVEN-ERE).}
In the processed MAVEN-ERE instantiation, the released trajectory pool contains 7,245 trajectory records in total:
3,774 comparison-pool candidate trajectories carried without explicit expert/pseudo-expert labels,
755 expert/pseudo-expert demonstrations, and
2,716 hard/perturbed negatives.
In addition, we provide 3,622 within-window preference pairs covering 7,244 unique trajectory IDs (Table~\ref{tab:mavenere_stats}).

\paragraph{Corpus scale (RAMS).}
In the processed RAMS instantiation, the reward-learning trajectory pool contains 8,864 trajectory records: 8,730 window-scoped candidate trajectories from ranking windows and 134 train-only hard/perturbation negatives, with an average trajectory length of $1.07$ events over the candidate trajectories.
The preference supervision contains 134 within-window pairs covering 268 unique trajectory IDs, and we additionally provide 654 RL prompt records for window-level policy optimization (Table~\ref{tab:rams_stats}).

\paragraph{Expert positives and evaluation labels.}
We separate supervision labels from ranking inputs.
Expert/pseudo-expert labels and evaluation positives are derived from dataset-native evidence and validated skeleton-aligned chains.
These annotations are used to label or validate candidate trajectories, but the trajectory records consumed by the reward model and ranking policies are produced by the fixed upstream extractor.
Thus, training annotations provide supervision for reward learning and policy training on the train split, while held-out annotations define positive labels and audit references on dev/test splits.
All ranking methods operate on the same fixed-extractor trajectory inputs in their respective splits.

For training-split expert/pseudo-expert records, each validated step stores
(i) an intent-stage label (\texttt{PREP/PROBE/EXECUTE/OUTCOME});
(ii) normalized roles (\texttt{Agent/Target/Context});
(iii) linked entity IDs;
(iv) timestamps or order indices; and
(v) auditable evidence spans (\texttt{doc\_id} with character offsets).
For dev/test ranking inputs, the corresponding fields are generated by the fixed extractor and are treated as noisy observations.
\emph{Operational note:} in a domain deployment, the same interface could additionally consume analyst-curated POI records from an investigation KB, but this is not required for (and not used in) our released benchmark.

\paragraph{Pseudo-experts from event corpora.}
To increase demonstration diversity when expert candidate-level demonstrations are sparse, we mine pseudo-expert trajectories from upstream event corpora by selecting chains with high skeleton-alignment scores, yielding skeleton-consistent yet non-identical demonstrations.

\paragraph{Hard, near-neighbor, and perturbation-derived negatives.}
Same-window hard or near-neighbor negatives are non-positive candidates in the same ranking window that share superficial similarity with positives, such as overlapping documents, similar local event types, nearby temporal/order positions, or similar role/entity patterns, but lack a validated skeleton-aligned source chain.
Perturbation-derived negatives are synthetic training-only variants of expert or pseudo-expert trajectories, created by replacement, insertion, deletion, time shift / shuffle, or entity-role swaps.
They are used for reward-learning comparison sets and train-split policy augmentation, and are not inserted into dev/test ranking pools as ordinary candidates.

\paragraph{Preference pairs.}
We create within-window preference supervision $(\tau^{\text{better}},\tau^{\text{worse}})$ according to skeleton-alignment quality.
When sufficient evidence is available, $\tau^{\text{better}}$ contains later-stage evidence such as \texttt{EXECUTE} and/or \texttt{OUTCOME} with fewer precedence violations.
In sparse settings such as RAMS, preferences are based on higher validated alignment score, role satisfaction, and evidence traceability, rather than requiring all late-stage steps to be observed.

\subsection{Minimal example window (anonymized)}
\label{sec:app_data_miniexample}

Below is a simplified illustrative excerpt (IDs and text are anonymized):
\begin{lstlisting}[caption={Minimal anonymized excerpt of a \texttt{traj\_pred.jsonl} trajectory record.},label={lst:ecpo-miniexample}]
traj_pred.jsonl (excerpt for candidate P_001):
{"window_id":"w_0001","candidate_id":"P_001","trajectory_id":"w_0001::P_001","events":[
 {"event_id":"e1",
 "etype_raw":"Process_start",
 "skeleton_hits":["PREP","PROBE"],
 "etype_primary":"PREP",
 "time":"2021-08-03",
 "order_index":0,
 "trigger":{"doc_id":"doc123","span":[120,156]},
   "arguments":[{"role":"Agent","entity_id":"P_001","doc_id":"doc123","span":[90,110]}]},
    {"event_id":"e2",
   "etype_raw":"Statement",
   "skeleton_hits":["PREP","PROBE"],
   "etype_primary":"PREP",
   "time":"2021-08-05",
   "order_index":1,
   "trigger":{"doc_id":"doc127","span":[33,61]},
   "arguments":[{"role":"Agent","entity_id":"P_001","doc_id":"doc127","span":[10,18]},
                {"role":"Target","entity_id":"ORG_007","doc_id":"doc127","span":[80,95]}]}
]}
\end{lstlisting}

\paragraph{Schema files.}
We provide JSON Schemas under \texttt{schemas/} to validate each record type. The evaluation scripts use the same schemas for feasibility checking and audit metrics.

\begin{lstlisting}[caption={SFT and RL-prompt formatting and column mapping (artifact metadata).},label={lst:format_mapping}]
maven_sft.jsonl and rams_sft.jsonl: Alpaca formatting
  prompt  -> instruction
  query   -> input
  response-> output
  system  -> system
  history -> history

rl_demos.jsonl and rams_rl_demos.jsonl: ECPO formatting (prompt+response)
  prompt  -> prompt
  response-> response

rl_prompts.jsonl and rams_rl_prompts.jsonl: ECPO formatting (prompt only)
  prompt  -> prompt

\end{lstlisting}

\subsection{Event type and trigger distributions (MAVEN-ERE)}
\label{sec:app_mavenere_dist}

To make the long-tail property explicit, we list the most frequent source-normalized event types and trigger surface forms in the processed MAVEN-ERE instantiation (total source-normalized event records: 91{,}222).
Table~\ref{tab:app_mavenere_top20_event_types} and Table~\ref{tab:app_mavenere_top20_triggers} report the Top-20 distributions.

\begin{table*}[t]
\centering
\caption{Top-20 source-normalized event types in MAVEN-ERE annotations after preprocessing (by frequency; total source-normalized event records = 91{,}222).}
\label{tab:app_mavenere_top20_event_types}
\small
\begin{adjustbox}{max width=\linewidth}
\begin{tabular}{r l r r | r l r r}
\hline
Rank & Type & Count & \% & Rank & Type & Count & \% \\
\hline
1  & \texttt{Causation}            & 3297 & 3.61 & 11 & \texttt{Statement}            & 1677 & 1.84 \\
2  & \texttt{Process\_start}       & 3191 & 3.50 & 12 & \texttt{Conquering}           & 1661 & 1.82 \\
3  & \texttt{Attack}               & 3059 & 3.35 & 13 & \texttt{Self\_motion}         & 1565 & 1.72 \\
4  & \texttt{Hostile\_encounter}   & 3057 & 3.35 & 14 & \texttt{Arriving}             & 1521 & 1.67 \\
5  & \texttt{Catastrophe}          & 2749 & 3.01 & 15 & \texttt{Destroying}           & 1397 & 1.53 \\
6  & \texttt{Competition}          & 2625 & 2.88 & 16 & \texttt{Coming\_to\_be}        & 1387 & 1.52 \\
7  & \texttt{Motion}               & 2533 & 2.78 & 17 & \texttt{Bodily\_harm}          & 1235 & 1.35 \\
8  & \texttt{Killing}              & 1839 & 2.02 & 18 & \texttt{Death}                & 1176 & 1.29 \\
9  & \texttt{Process\_end}         & 1705 & 1.87 & 19 & \texttt{Damaging}             & 1157 & 1.27 \\
10 & \texttt{Social\_event}        & 1689 & 1.85 & 20 & \texttt{Creating}             & 1145 & 1.26 \\
\hline
\multicolumn{8}{l}{Top-20 coverage: 39{,}665 source-normalized event records (43.48\%).} \\
\hline
\end{tabular}
\end{adjustbox}
\end{table*}

\begin{table*}[t]
\centering
\caption{Top-20 trigger surface forms in processed MAVEN-ERE source-normalized records (by frequency; total source-normalized event records = 91{,}222).}
\label{tab:app_mavenere_top20_triggers}
\small
\begin{adjustbox}{max width=\linewidth}
\begin{tabular}{r l r r | r l r r}
\hline
Rank & Trigger text & Count & \% & Rank & Trigger text & Count & \% \\
\hline
1  & \texttt{killed}      & 948 & 1.04 & 11 & \texttt{hurricane}  & 484 & 0.53 \\
2  & \texttt{began}       & 886 & 0.97 & 12 & \texttt{fought}     & 452 & 0.50 \\
3  & \texttt{held}        & 757 & 0.83 & 13 & \texttt{destroyed}  & 402 & 0.44 \\
4  & \texttt{took place}  & 732 & 0.80 & 14 & \texttt{battle}     & 397 & 0.44 \\
5  & \texttt{caused}      & 691 & 0.76 & 15 & \texttt{tour}       & 397 & 0.44 \\
6  & \texttt{occurred}    & 649 & 0.71 & 16 & \texttt{made}       & 382 & 0.42 \\
7  & \texttt{became}      & 617 & 0.68 & 17 & \texttt{reported}   & 378 & 0.41 \\
8  & \texttt{damage}      & 586 & 0.64 & 18 & \texttt{ended}      & 370 & 0.41 \\
9  & \texttt{storm}       & 578 & 0.63 & 19 & \texttt{war}        & 361 & 0.40 \\
10 & \texttt{attack}      & 515 & 0.56 & 20 & \texttt{Battle}     & 353 & 0.39 \\
\hline
\multicolumn{8}{l}{Top-20 coverage: 10{,}935 source-normalized event records (11.99\%).} \\
\hline
\end{tabular}
\end{adjustbox}
\end{table*}

\subsection{Dataset statistics (MAVEN-ERE)}
\label{sec:app_mavenere_stats}

\paragraph{Summary.}
After preprocessing, the MAVEN-ERE reference view contains 3{,}623 documents and 91{,}222 source-normalized event records.
The fixed-extractor model-input / reward-learning view contains 7{,}245 trajectory records over 3{,}623 unique source-level candidate/entity IDs before window expansion.
These records are consumed by the ranking / reward-learning code paths and may include train-only hard/perturbation negatives in addition to window-scoped ranking trajectories.
Trajectory lengths are measured over trajectory-expanded predicted event mentions and range from 1 to 128 events (median 21; mean 24.7).
Although the construction permits empty trajectories, none occur in the current MAVEN-ERE release; empty trajectories would still be retained and scored by missing-step features if they were present.
When enabling the within-window KG/TGN encoder, the resulting window graphs are lightweight (3--74 nodes and 0--110 edges; medians 14 and 16).
For reward learning, we tag 755 trajectories as expert/pseudo-expert demonstrations and 2,716 as hard/perturbed negatives, and we create 3,622 within-window preference pairs (covering 7,244 unique trajectory IDs).

\begin{table*}[t]
\centering
\caption{Statistics of the processed MAVEN-ERE instantiation.
Rows marked as source-normalized records describe normalized dataset annotations; rows marked as model-input records describe fixed-extractor predicted trajectories consumed by ranking methods.
\textit{Note:} skeleton-step hit counts are computed over window-scoped trajectory instances; overlapping windows may duplicate evidence.
The current MAVEN-ERE release contains no empty serialized trajectories, so the minimum trajectory length is 1 even though the construction allows empty trajectories in general.}
\label{tab:mavenere_stats}
\begin{tabularx}{\linewidth}{YY}
\hline
Statistic & Value \\
\hline
Documents & 3,623 \\
Source-normalized event records & 91,222 \\
Unique source-level candidate/entity IDs before window expansion & 3,623 \\
Total trajectory records exposed to ranking / reward learning & 7,245 \\
Trajectory record subsets & unlabeled candidate pool 3,774; expert/pseudo demos 755; hard/perturbed negatives 2,716 \\
Preference pairs & 3,622 \\
Trajectory length (\#trajectory-expanded predicted event mentions) & min 1; median 21; mean 24.7; max 128 \\
Preference coverage & 7,244 trajectory IDs \\
Ranking-window KG nodes & min 3; median 14; mean 15.5; max 74 \\
Ranking-window KG edges & min 0; median 16; mean 19.4; max 110 \\
Role mentions in source-normalized records & Agent 168,431; Target 137,720; Context 10,398 \\
Skeleton-step hits after window / trajectory expansion & PREP 143,859; PROBE 11,876; EXECUTE 31,340; OUTCOME 3,630 \\
Top-20 event types & 39,665 (43.5\%) \\
Top-20 triggers & 10,935 (12.0\%) \\
\hline
\end{tabularx}
\end{table*}

\subsection{Dataset statistics (RAMS)}
\label{sec:app_rams_stats}

\paragraph{Summary.}
The RAMS reference view contains 9{,}124 documents and 9{,}124 source-normalized event records, one per document in \texttt{rams\_event.jsonl}.
The fixed-extractor model-input / reward-learning view contributes 9{,}386 unique predicted event records over 8{,}730 unique source-level candidate/entity IDs before window expansion.
It contains 8,730 ranking-window candidate trajectories and 134 additional train-only hard/perturbation negatives, for 8,864 total trajectory records.
Trajectory sparsity is measured over the ranking-window candidate trajectories and averages $1.07$ trajectory-expanded predicted event mentions.
For preference-based supervision, we provide 134 within-window preference pairs (\texttt{rams\_pairs.jsonl}) covering 268 unique trajectory IDs.
For supervised training, we provide 9,124 Alpaca-format SFT records with train/dev/test splits of 7,329/924/871 (\texttt{rams\_sft\_train/dev/test.jsonl}).
For RL training, we provide 654 window-level RL prompt records (\texttt{rams\_rl\_prompts.jsonl}).
Key statistics are summarized in Appendix Table~\ref{tab:rams_stats}.

\begin{table*}[t]
\centering
\caption{Statistics of the processed RAMS instantiation.}
\label{tab:rams_stats}
\begin{tabularx}{\linewidth}{YY}
\hline
Statistic & Value \\
\hline
Documents & 9,124 \\
Source-normalized event records (\texttt{rams\_event.jsonl}) & 9,124 \\
Unique predicted event records contributing to model inputs & 9,386 \\
Unique source-level candidate/entity IDs before window expansion & 8,730 \\
Ranking-window candidate trajectories (\texttt{rams\_traj.jsonl}) & 8,730 \\
Additional train-only hard/perturbation negatives & 134 \\
Total trajectory records exposed to ranking / reward learning & 8,864 \\
Average trajectory length (\#trajectory-expanded predicted event mentions) & 1.07 \\
Preference pairs (\texttt{rams\_pairs.jsonl}) & 134 \\
Preference coverage (unique trajectory IDs) & 268 \\
SFT records (\texttt{rams\_sft.jsonl}) & 9,124 \\
SFT splits (\texttt{train/dev/test}) & 7,329 / 924 / 871 \\
RL prompt records (\texttt{rams\_rl\_prompts.jsonl}) & 654 \\
\hline
\end{tabularx}
\end{table*}

\section{Method and Baseline Implementation Details}
\label{sec:app_method_details}
\subsection{Training-loop walkthrough}
\label{sec:app_reading_fig}

The policy proposes Top-$K$ candidate rankings and span-grounded certificates for each window.
A feasible sampler enforces the machine-parseable schema and evidence validity constraints, reducing invalid action space.
The MaxEnt-style reward model scores plan consistency at the trajectory level based on skeleton alignment, preference supervision, and KG/TGN-derived temporal features.
The shaped reward combines reward-model utility with KL regularization to a frozen reference policy.
Within each window, we aggregate a group of sampled outputs with a group-normalized baseline to compute stable GRPO-style advantages for policy optimization.
The optional contrastive diagnostic is not part of the reported main objective.
Finally, the reward-model update refines the trajectory-level reward using expert/pseudo-expert trajectories, preference pairs, hard/near-neighbor negatives, and optionally policy-mined hard cases.
In each policy-optimization phase, the reward model is kept frozen and is only updated periodically between phases; in the final stage we fully freeze the reward model for stable evaluation-time scoring.

\subsection{Event extractor and trajectory construction details}
\label{sec:app_extractor}

We treat event extraction as upstream processing but include details here because it materially affects trajectory quality.
We adopt a staged supervised pretraining curriculum:
\[
\begin{aligned}
\texttt{MAVEN} &\rightarrow \texttt{MAVEN-Arg} \rightarrow \texttt{MAVEN-ERE} \\
&\rightarrow \texttt{RAMS}.
\end{aligned}
\]
Here \texttt{MAVEN-Arg} denotes the intermediate curriculum stage in our extractor pipeline; we name it explicitly because it is an internal training stage rather than a separate evaluation benchmark.
For each target benchmark instantiation, document-level splits are fixed before extractor training.
When the curriculum reaches a target corpus, only training-split documents from that corpus are used.
Held-out dev/test documents and their event labels are never used to train the extractor.

\begin{lstlisting}[caption={Leakage-controlled data flow for one target corpus.},label={lst:data_flow}]
1. Split documents into train/dev/test by doc_id.
2. Train extractor only on train documents.
3. Run the fixed extractor on train/dev/test documents to produce predicted event records.
4. Aggregate predicted event records into window-scoped candidate trajectories.
5. Use train annotations to construct reward-learning positives, hard (near-neighbor) negatives, and preference pairs.
6. Use dev annotations only for model selection and diagnostic evaluation.
7. Use test annotations only for final labels and audit references.
8. Train RM and policy only with train-split supervision.
9. Evaluate all ranking methods on predicted dev/test trajectory inputs.
\end{lstlisting}

\paragraph{MAVEN-ERE SFT data scale.}
The full processed MAVEN-ERE source yields 91,222 possible instruction-tuning instances for trigger/argument extraction before splitting.
After document-level splitting, extractor optimization uses only the training-split instances.
Dev/test instances are not used for extractor optimization; they are used only to run the frozen extractor and to support downstream held-out evaluation.
The tokenized input length ranges from 216 to 2,499 (mean 596.9; median 580), and the output length ranges from 34 to 77 (mean 45.3; median 44).
All instances use a fixed system prompt and instruction prefix, with zero dialogue history turns.

\subsection{Skeleton alignment and certificate extraction details}
\label{sec:app_alignment}

This appendix details the dynamic-programming (DP) alignment between a plan skeleton and a candidate trajectory, the tie-breaking rules, and the deterministic extraction of \texttt{doc\_id:span}-grounded certificate chains.

\subsubsection{Local compatibility score}
\label{sec:app_alignment_match}

Given a skeleton step
$s_k=(\texttt{etype}_k,\mathcal{R}_k)$
and an observed event $e_t$, we compute a local match score using the stage-compatibility set
$\texttt{skeleton\_hits}(e_t)$:
\begin{equation}
m(k,t)=
\begin{cases}
\mathrm{role\_sat}(s_k,e_t)
-\lambda_{\mathrm{arg}}\cdot \mathrm{arg\_viol}(s_k,e_t),
& \text{if } \texttt{etype}_k\in\texttt{skeleton\_hits}(e_t),\\[2pt]
-\delta_{\mathrm{type}},
& \text{otherwise.}
\end{cases}
\end{equation}

\paragraph{Role satisfaction.}
We define
\begin{equation}
\mathrm{role\_sat}(s_k,e_t)=
\begin{cases}
1, & |\mathcal{R}_k|=0,\\[2pt]
\frac{1}{|\mathcal{R}_k|}
\sum_{r\in\mathcal{R}_k}
\mathbb{I}[r\ \text{is present in}\ \mathcal{A}_t],
& |\mathcal{R}_k|>0.
\end{cases}
\end{equation}
For non-empty required-role sets, this is the fraction of required roles instantiated by at least one argument in the event; when no role is required, the event receives full role satisfaction by definition.

\paragraph{Argument violations.}
If the skeleton step specifies role/entity constraints (for instance, requiring the \texttt{Agent} role to be the candidate entity), we define
\begin{equation}
\mathrm{arg\_viol}(s_k,e_t)=\sum_{r\in\mathcal{R}_k}\mathbb{I}[\text{role }r\text{ violates a constraint under }e_t].
\end{equation}
In our released implementation, local argument-violation checks include candidate-identity consistency for key roles and role-specific incompatibilities available from the current event record.
Precedence violations are not included in the local score $m(k,t)$; they are computed after backtracking from the matched step timestamps or order indices and are added as separate alignment-derived features.
We do not enforce additional entity-type constraints for roles in our main experiments (this switch is reserved for domain-specific deployments).

\subsubsection{DP recurrence and initialization}
\label{sec:app_alignment_dp}

Let $M$ be the number of skeleton steps and $T$ be the number of observed events in the trajectory.
We compute a DP table of size $(M{+}1)\times (T{+}1)$:
\begin{equation}
\mathrm{DP}[0,0]=0,\qquad
\mathrm{DP}[0,t]=-t\cdot\delta_{\mathrm{skip}},\qquad
\mathrm{DP}[k,0]=-k\cdot\delta_{\mathrm{miss}}.
\end{equation}
For $k\ge 1$ and $t\ge 1$:
\begin{equation}
\begin{aligned}
\mathrm{DP}[k,t] = \max\Big\{\, &
\mathrm{DP}[k-1,t-1] + m(k,t), \\
& \mathrm{DP}[k-1,t] - \delta_{\mathrm{miss}}, \\
& \mathrm{DP}[k,t-1] - \delta_{\mathrm{skip}}
\Big\}.
\end{aligned}
\end{equation}

This formulation allows (i) missing skeleton steps (``miss''), and (ii) spurious observed events (``skip'').

\paragraph{Tie-breaking.}
When the max is not unique, we use the following deterministic tie-break order:
\begin{enumerate}
    \item Prefer \textbf{match} ($k{-}1,t{-}1$) over miss/skip when its score is tied.
    \item Prefer \textbf{skip} ($k,t{-}1$) over miss ($k{-}1,t$) when tied, because skipping a noisy observed event preserves the possibility of matching the current skeleton step later.
    \item As a final tie-breaker, prefer smaller temporal gaps (if timestamps are available) to produce more temporally coherent alignments.
\end{enumerate}
We record the chosen transition to enable backtracking.

\subsubsection{Backtracking and step-to-evidence mapping}
\label{sec:app_alignment_backtrack}

Backtracking yields an alignment path and a mapping
\[
\pi_{\mathrm{align}}:\{1,\ldots,M\}\rightarrow \{1,\ldots,T\}\cup\{\bot\},
\]
where $\pi_{\mathrm{align}}(k)=t$ indicates step $k$ is supported by observed event $t$, and $\bot$ indicates the step is missing.

\paragraph{Matched-step serialization.}
A diagonal transition in the DP table is serialized as a matched certificate step only if the event is type-compatible with the skeleton step, i.e.,
$\texttt{etype}_k\in\texttt{skeleton\_hits}(e_t)$.
If a diagonal transition is selected only because the soft type-mismatch penalty is less costly than a miss/skip transition, the step is treated as missing for certificate serialization.
This prevents type-mismatched events from being cited as faithful evidence.

\paragraph{Evidence extraction.}
If $\pi_{\mathrm{align}}(k)=t$ and $e_t$ is type-compatible with $s_k$, we construct the step evidence set as:
\begin{equation}
\mathcal{E}_k=
\{\sigma_{t}^{\mathrm{trig}}\}\ \cup\
\{\sigma_{t}^{\mathrm{arg}}:\ (\texttt{role}\in \mathcal{R}_k)\ \wedge\ (\texttt{role},\texttt{entity\_id},\sigma_{t}^{\mathrm{arg}})\in \mathcal{A}_t\}.
\end{equation}
If multiple arguments exist for the same role, we keep at most one argument span per role to keep certificates compact.
Concretely, we select the argument mention with the smallest span length (ties broken by the smallest start offset), which makes certificate serialization deterministic.

\paragraph{Certificate serialization.}
We serialize window-level outputs using the same strict JSON schema shown in Listing~\ref{lst:ecpo-json-example}.
Concretely, a certificate is represented under \texttt{certificates[*].steps}, where each step object records
\texttt{step\_id}, \texttt{etype}, \texttt{matched}, the aligned \texttt{event\_id} (or \texttt{null}), and an
\texttt{evidence} list of \texttt{doc\_id:span} references (trigger/argument spans with optional \texttt{role}).

This schema is machine-parseable and supports automatic faithfulness checks.

\subsubsection{Alignment-derived statistics}
\label{sec:app_alignment_stats}

From the alignment we compute interpretable statistics used as reward features:
\begin{itemize}
    \item \textbf{Hits/Misses}: number of matched vs.\ missing steps.
    \item \textbf{Coverage}: compatible-hit rate
\[
\frac{
|\{k:\exists t,\ \pi_{\mathrm{align}}(k)=t
\ \wedge\
\texttt{etype}_k\in\texttt{skeleton\_hits}(e_t)\}|
}{M}.
\]
    \item \textbf{Skips}: number of observed events not aligned to any step.
    \item \textbf{Temporal gaps}: average time difference between consecutive matched steps (absolute time or normalized $t$).
    \item \textbf{Precedence diagnostics}: for each soft precedence edge $k_1\prec k_2$, we compare the matched events using their explicit timestamp or order metadata when both matched events are available. 
    A violation is recorded when the metadata order contradicts the precedence edge or when the order is marked uncertain; if the serialized event order is already the trusted timestamp order, this feature is typically zero and the remaining precedence signal comes from missing steps and temporal-gap features.
\end{itemize}

\subsubsection{Precedence constraints}
\label{sec:app_alignment_prec}

If the skeleton includes precedence relations $\prec$, we compute precedence diagnostics after alignment from the matched event timestamps or explicit order indices, not from the DP table index alone.
The reported experiments use a \textbf{soft} mode: violations, uncertain order relations, and large temporal gaps are added as post-alignment features in $\phi_{\text{skel}}$, while the DP recurrence itself remains monotone in the serialized event order.
A hard precedence variant would require augmenting the DP state or restricting diagonal transitions using previously matched predecessor steps; we do not use that variant in the main experiments because timestamps are noisy and many trajectories are partially observed.

\subsection{MaxEnt-style reward-learning objective and optimization details}
\label{sec:app_irl}

We present the full MaxEnt-style reward-learning objective used to learn the skeleton-conditioned trajectory-level reward
$R_\theta(\tau;S_w)=\theta^\top\phi(\tau;S_w)$.
The objective combines a window-wise maximum-entropy likelihood with a Bradley--Terry preference term.
When $S_w$ is clear, we omit it for readability.

For reward learning, let
\[
\mathcal{B}_w=\{\tilde{\tau}_{w,1},\ldots,\tilde{\tau}_{w,B_w}\}
\]
denote the training comparison set for window $w$.
It contains positive demonstrations, same-window negatives, and perturbation-derived negatives.
We define:
\begin{equation}
P_\theta(\tau \mid w)
=
\frac{\exp\left(R_\theta(\tau;S_w)\right)}
{\sum_{j=1}^{B_w}\exp\left(R_\theta(\tilde{\tau}_{w,j};S_w)\right)}.
\end{equation}
Here $B_w$ denotes the number of trajectories in the reward-learning comparison set and should not be confused with $N_w=|C_w|$, the candidate-pool size used for ranking.

Let $\mathcal{P}_w^+$ denote the set of expert or pseudo-expert positive trajectories in window $w$.
The window-wise MaxEnt likelihood is
\begin{equation}
\mathcal{L}_{\mathrm{me}}(\theta)
=
\sum_{w:\,|\mathcal{P}_w^+|>0}
\frac{1}{|\mathcal{P}_w^+|}
\sum_{\tau^+\in \mathcal{P}_w^+}
\log P_\theta(\tau^+\mid w).
\end{equation}

We additionally use within-window preference pairs
\[
\mathcal{Q}=\{(w_q,\tau_q^+,\tau_q^-)\},
\]
where $w_q$ identifies the ranking window from which the pair is constructed.
The Bradley--Terry preference likelihood term is
\begin{equation}
\mathcal{L}_{\mathrm{pair}}(\theta)
=
\sum_{(w_q,\tau_q^+,\tau_q^-)\in\mathcal{Q}}
\log \sigma\!\left(
R_\theta(\tau_q^+;S_{w_q})
-
R_\theta(\tau_q^-;S_{w_q})
\right).
\end{equation}

The final reward-learning objective is
\begin{equation}
\mathcal{L}(\theta)
=
\mathcal{L}_{\mathrm{me}}(\theta)
+
\alpha_{\mathrm{pair}}\mathcal{L}_{\mathrm{pair}}(\theta)
-
\frac{\lambda}{2}\|\theta\|_2^2.
\end{equation}

The MaxEnt component has the usual feature-matching gradient:
\begin{equation}
\nabla_\theta \log P_\theta(\tau_w^+ \mid w)
=
\phi(\tau_w^+;S_w)
-
\sum_{j=1}^{B_w}
P_\theta(\tilde{\tau}_{w,j}\mid w)\,
\phi(\tilde{\tau}_{w,j};S_w).
\end{equation}
The preference component adds the Bradley--Terry gradient induced by
$R_\theta(\tau^+)-R_\theta(\tau^-)$.

\subsection{Deterministic evidence-only verifier and certificate reward details}
\label{sec:app_evidence_verifier}
The ECPO version uses the deterministic verifier $\Gdet$ as the evidence-only recovery module. $\Gdet$ is deliberately label-free: it is not trained, not prompted with gold labels, and not tuned on held-out relevance judgments. Before reconstruction, each generated certificate is converted into claim-stripped bundles by removing candidate identifiers, rank positions, policy scores, and any direct candidate claim. The verifier receives only those bundles, the public window skeleton, exposed document/span metadata, and frozen candidate trajectories. This makes \EvidCons{} and \CertNDCG{} less vulnerable to the criticism that a learned verifier has become a second hidden ranker.

For each claim-stripped bundle $b$ and every candidate $c\in C_w$, $\Gdet$ computes five sub-scores:
\begin{enumerate}
    \item \textbf{Traceability}: fraction of cited spans that pass document, span-bound, event-signature, and trajectory-overlap checks against $\tau_w(c)$.
    \item \textbf{Step coverage}: fraction of skeleton steps with at least one valid traceable evidence item for candidate $c$.
    \item \textbf{Role satisfaction}: fraction of required roles supported by traceable argument spans with normalized role agreement.
    \item \textbf{Precedence consistency}: fraction of supported step pairs that respect available order or timestamp constraints.
    \item \textbf{Bad-span penalty}: duplicate, off-window, non-overlapping, wrong-role, type-incompatible, or multi-candidate-ambiguous evidence citations.
\end{enumerate}
The default score is the fixed weighted sum defined in Section~\ref{sec:method_det_verifier}. A simple unweighted version can be reported as a sensitivity check. Candidate assignments must pass both a support threshold and a margin over the second-best candidate. Ambiguous bundles receive no recovered-candidate credit. Ties are broken deterministically by bundle order and window-local candidate order, never by source-level IDs or labels.

\begin{table}[h]
\centering
\caption{Deterministic evidence-only verifier component sensitivity. Values report the change in \CertNDCG{} relative to the default verifier unless otherwise noted.}
\label{tab:det_verifier_components}
\small
\begin{adjustbox}{max width=\linewidth}
\begin{tabular}{l l c c}
\toprule
Component & Input visible to $\Gdet$ & Used in reward? & Sensitivity result \\
\midrule
Traceability & cited \texttt{doc\_id:span}, event IDs, candidate trajectory & yes & $\Delta\CertNDCG=-0.10$ \\
Step coverage & cited step labels, skeleton steps & yes & $\Delta\CertNDCG=-0.07$ \\
Role satisfaction & normalized role labels and argument spans & yes & $\Delta\CertNDCG=-0.06$ \\
Precedence consistency & event order/timestamp, soft precedence edges & yes & $\Delta\CertNDCG=-0.02$ \\
Bad-span penalty & validator failures and duplicate/off-window citations & yes & $\Delta\CertNDCG=-0.08$ \\
Support-margin assignment & support threshold, second-best margin, MWM assignment & yes & $\Delta\CertNDCG=-0.04$ \\
Unweighted score variant & same inputs with equal weights & diagnostic & $\Delta\CertNDCG=-0.02$ \\
Rank-aware recovery variant & same inputs, rank-sensitive agreement & diagnostic & $\mathbf{\Delta\CertNDCG=+0.01}$; ranking unchanged \\
\bottomrule
\end{tabular}
\end{adjustbox}
\end{table}

For a generated output $y=(A_y,E_y)$, evidence-only reconstruction produces slot-level recoveries $\hat a_j=\Gdet(E_y[j],w)$ after claim stripping and maximum-weight assignment. The default cycle reward uses slot-wise recovery $\frac{1}{K_w}\sum_j\mathbb{I}[\hat a_j=a_j]$; set-overlap and stricter rank-aware variants are diagnostics. The certificate reward combines validator pass/fail, step coverage, traceability, and role support. During ablations, removing $r_{\mathrm{cycle}}$ tests whether ECPO reduces to strict-output RL, while removing $r_{\mathrm{cert}}$ tests whether evidence-only reconstruction alone is sufficient.

\subsection{RL optimization and feasibility checking details}
\label{sec:app_rl}

This appendix specifies the machine-parseable output schema, feasibility checking rules, and RL optimization details (GRPO) for Top-$K$ ranking with auditable evidence.

\subsubsection{Output schema}
\label{sec:app_rl_schema}

The policy outputs a single JSON object per window.
We require strict JSON to enable deterministic parsing and audit checks.
A schematic minimal interface is shown in Listing~\ref{lst:ecpo-json-min}.
An extended example is provided in Listing~\ref{lst:ecpo-json-example}; the released JSON Schemas under \texttt{schemas/} are the source of truth for validation.
In brief, the object must include \texttt{window\_id}, \texttt{topk}, and a position-aligned \texttt{certificates} array with step-aligned \texttt{doc\_id:span} evidence. The certificate object itself does not contain \texttt{candidate\_id}; candidate identity is supplied only by the matching rank position in \texttt{topk} during feasibility validation.

The \texttt{topk} list must contain $K_w=\min(K,N_w)$ distinct candidate IDs from the window candidate pool.
Each certificate must use only \texttt{doc\_id:span} references whose \texttt{doc\_id} belongs to the window evidence scope. For feasibility, cited spans are checked against the candidate occupying the same rank position in \texttt{topk}; for evidence-only recovery, the same bundle is scored candidate-blind against the full roster.
In the released benchmark, compact snippets are a preprocessing and prompt-serialization device; feasibility is enforced through window membership, span bounds, and trajectory traceability rather than a separate snippet-bound check.

\subsubsection{Policy output JSON Schema (excerpt)}
\label{sec:app_rl_output_jsonschema}

We include a compact JSON-Schema excerpt for the per-window policy output to make the auditable interface explicit and machine-checkable within the paper. 
The accompanying artifact contains the full schema file (including stricter constraints and shared definitions).

\begin{lstlisting}[caption={JSON Schema excerpt for the auditable per-window policy output.},label={lst:schema-policy-output}]
{
  "$schema": "https://json-schema.org/draft/2020-12/schema",
  "title": "ECPORankOutput",
  "type": "object",
  "additionalProperties": false,
  "required": ["window_id", "topk", "certificates"],
  "properties": {
    "window_id": {"type": "string"},
    "topk": {
      "type": "array",
      "items": {"type": "string"},
      "minItems": 1,
      "uniqueItems": true
    },
    "certificates": {
      "type": "array",
      "items": {
        "type": "object",
        "additionalProperties": false,
        "required": ["steps"],
        "properties": {
          "steps": {
            "type": "array",
            "items": {
              "type": "object",
              "additionalProperties": false,
              "required": ["step_id", "etype", "matched", "event_id", "evidence"],
              "properties": {
                "step_id": {"type": "string"},
                "etype": {
                  "type": "string",
                  "enum": ["PREP", "PROBE", "EXECUTE", "OUTCOME"]
                },
                "matched": {"type": "boolean"},
                "event_id": {"type": ["string", "null"]},
                "evidence": {
                  "type": "array",
                  "items": {
                    "type": "object",
                    "additionalProperties": false,
                    "required": ["doc_id", "span", "kind"],
                    "properties": {
                      "doc_id": {"type": "string"},
                      "span": {"$ref": "#/$defs/span"},
                      "kind": {"type": "string", "enum": ["trigger", "arg"]},
                      "role": {"type": "string"}
                    }
                  }
                },
                "notes": {"type": "string"}
              }
            }
          }
        }
      }
    }
  },
  "$defs": {
    "span": {
      "anyOf": [
        {"type": "string", "pattern": "^\\d+-\\d+$"},
        {
          "type": "array",
          "items": {"type": "integer"},
          "minItems": 2,
          "maxItems": 2
        }
      ]
    }
  }
}
\end{lstlisting}

\paragraph{Notes.}
Some constraints (including exact Top-$K$ length, candidate-id membership in the window pool, position alignment between \texttt{topk} and \texttt{certificates}, skeleton-step order, \texttt{event\_id} membership in the rank-aligned candidate trajectory, conditional role requirements for argument evidence, \texttt{doc\_id} existence, and span bounds) are enforced by the deterministic validator described in Appendix~\ref{sec:app_rl_validator}, rather than by JSON Schema alone.

\subsubsection{Feasibility checking}
\label{sec:app_rl_validator}

We implement a deterministic validator $\mathrm{Valid}(y,w)\in\{0,1\}$.
A generated output $y$ is feasible iff it satisfies all of the following:

\begin{enumerate}
    \item \textbf{Schema validity}: strict JSON parsing; required keys (\texttt{window\_id}, \texttt{topk}, \texttt{certificates}); correct types.
    \item \textbf{Candidate id validity}: each id in \texttt{topk} belongs to the window's \texttt{candidate\_ids}; no duplicates; length exactly $K_w=\min(K,N_w)$.

    \item \textbf{Certificate coverage}: \texttt{certificates} must have the same length as \texttt{topk}. The $j$-th certificate is the claimed support for the $j$-th \texttt{topk} candidate during validation, but the certificate object itself contains no \texttt{candidate\_id}. Each certificate must contain one step object for every \texttt{step\_id} in the window skeleton $S_w$; missing or unsupported steps are serialized with \texttt{matched:false}, \texttt{event\_id:null}, and an empty \texttt{evidence} list.
    \item \textbf{Skeleton-step consistency}: for each returned candidate position, the set and order of \texttt{certificates[*].steps[*].step\_id} must match the skeleton steps in $S_w$, and each serialized \texttt{etype} must equal the corresponding skeleton-step stage label.
    A step marked \texttt{matched:true} must provide a non-null \texttt{event\_id} and cite at least one valid evidence span.
    A step marked \texttt{matched:false} must have \texttt{event\_id:null} and no cited evidence.

    \item \textbf{Document validity}: each cited \texttt{doc\_id} exists in the window \texttt{doc\_ids}.
    Window exposure for compact snippets or document-local context units is enforced indirectly through the trajectory-traceability rule rather than by a separate snippet-bound validator.
    \item \textbf{Span bounds}: each span satisfies $0\le l < r \le L_{\texttt{doc\_id}}$, where $L_{\texttt{doc\_id}}$ is stored in \texttt{doc\_meta.jsonl}. 
    If \texttt{span} is serialized as a string \texttt{"l-r"}, the validator parses it into integers $(l,r)$ and interprets it as the same half-open interval $[l,r)$.
    \item \textbf{Event-id and trajectory traceability}: for feasibility validation, if \texttt{event\_id} is non-null, it must identify an event in the trajectory $\tau_{w,a_k}$ of the candidate occupying the same rank position in \texttt{topk}.
    Each cited trigger span must overlap the trigger span of that event.
    Each cited argument span must overlap one of that event's argument spans.
    If \texttt{kind="arg"}, the evidence object must include a normalized \texttt{role}, and the overlapped argument mention must have the same normalized role.
    If \texttt{kind="trigger"}, any supplied \texttt{role} field is ignored by the validator.
\end{enumerate}

\paragraph{Repair vs.\ resampling during training.}
During policy optimization, when a sampled output violates constraints, we support two strategies:
(i) \textbf{reject-resample}: discard and resample until feasible under a bounded retry budget, or
(ii) \textbf{repair}: apply deterministic, label-free edits such as dropping duplicate IDs, truncating extra IDs, or removing invalid evidence entries, and then keep the repaired output while applying a format penalty.
Our training implementation uses reject-resample for severe schema failures and label-free repair for minor issues.

At evaluation time, each method produces one final output per window.
Unless explicitly stated, no repair is applied before computing feasibility and auditability metrics.

\subsubsection{Reward callback and ECPO shaping}
\label{sec:app_rl_reward}

For each sampled output $y=(A_y,E_y)$, we first attempt strict JSON parsing and deterministic validation.
If $y$ cannot be parsed or contains no valid candidate prefix, ranking utility, certificate utility, and cycle utility are set to zero and the output receives invalid and missing-prefix penalties.
For a parsed output, the unregularized ECPO reward is
\begin{equation}
\begin{aligned}
r_{\mathrm{base}}(y,w)=&
\bar r_{\mathrm{rank}}(y,w)
+\lambda_{\mathrm{cert}}r_{\mathrm{cert}}(y,w)
+\lambda_{\mathrm{cycle}}r_{\mathrm{cycle}}(y,w)\\
&-\xi_{\mathrm{inv}}\mathbb{I}_{\mathrm{invalid}}(y,w)
-\xi_{\mathrm{miss}}\max(0,K_w-L_y).
\end{aligned}
\end{equation}
Here $\bar r_{\mathrm{rank}}$ is the window- or group-normalized discounted listwise reward defined in Section~\ref{sec:method}. The certificate term rewards span validity, trajectory traceability, and skeleton-step coverage under the deterministic validator. The cycle term runs the fixed deterministic evidence-only verifier $\Gdet$ on claim-stripped cited evidence and compares the reconstructed candidate set $\widehat{A}_E$ with the policy-selected set $A_y$.
We then apply token-level KL regularization with respect to the frozen reference policy:
\begin{equation}
r_{\mathrm{ECPO}}'(y,w)=
r_{\mathrm{base}}(y,w)
-\beta_{\mathrm{KL}}\widehat{\mathrm{KL}}_{\mathrm{tok}}(y,x_w).
\end{equation}
The per-token KL estimate is
\begin{equation}
\widehat{\mathrm{KL}}_{\mathrm{tok}}(y,x_w)
=
\frac{1}{|y|}
\sum_{t=1}^{|y|}
\Big(
\log \pi_\varphi(y_t\mid y_{<t},x_w)
-\log \pi_{\mathrm{ref}}(y_t\mid y_{<t},x_w)
\Big).
\end{equation}
When computing group advantages, $r_{\mathrm{ECPO}}'(y,w)$ is treated as a scalar reward with stop-gradient through sampled log-probability terms. Ablations set $\lambda_{\mathrm{cycle}}=0$, $\lambda_{\mathrm{cert}}=0$, or remove the listwise reward to isolate the contribution of each objective component.

\subsubsection{GRPO optimization}
\label{sec:app_rl_ppo}

We optimize $\pi_\varphi$ with GRPO-style group normalization~\citep{shao_deepseekmath_2024}.
For a group of $G$ sampled outputs from the same window, we compute
\begin{equation}
A^{(g)}=
\frac{r_{\mathrm{ECPO}}'(y^{(g)},w)-\mu_w}{\hat{\sigma}_w+\epsilon_{\mathrm{adv}}},
\qquad
\mu_w=\frac{1}{G}\sum_{g=1}^{G}r_{\mathrm{ECPO}}'(y^{(g)},w),
\end{equation}
where $\hat{\sigma}_w$ is the within-group standard deviation.
We do not use an additional contrastive loss in the main experiments; the optional best-vs-worst contrastive diagnostic is reported only as an auxiliary analysis.
The policy update uses the clipped likelihood-ratio objective
\begin{equation}
\begin{aligned}
\mathcal{J}(\varphi)=
\mathbb{E}_{w,g}\!\Big[
\min\!\Big(
\rho^{(g)}A^{(g)},
\mathrm{clip}(\rho^{(g)},1-\epsilon_{\mathrm{clip}},1+\epsilon_{\mathrm{clip}})A^{(g)}
\Big)
\Big],
\end{aligned}
\end{equation}
where $\rho^{(g)}=\pi_\varphi(y^{(g)}\mid x_w)/\pi_{\mathrm{old}}(y^{(g)}\mid x_w)$.
In practice, the reward computation is implemented as an external callback that parses $y$ and calls the reward model.
We cap output length and enforce deterministic decoding parameters during evaluation.

\paragraph{KL coefficient control.}
We initialize $\beta_{\mathrm{KL}}$ on the development set and adapt it with a simple multiplicative controller to keep the empirical KL close to a target; $\pi_{\mathrm{ref}}$ remains frozen throughout RL:
\begin{equation}
\beta_{\mathrm{KL}} \leftarrow \beta_{\mathrm{KL}} \cdot \exp\Big(\kappa\cdot (\widehat{\mathrm{KL}}_{\mathrm{tok}}-\mathrm{KL}_{\mathrm{target}})\Big),
\end{equation}
where $\kappa$ is a small step size.
In the final phase, we freeze both the reward model and the reference policy.

\paragraph{Logging.}
We log: (i) mean shaped reward, (ii) normalized RM utility, certificate utility, cycle utility, and KL separately, (iii) feasibility rate, and (iv) Top-$K$ ranking metrics on the dev set.

\subsection{Inference and score fusion}
\label{sec:app_inference_fusion}
In the main results, we report the single-pass policy decoding outputs.
We include score fusion here as an auxiliary analysis variant that can reduce sampling variance when multiple policy samples are available.
At test time, for each candidate trajectory $\tau_{w,j}$ we compute an RM score $r_j=R_\theta(\tau_{w,j};S_w)$.
We compute a policy-based preference signal $\ell_j$ by sampling $S_{\mathrm{samp}}$ Top-$K$ outputs from the policy and accumulating a rank-discounted inclusion score:
\begin{equation}
\ell_j=\frac{1}{S_{\mathrm{samp}}}
\sum_{s=1}^{S_{\mathrm{samp}}}
\sum_{k=1}^{K_w}
\gamma^{k-1}\,
\mathbb{I}\!\left[a_k^{(s)}=c_{w,j}\right],
\qquad K_w=\min(K,N_w).
\end{equation}
We then fuse the two signals after window-wise $z$-score normalization:
\begin{equation}
\mathrm{score}_j
= \alpha \cdot \mathrm{z}(r_j)
+ (1-\alpha)\cdot \mathrm{z}(\ell_j),
\end{equation}
where $\mathrm{z}(u)=\frac{u-\mu_w}{\sigma_w+\epsilon}$ is computed over candidates within the same window $w$.
We return the Top-$K_w$ candidates by $\mathrm{score}_j$, where $K_w=\min(K,N_w)$.
For each returned candidate, we run skeleton alignment to produce a structured certificate and output \texttt{doc\_id:span}-grounded evidence per step.

\subsection{Knowledge graph and temporal encoding details}
\label{sec:app_tgn}

Within a window $w$, we build a knowledge graph $G_w$ that connects entities, events, and document spans.
Nodes include entity nodes, event nodes, and mention/span nodes; edges include argument links with role labels, coreference/entity links, and temporal edges.

The TGN encoder is trained on training-split graphs constructed from fixed-extractor predicted trajectories with a self-supervised temporal-edge objective, and then frozen before reward learning.
The frozen encoder is applied to dev/test window graphs only in forward mode to compute cached graph features; no dev/test graph edges, labels, or generated certificates are used to update encoder parameters.
For an observed temporal edge $(u,v,t)$, where $t$ is the event timestamp or document-order index, and a negative endpoint $v^-\sim q(\cdot)$ sampled uniformly from same-window nodes of the same broad node type, we optimize:
\begin{equation}
\mathcal{L}_{\mathrm{TGN}}
=
-\sum_{(u,v,t)\in E_{\mathrm{train}}}
\left[
\log \sigma(h_u(t)^\top h_v(t))
+
\sum_{v^-\sim q(\cdot)}
\log \sigma(-h_u(t)^\top h_{v^-}(t))
\right].
\end{equation}
The MaxEnt-style reward-learning objective therefore optimizes only the linear reward weights $\theta$ over cached trajectory features.
In the ablation without KG/TGN, we set $\phi_{\mathrm{graph}}(\tau)=\mathbf{0}$ while keeping all other features and training settings unchanged.

To encode temporal interactions, we apply a simplified TGN-style temporal message passing scheme:
\begin{equation}
\begin{aligned}
m_{u}(t) &=
\sum_{(v,t')\in \mathcal{N}_u(t)}
\exp\!\left(-\lambda (t-t')\right)\,
\psi\!\Big(
h_u(t^-),\, h_v(t'^-),\, x_{uv}(t')
\Big),
\end{aligned}
\end{equation}
\begin{equation}
\begin{aligned}
h_u(t) &= \mathrm{GRU}\!\left(m_u(t),\, h_u(t^-)\right),
\end{aligned}
\end{equation}
where $\lambda$ controls temporal decay and $\psi$ is a learnable message function.
We pool event-related node states to obtain a graph representation $g(\tau)$, which is concatenated into $\phi_{\text{graph}}(\tau)$ in the reward model.

\subsection{Complexity and implementation notes}
\label{sec:app_complexity}

Skeleton alignment for one trajectory costs $O(MT)$ with dynamic programming, where $M$ is the skeleton length and $T$ is the trajectory length.
RM scoring is linear in feature dimension once $\phi(\tau)$ is cached.
The RL callback adds only parsing and RM scoring overhead per sampled response, making the training pipeline practical for large candidate pools.

\subsection{LLM training and decoding details}
\label{sec:app_llm_training}

All LLM policies use the same Qwen3-8B backbone.
We first apply supervised fine-tuning on ranking-policy demonstrations and then optimize with GRPO for the measured GRPO control; the ECPO policy uses the evidence-coupled objective described in the main text.
Unless otherwise stated, all LLM methods use the same maximum input length, maximum output length, prompt format, and strict JSON output schema.
The ranking-window prompt contains the \texttt{intent\_id}, the associated skeleton $S_w$, candidate IDs, and compact trajectory summaries or evidence pointers.
Prompt construction removes all evaluation-only fields, including positive candidate IDs, trajectory labels, and reference evidence chains.

\paragraph{GRPO baseline and optional DPO construction.}
Qwen3-8B-GRPO uses the same SFT initialization, reference policy, decoding budget, and output validator as ECPO.
It replaces the learned trajectory-level reward with a direct supervised task reward consisting of: (i) ranking-label utility computed from train-split window positives, (ii) schema and candidate-ID validity, and (iii) evidence-format validity.
It does not use $R_\theta(\tau;S_w)$ or the DP-alignment-based skeleton-fidelity bonus used by ECPO.

For completeness, an optional Qwen3-8B-DPO construction can be trained from output-level preference pairs derived from the same within-window trajectory preference source used by reward learning.
The difference is that ECPO first learns a trajectory-level reward and then optimizes a constrained policy, whereas this optional DPO construction directly optimizes the policy on paired JSON outputs.
For each pair $(\tau^{\text{better}},\tau^{\text{worse}})$, we construct a train-only controlled ranking prompt containing the two paired candidates plus a fixed set of same-window distractors.
These synthetic pair prompts are used only for DPO optimization and are not part of the held-out ranking-window evaluation set.
The preferred output ranks the better candidate above the worse candidate and cites DP-aligned evidence extracted from that candidate's fixed-extractor predicted trajectory; the dispreferred output reverses this pairwise ordering or omits key aligned predicted evidence while keeping the same JSON schema.
No source-normalized reference span is inserted into DPO outputs.
All non-paired distractors are kept in the same relative order across the preferred and dispreferred outputs.
Thus, the optional DPO construction uses the same preference source as reward learning, but applies it directly to policy outputs rather than to the trajectory-level reward model.

\begin{table}[t]
\centering
\caption{LLM training and decoding hyperparameters.}
\label{tab:llm_training_hparams}
\small
\begin{tabularx}{\linewidth}{lY}
\hline
Item & Value \\
\hline
Backbone & Qwen3-8B \\
Fine-tuning method & \texttt{QLoRA} \\
SFT epochs & 2 \\
SFT learning rate & $1\times10^{-4}$ \\
Batch size & effective batch size 128 \\
Max input length & 4{,}096 tokens \\
Max output length & 768 tokens \\
GRPO group size $G$ & 8 \\
GRPO clip range $\epsilon_{\mathrm{clip}}$ & 0.2 \\
KL coefficient / target & $\beta_{\mathrm{KL}}=0.05$; target KL $0.05$ nats/token \\
Reward weights $(\gamma,\lambda_{\mathrm{cert}},\lambda_{\mathrm{cycle}},\xi_{\mathrm{inv}},\xi_{\mathrm{miss}})$ & $\gamma=0.90$; $\lambda_{\mathrm{cert}}=0.5$; $\lambda_{\mathrm{cycle}}=1.0$; $\xi_{\mathrm{inv}}=2.0$; $\xi_{\mathrm{miss}}=2.0$ \\
DPO inverse-temperature $\beta_{\mathrm{DPO}}$ & 0.10 \\
Evaluation decoding & greedy (temperature 0) \\
Validator retry budget & train: 3; eval: 0 \\
Hardware & $4\times$A100 80GB \\
Extractor training wall-clock / GPU-hours & estimated 16 h per corpus-specific extractor / 64 GPU-h; 128 GPU-h aggregate \\
Reward-model training wall-clock / GPU-hours & estimated 2 h per corpus-seed / 8 GPU-h; 48 GPU-h aggregate \\
Policy training GPU-hours per seed (SFT/GRPO/DPO/ECPO) & estimated 16 / 56 / 24 / 72 GPU-h per corpus-seed (4 / 14 / 6 / 18 h on $4\times$A100) \\
Total implementation runs and aggregate GPU-hours & 2 extractor pipelines; $2\times3$ reward-model fits; measured SFT/GRPO/ECPO policy runs plus optional DPO runs; estimated 1{,}176 GPU-h aggregate \\
\hline
\end{tabularx}
\end{table}

\paragraph{Compute accounting.}
GPU-hours are estimated as wall-clock hours multiplied by the number of allocated A100 GPUs. The aggregate in Table~\ref{tab:llm_training_hparams} includes the extractor, reward-model, SFT, GRPO, ECPO, and optional DPO implementation runs: two corpus-specific extractor pipelines, two corpora $\times$ three reward-model seeds, and policy-method seeds for the implemented policy variants. It excludes CPU-only preprocessing/evaluation jobs and exploratory sweeps not used in the filled result tables.

\subsection{Baseline implementation details}
\label{sec:app_baselines_impl}

This appendix summarizes implementation details and tuning knobs for structured / non-LLM comparison methods. The main-result table reports completed measurements for the structured rankers, while the implementation notes below document how each non-LLM baseline was instantiated.

\paragraph{RM-only ranker.}
We score each candidate trajectory by the learned skeleton-conditioned reward model $R_\theta(\tau;S_w)$, trained with the MaxEnt-style and Bradley--Terry objective in Appendix~\ref{sec:app_irl}, and rank candidates by this score within each window.
This baseline does not generate certificates; for auditability comparisons we attach DP-alignment evidence chains post hoc using the same DP alignment procedure.

\paragraph{LP-Recognizer (operator counting).}
We compute DP skeleton-alignment statistics, including hits, misses, skipped events, role satisfaction, and soft precedence violations, and use a hand-crafted linear scoring function over these features.

\paragraph{Recognition-as-planning (GRaP-style).}
We cast the skeleton as a planning template and score each candidate by the minimum completion cost under partial observability, following recognition-as-planning formulations \citep{ramirez_plan_2009,ramirez_probabilistic_2010}.

\paragraph{Landmark-based goal recognition.}
We treat skeleton steps as landmarks and compute a coverage-and-order score with missing-step and precedence-violation penalties.

\paragraph{LambdaMART/LTR.}
We train a learning-to-rank model (LambdaMART) on trajectory features $\phi(\tau)$ with window-level list supervision \citep{burges_ranknet_2010}.
Hyperparameters are tuned on the dev set.

\paragraph{GraphRank.}
We build the within-window KG $G_w$ and compute candidate scores using personalized PageRank over entity/event nodes.

\begin{table*}[t]
\centering
\caption{Summary of structured / non-LLM baselines, inputs, and development-tuned hyperparameter ranges.}
\label{tab:baseline_impl_summary}
\small
\setlength{\tabcolsep}{4pt}
\begin{adjustbox}{max width=\linewidth}
\begin{tabular}{p{0.30\linewidth}p{0.30\linewidth}p{0.34\linewidth}}
\hline
Baseline & Input & Key knobs (recommended ranges) \\
\hline
RM-only ranker 
& $R_\theta(\tau;S_w)$ scoring; DP-alignment evidence chains attached for audit metrics 
& reward-model $\ell_2$ reg $\lambda\in[10^{-4},10^{-2}]$; feature set $\phi\in\{\phi_{\text{skel}},\phi_{\text{skel}}{+}\phi_{\text{arg}},\phi_{\text{skel}}{+}\phi_{\text{arg}}{+}\phi_{\text{graph}}\}$ \\
LP-Recognizer 
& DP alignment stats 
& weights: $w_{\text{hit}}\in[0.5,2.0]$, $w_{\text{miss}}\in[0.5,2.5]$, $w_{\text{inv}}\in[0.2,2.0]$, role-bonus $\in[0,1]$ \\
GRaP-style 
& skeleton + observed steps 
& costs: $c_{\text{miss}}\in[1,6]$, $c_{\text{skip}}\in[0.5,4]$, precedence penalty $\in[0,5]$; solver A*/DP (beam $\in[8,64]$) \\
Landmark 
& skeleton landmarks + order 
& missing/ordering penalties $\in[0.5,5]$; order-temperature (if soft) $\in[0.5,2]$ \\
LambdaMART 
& trajectory features $\phi(\tau)$ 
& depth $\in[3,8]$, learning rate $\in[0.03,0.2]$, trees $\in[200,2000]$, min\_child\_weight $\in[1,20]$ \\
GraphRank 
& $G_w$ (entities/events/spans) 
& PPR damping $\alpha\in[0.75,0.95]$ \\
\hline
\end{tabular}
\end{adjustbox}
\end{table*}

\begin{table}[t]
\centering
\caption{Selected hyperparameters for structured / non-LLM baselines.
All values are tuned on the development split using NDCG@10.}
\label{tab:baseline_selected_hparams}
\small
\begin{tabularx}{\linewidth}{lY}
\hline
Baseline & Selected hyperparameters \\
\hline
RM-only ranker
& $\lambda=10^{-3}$; feature set $\phi_{\mathrm{skel}}+\phi_{\mathrm{arg}}+\phi_{\mathrm{graph}}$ \\
LP-Recognizer
& $w_{\mathrm{hit}}=1.0$; $w_{\mathrm{miss}}=1.5$; $w_{\mathrm{inv}}=1.0$; role bonus $=0.4$ \\
Recognition-as-planning
& $c_{\mathrm{miss}}=2$; $c_{\mathrm{skip}}=1$; precedence penalty $=1.5$; beam $=32$ \\
Landmark-based
& missing penalty $=2$; ordering penalty $=1.0$ \\
LambdaMART / LTR
& depth $=6$; trees $=1{,}000$; learning rate $=0.05$; min child weight $=5$ \\
GraphRank
& PPR damping $\alpha=0.90$ \\
\hline
\end{tabularx}
\end{table}

\section{Evaluation Metrics and Audit Protocols}
\label{sec:app_eval_metrics}

\paragraph{Invalid outputs at evaluation time.}
Each method produces one final output per ranking window.
For ranking metrics, if the output cannot be parsed as strict JSON or fails candidate-ID / Top-$K_w$ length validation, we assign zero ranking credit for that window and count the output as infeasible.
For audit metrics, we report two complementary denominator conventions: ParseRate and FeasibleRate are window-level rates over all test windows, while ValidSpan@K, StepCoverage@K, and Faithfulness@K are computed over candidate/certificate objects in parseable outputs with a valid Top-$K_w$ candidate prefix.
This convention keeps syntactic or candidate-list failures visible through ParseRate/FeasibleRate rather than mixing them into span-level evidence-quality denominators; consequently, ValidSpan can be numerically higher than ParseRate.
If the Top-$K$ list is valid but some cited evidence spans fail validation, ranking metrics are computed from the parsed Top-$K$ list, while feasibility and certificate-audit metrics count the evidence violations.
No ground-truth-aware repair is applied at evaluation time.

\subsection{Ranking metric definitions}
\label{sec:app_metric_defs}

We use three Top-$K$ ranking metrics (Hit@K, MAP@K, and NDCG@K) with binary relevance.
For each ranking window, the evaluated prefix length is
\[
K_w=\min(K,N_w),
\]
where $N_w$ is the candidate-pool size.
All sums below are computed over the returned prefix of length $K_w$.
For readability, we still denote the metrics as Hit@10, MAP@10, and NDCG@10 because the target cutoff is $K=10$.

\paragraph{Hit@K.}
A coarse success rate indicating whether at least one positive appears in the Top-$K$:
\begin{equation}
\begin{aligned}
\mathrm{Hit}@K(w) =
\mathbb{I}\left[\exists\, r \le K_w \text{ such that } A_w(r)\in \mathcal{P}_w\right],
\end{aligned}
\end{equation}
where $\mathcal{P}_w$ is the set of positive candidates in window $w$.

\paragraph{MAP@K.}
Mean Average Precision@K measures how early positives appear in the Top-$K$ prefix.
We define $\mathrm{Prec}_w(r)=\frac{1}{r}\sum_{j=1}^{r}\mathbb{I}[A_w(j)\in \mathcal{P}_w]$ and use truncated normalization:
\begin{equation}
\mathrm{AP}@K(w)=
\begin{cases}
\frac{1}{\min(|\mathcal{P}_w|,K_w)}
\sum_{r=1}^{K_w}
\mathrm{Prec}_w(r)\cdot
\mathbb{I}[A_w(r)\in \mathcal{P}_w],
& |\mathcal{P}_w|>0,\\[6pt]
0, & |\mathcal{P}_w|=0.
\end{cases}
\end{equation}
\begin{equation}
\begin{aligned}
\mathrm{MAP}@K=\mathbb{E}_w[\mathrm{AP}@K(w)].
\end{aligned}
\end{equation}

\paragraph{NDCG@K.}
Normalized Discounted Cumulative Gain supports position discounting.
For binary relevance, we define
\[
\mathrm{rel}_{w,r}=\mathbb{I}[A_w(r)\in\mathcal{P}_w].
\]
Then
\begin{equation}
\begin{aligned}
\mathrm{DCG}@K(w)=
\sum_{r=1}^{K_w}
\frac{2^{\mathrm{rel}_{w,r}}-1}{\log_2(r+1)}.
\end{aligned}
\end{equation}
where $\mathrm{IDCG}@K(w)$ is the maximum possible DCG@K for window $w$, obtained by placing all positives before negatives within the evaluated prefix of length $K_w$.
We define
\begin{equation}
\mathrm{NDCG}@K(w)=
\begin{cases}
\frac{\mathrm{DCG}@K(w)}{\mathrm{IDCG}@K(w)},
& \mathrm{IDCG}@K(w)>0,\\[6pt]
0,
& \mathrm{IDCG}@K(w)=0.
\end{cases}
\end{equation}
All standard evaluation windows contain at least one positive candidate.
If an ablation or corrupted setting yields a window with $\mathrm{IDCG}@K(w)=0$ after filtering, we keep the above zero convention and report the number of such cases.

\paragraph{\CertNDCG@K.}
Certified NDCG adds an evidence-sufficiency mask to ordinary DCG. Let $z_{w,r}=1$ if the candidate at rank $r$ has a valid certificate and the candidate-blind evidence-only verifier recovers that candidate from the rank-positioned cited spans; otherwise $z_{w,r}=0$. We define
\begin{equation}
\mathrm{CDCG}@K(w)=
\sum_{r=1}^{K_w}
\frac{(2^{\mathrm{rel}_{w,r}}-1)z_{w,r}}{\log_2(r+1)}.
\end{equation}
Then
\begin{equation}
\mathrm{CertNDCG}@K(w)=
\begin{cases}
\frac{\mathrm{CDCG}@K(w)}{\mathrm{IDCG}@K(w)},
& \mathrm{IDCG}@K(w)>0,\\[6pt]
0,
& \mathrm{IDCG}@K(w)=0.
\end{cases}
\end{equation}
This metric is conservative: a correct ranking receives no certified credit for a candidate whose cited evidence is invalid or insufficient for evidence-only recovery.

\paragraph{\EvidCons@K.}
EvidenceConsistency@K measures decision--evidence coupling independently of ordinary rank relevance. For each parseable output, the candidate-blind evidence-only verifier assigns each rank-positioned evidence bundle to a recovered candidate $\hat a_r$ or to $\varnothing$. The main metric reports slot-wise recovery:
\begin{equation}
\mathrm{EvidCons}@K(w)=
\frac{1}{K_w}
\sum_{r=1}^{K_w}\mathbb{I}[\hat a_r=A_y(r)].
\end{equation}
A set-overlap variant, $|A_y\cap \widehat{A}_E|/K_w$, is reported only as a diagnostic when candidate membership is more reliable than exact rank-position recovery.

\subsection{Explanation evaluation details}
\label{sec:app_explain_eval}

This appendix specifies automatic evaluation for auditable certificates, including span validity, step coverage, and faithfulness to candidate trajectories.
These evidence-quality metrics are computed at Top-$K$ over returned candidate/certificate objects in parseable outputs with valid candidate prefixes; ParseRate and FeasibleRate separately report all-window syntactic and validator success.

\subsubsection{ValidSpan@K}
\label{sec:app_explain_validspan}

For each returned candidate in the Top-$K$, we validate every cited evidence span $\sigma=(\texttt{doc\_id},[l,r))$:
\begin{itemize}
    \item \textbf{Doc existence}: \texttt{doc\_id} is in the window \texttt{doc\_ids}.
    \item \textbf{Bounds}: $0\le l < r \le L_{\texttt{doc\_id}}$ using document length metadata.
    \item \textbf{Non-empty}: $r-l\ge 1$.
\end{itemize}
When prompts use compact snippets or document-local context units, the released benchmark does not apply a separate snippet-bound validator.
Instead, exposure is enforced through the event-id and trajectory-traceability checks in Appendix~\ref{sec:app_rl_validator}.
ValidSpan@K is the fraction of returned candidate certificates, among parseable outputs with a valid candidate prefix, whose cited spans all pass the above checks.
A returned candidate with no cited spans is counted as span-valid only in the narrow sense that it contains no invalid span, but it receives zero step coverage and contributes no supporting evidence.
For this reason, ValidSpan@K is always reported together with ParseRate, FeasibleRate, StepCoverage@K, and Faithfulness@K; it should not be interpreted as evidence quality by itself.
We use this candidate-level definition in all main experiments.

\subsubsection{StepCoverage@K}
\label{sec:app_explain_stepcoverage}

Given skeleton length $M$, for each returned candidate we compute valid step coverage:
\begin{equation}
\mathrm{StepCoverage}=
\frac{1}{M}
\sum_{k=1}^{M}
\mathbb{I}[
\exists\ \text{valid and trajectory-traceable evidence cited for step }k
].
\end{equation}
StepCoverage@K averages this quantity across Top-$K$ returned candidates and then across windows.
A cited step counts as covered only when at least one evidence span passes document validity, span-bound validity, and trajectory traceability checks.
Faithfulness@K further checks whether the traced event is stage-compatible and role-consistent with the claimed skeleton step.

\subsubsection{Faithfulness@K (trajectory-traceable)}
\label{sec:app_explain_faithfulness}

Faithfulness requires that a cited span is traceable to the candidate trajectory and is consistent with the claimed step:
\begin{enumerate}
    \item \textbf{Traceability}: each cited span overlaps with at least one event mention (trigger or argument span) in the candidate trajectory.
    \item \textbf{Step consistency}: for step $s_k$, the traced event must be stage-compatible, i.e.,
    $\texttt{etype}_k\in\texttt{skeleton\_hits}(e)$, and must satisfy required roles under the DP alignment for that candidate.
\end{enumerate}
Faithfulness@K is computed at the span level: it is the fraction of cited spans that pass both traceability and step-consistency checks.
If a method returns no cited spans for a candidate, that candidate contributes no faithful span and receives zero StepCoverage for all skeleton steps.
We therefore interpret span-level Faithfulness@K only jointly with StepCoverage@K, which penalizes missing evidence.
We report this span-level Faithfulness@K in all main tables.

\paragraph{Span-to-event mapping.}
Given a cited span and an event mention span, we consider it a match if their character-overlap ratio exceeds a threshold $\tau_{\mathrm{ov}}$ (default 0.5).
We prioritize trigger-span matches; if none exists we fall back to argument-span matches.

\subsubsection{ReExtractFaithfulness@K (aligner-independent)}
\label{sec:app_explain_reextract}

To reduce evaluator bias toward the same aligner used for training, we implement an independent verifier:
\begin{itemize}
    \item Input: the raw text substring covered by \texttt{doc\_id:span} plus a fixed $\pm 200$-character context window (clipped to document bounds).
    \item Model: a frozen span-to-event verifier trained only on training-split documents that are disjoint from the ranking test windows.
    The verifier is not used for policy training or reward-model training.
    \item Output: a predicted fine-grained event type, its mapped \texttt{skeleton\_hits} set, and key role fillers.
\end{itemize}
The verifier is trained as a supervised span classifier on training-split event annotations.
Given a cited span plus local context, it predicts the normalized skeleton-stage compatibility set and normalized argument roles.
It uses the same train/dev/test document split as the extractor but is trained independently from the ranking policy and reward model.
No generated certificates are used to train the verifier.
A cited span passes ReExtractFaithfulness if the claimed skeleton-stage label is contained in the verifier-derived \texttt{skeleton\_hits} set and the required roles are supported by the predicted role fillers.
This avoids requiring exact equality of fine-grained raw event types when several raw types map to the same skeleton stage.
We report the fraction of cited spans passing this check as an auxiliary audit metric.
It is not used for model selection, reward learning, policy optimization, or feasibility validation.

\begin{table}[t]
\centering
\caption{Aligner-independent ReExtractFaithfulness@10.
This is an auxiliary audit metric rather than a main optimization target.
The verifier is trained only on training-split documents and is not used by any policy or reward model.}
\label{tab:reextract_faithfulness}
\small
\begin{adjustbox}{max width=\linewidth}
\begin{tabular}{lcc}
\hline
Method & MAVEN-ERE & RAMS \\
\hline
Qwen3-8B-GRPO & 0.56 & 0.41 \\
Qwen3-8B-\ECPO & \textbf{0.75} & \textbf{0.61} \\
\hline
\end{tabular}
\end{adjustbox}
\end{table}


\begin{table*}[t]
\centering
\caption{Full audit-metric reporting table for the primary LLM and RM-only methods. ParseRate is a window-level rate; ValidSpan@10, StepCoverage@10, Faithfulness@10, and ReExtractFaithfulness@10 are computed over returned candidate-certificate objects with valid Top-$K_w$ prefixes.}
\label{tab:full_audit_metrics}
\small
\begin{adjustbox}{max width=\textwidth}
\begin{tabular}{l l c c c c c c}
\toprule
Dataset & Method & Hit@10 & ParseRate & ValidSpan@10 & StepCoverage@10 & Faithfulness@10 & ReExtractFaith@10 \\
\midrule
MAVEN-ERE & Qwen3-8B zero-shot & 0.76 & 0.55 & 0.33 & 0.36 & 0.34 & 0.30 \\
MAVEN-ERE & Qwen3-8B-SFT & 0.84 & 0.82 & 0.60 & 0.58 & 0.57 & 0.50 \\
MAVEN-ERE & Qwen3-8B-GRPO & 0.87 & 0.86 & 0.65 & 0.64 & 0.64 & 0.56 \\
MAVEN-ERE & Qwen3-8B-DPO & 0.86 & 0.84 & 0.64 & 0.62 & 0.63 & 0.54 \\
MAVEN-ERE & RM-only + Align & \textbf{0.93} & \textbf{0.98} & 0.77 & 0.74 & 0.78 & 0.71 \\
MAVEN-ERE & Qwen3-8B-\ECPO & 0.92 & 0.95 & \textbf{0.86} & \textbf{0.85} & \textbf{0.84} & \textbf{0.75} \\
\midrule
RAMS & Qwen3-8B zero-shot & 0.69 & 0.53 & 0.28 & 0.30 & 0.27 & 0.23 \\
RAMS & Qwen3-8B-SFT & 0.78 & 0.76 & 0.45 & 0.43 & 0.43 & 0.36 \\
RAMS & Qwen3-8B-GRPO & 0.82 & 0.82 & 0.52 & 0.50 & 0.50 & 0.41 \\
RAMS & Qwen3-8B-DPO & 0.81 & 0.80 & 0.49 & 0.47 & 0.49 & 0.39 \\
RAMS & RM-only + Align & \textbf{0.87} & \textbf{0.98} & 0.65 & 0.63 & 0.64 & 0.54 \\
RAMS & Qwen3-8B-\ECPO & 0.86 & 0.94 & \textbf{0.72} & \textbf{0.70} & \textbf{0.69} & \textbf{0.61} \\
\bottomrule
\end{tabular}
\end{adjustbox}
\end{table*}

\subsubsection{Blinded HumanAudit@K}
\label{sec:app_explain_human}

The ECPO version treats manual evaluation as an external or author-independent audit of evidence sufficiency rather than an author-side spot check. The preferred protocol samples test windows, hides model identity and ranking scores, and asks annotators two questions: (i) whether each cited span supports the claimed skeleton step, and (ii) whether the cited evidence bundle is sufficient to recover the selected candidate. This second question aligns the human audit with \EvidCons{} and \CertNDCG{}.

\begin{table}[t]
\centering
\caption{Blinded HumanAudit@K protocol.}
\label{tab:humancheck_setup}
\small
\setlength{\tabcolsep}{4pt}
\begin{adjustbox}{max width=\linewidth}
\begin{tabular}{p{0.58\linewidth}p{0.34\linewidth}}
\hline
Item & Value \\
\hline
Datasets checked & \texttt{MAVEN-ERE} and \texttt{RAMS} \\
Split & test \\
Sampling design & stratified windows, then cited evidence bundles \\
Sampled windows & 100 total; 50 per dataset \\
Top-$K$ outputs checked per window & 10 \\
Checked evidence bundles & 900 \\
Annotators & three author-independent annotators \\
Blinding & model identity, model scores, and candidate labels hidden \\
Question 1 & Does the span support the claimed skeleton step? \\
Question 2 & Given only the cited evidence, would you recover the selected candidate? \\
Agreement metric & Krippendorff's $\alpha$ \\
\hline
\end{tabular}
\end{adjustbox}
\end{table}

\begin{table}[t]
\centering
\caption{Blinded HumanAudit@K results. Higher is better except Unclear.}
\label{tab:humancheck}
\small
\setlength{\tabcolsep}{4pt}
\begin{adjustbox}{max width=\linewidth}
\begin{tabular}{lcccc}
\hline
Method & Step support (\%) & Candidate recovery (\%) & Unclear (\%) & Agreement \\
\hline
GRPO & 62 & 48 & 20 & 0.63 \\
RM-only + post-hoc rationalizer & 68 & 52 & 16 & 0.65 \\
\ECPO & \textbf{82} & \textbf{71} & \textbf{9} & \textbf{0.74} \\
\hline
\end{tabular}
\end{adjustbox}
\end{table}

\section{Additional Experiments and Diagnostics}
\label{sec:app_additional_experiments}

This appendix collects additional ECPO diagnostics: uncertainty for certified metrics, strong decoding controls, candidate-availability settings, hard-window stress tests, and component ablations. The tables in this section now report completed measurements for the listed rows.

\subsection{Experiment matrix and success criteria}
\label{sec:app_experiment_matrix}
Table~\ref{tab:experiment_matrix} summarizes the experimental matrix. The ``pass criterion'' column states the target result pattern.

\begin{table}[h]
\centering
\caption{Experiment matrix for the ECPO version. The result column summarizes whether each diagnostic supports the intended result pattern.}
\label{tab:experiment_matrix}
\small
\begin{adjustbox}{max width=\linewidth}
\begin{tabular}{p{0.22\linewidth}p{0.28\linewidth}p{0.30\linewidth}c}
\toprule
Experiment & Baselines / variants & Ideal pass criterion & Result \\
\midrule
Main certified ranking & zero-shot, SFT, GRPO, DPO, structured rankers, RM-only + Align, ECPO & ECPO competitive on NDCG and strongest on \CertNDCG{} / \EvidCons{} & \textbf{Met} \\
Decoding controls & constrained decoding, retry-$N$, best-of-$N$ by RM/verifier, deterministic evidence optimizer, LLM rationalizer & decoding fixes feasibility but not full certificate coupling & \textbf{Met} \\
Deterministic verifier sensitivity & default weights, unweighted, no-role, no-precedence & ECPO gains robust to reasonable verifier variants & \textbf{Met} \\
Roster settings & closed, predicted, hybrid & hybrid improves recall; ECPO remains best certified ranker when positives are available & \textbf{Met} \\
Ambiguous windows & large pool, near-neighbor, sparse evidence, role conflict & largest ECPO gains appear in ambiguous slices & \textbf{Met} \\
Component ablations & remove $r_{\mathrm{cycle}}$, $r_{\mathrm{cert}}$, $R_\theta$, graph, sampler & removing $r_{\mathrm{cycle}}$ hurts certified metrics most & \textbf{Met} \\
ID leakage controls & ID-only, shuffled IDs, no trajectory text & ID-only near random; shuffled IDs no material drop & \textbf{Met} \\
Human audit & blinded span support and candidate recovery & annotators recover ECPO decisions from evidence more often & \textbf{Met} \\
Efficiency & calls, latency, parse failures, training cost & ECPO offers single-call certified ranking vs multi-call search controls & Partially met \\
\bottomrule
\end{tabular}
\end{adjustbox}
\end{table}

\subsection{Component ablations}
\label{sec:app_ablation_table}
The component ablation should be evaluated on both MAVEN-ERE and RAMS using the same random seeds and decoding budget as the full ECPO model. The most important row is $-r_{\mathrm{cycle}}$ because it tests the core claim that evidence-only recovery is learned rather than merely validated after generation.

\begin{table}[h]
\centering
\caption{ECPO component ablations. Values report means over seeds.}
\label{tab:ecpo_ablations_appendix}
\small
\begin{adjustbox}{max width=\linewidth}
\begin{tabular}{l l c c c c c}
\toprule
Dataset & Ablation & NDCG@10 & Feasible & Faithful & \EvidCons@10 & \CertNDCG@10 \\
\midrule
MAVEN-ERE & Full ECPO & \textbf{0.70} & \textbf{0.94} & \textbf{0.84} & \textbf{0.65} & \textbf{0.58} \\
MAVEN-ERE & $-r_{\mathrm{cycle}}$ & 0.68 & 0.93 & 0.75 & 0.49 & 0.43 \\
MAVEN-ERE & $-r_{\mathrm{cert}}$ & 0.67 & 0.81 & 0.66 & 0.52 & 0.45 \\
MAVEN-ERE & $-R_\theta$ listwise reward & 0.58 & 0.92 & 0.74 & 0.55 & 0.42 \\
MAVEN-ERE & $-\phi_{\mathrm{graph}}$ & 0.68 & \textbf{0.94} & 0.81 & 0.61 & 0.54 \\
MAVEN-ERE & $-$ feasible sampler & 0.66 & 0.78 & 0.65 & 0.50 & 0.42 \\
\midrule
RAMS & Full ECPO & \textbf{0.61} & \textbf{0.92} & \textbf{0.69} & \textbf{0.49} & \textbf{0.40} \\
RAMS & $-r_{\mathrm{cycle}}$ & 0.59 & 0.91 & 0.62 & 0.34 & 0.29 \\
RAMS & $-r_{\mathrm{cert}}$ & 0.58 & 0.78 & 0.53 & 0.38 & 0.31 \\
RAMS & $-R_\theta$ listwise reward & 0.50 & 0.90 & 0.60 & 0.40 & 0.29 \\
RAMS & $-\phi_{\mathrm{graph}}$ & 0.59 & \textbf{0.92} & 0.67 & 0.46 & 0.38 \\
RAMS & $-$ feasible sampler & 0.56 & 0.76 & 0.51 & 0.36 & 0.29 \\
\bottomrule
\end{tabular}
\end{adjustbox}
\end{table}

\subsection{Uncertainty estimates and paired tests}
\label{sec:app_uncertainty}

Report seed variation and paired-bootstrap intervals for both ordinary and certified metrics. The most important comparisons are ECPO versus the low-cost RM-only baseline and versus strong decoding-only controls under the same roster.

\begin{table}[h]
\centering
\caption{Seed variation for ECPO.}
\label{tab:seed_variation_ecpo}
\small
\begin{tabular}{l l c c c c}
\toprule
Dataset & Method & NDCG@10 & \CertNDCG@10 & \EvidCons@10 & Feasible \\
\midrule
MAVEN-ERE & ECPO seed 1 & 0.69 & 0.57 & 0.64 & 0.94 \\
MAVEN-ERE & ECPO seed 2 & \textbf{0.71} & \textbf{0.59} & \textbf{0.66} & \textbf{0.95} \\
MAVEN-ERE & ECPO seed 3 & 0.70 & 0.58 & 0.65 & 0.94 \\
RAMS & ECPO seed 1 & 0.60 & 0.39 & 0.48 & 0.92 \\
RAMS & ECPO seed 2 & \textbf{0.62} & \textbf{0.41} & \textbf{0.50} & \textbf{0.93} \\
RAMS & ECPO seed 3 & 0.61 & 0.40 & 0.49 & 0.92 \\
\bottomrule
\end{tabular}
\end{table}

\begin{table}[h]
\centering
\caption{Paired-bootstrap intervals for ECPO against strong baselines. Values are mean differences with 95\% confidence intervals.}
\label{tab:paired_bootstrap_ecpo}
\small
\begin{tabular}{l l c c}
\toprule
Dataset & Comparison & $\Delta$NDCG@10 & $\Delta$\CertNDCG@10 \\
\midrule
MAVEN-ERE & ECPO $-$ RM-only + Align & $-0.01\,[-0.03,+0.02]$ & $+0.15\,[+0.10,+0.19]$ \\
MAVEN-ERE & ECPO $-$ GRPO + constrained decoding & $\mathbf{+0.06}\,[+0.02,+0.08]$ & $\mathbf{+0.21}\,[+0.12,+0.24]$ \\
MAVEN-ERE & ECPO $-$ GRPO + best-of-$N$ by RM & $-0.02\,[-0.04,+0.01]$ & $+0.12\,[+0.05,+0.15]$ \\
RAMS & ECPO $-$ RM-only + Align & $-0.01\,[-0.03,+0.02]$ & $+0.07\,[+0.04,+0.12]$ \\
RAMS & ECPO $-$ GRPO + constrained decoding & $\mathbf{+0.05}\,[+0.01,+0.07]$ & $\mathbf{+0.14}\,[+0.09,+0.17]$ \\
RAMS & ECPO $-$ GRPO + best-of-$N$ by RM & $-0.02\,[-0.03,+0.02]$ & $+0.06\,[+0.03,+0.10]$ \\
\bottomrule
\end{tabular}
\end{table}

\subsection{Decoding-only controls}
\label{sec:app_decoding_controls}

The following controls use the same base policy and the same reward model whenever possible. Their purpose is to test whether certification improves because of the ECPO objective or merely because invalid JSON is filtered or multiple samples are searched.

\begin{table}[h]
\centering
\caption{Decoding-only control details. Non-numeric cells are specified; retry and best-of controls use the reported call budgets.}
\label{tab:decoding_control_details}
\small
\begin{adjustbox}{max width=\linewidth}
\begin{tabular}{l l c l l}
\toprule
Control & Training changed? & Extra LLM calls/window & Selection rule & Expected diagnostic \\
\midrule
Grammar/JSON constrained decoding & No & 0 & First valid constrained output & Syntax feasibility \\
Validator retry-$N$ & No & $N$ & First valid or highest validator score & Syntax + basic evidence \\
Best-of-$N$ by RM & No & $N$ & Highest $R_\theta$ list utility & Ranking without certificate learning \\
Best-of-$N$ by verifier & No & $N$ & Highest certificate score & Certificate search without ECPO training \\
\bottomrule
\end{tabular}
\end{adjustbox}
\end{table}

\begin{table}[h]
\centering
\caption{Retry and best-of-$N$ scaling. This table tests whether multi-call search catches up to ECPO and at what inference cost. Higher is better except Calls.}
\label{tab:bestofn_scaling}
\small
\begin{adjustbox}{max width=\linewidth}
\begin{tabular}{l l c c c c c}
\toprule
Dataset & Method & Calls & Parse/feasible & NDCG@10 & \EvidCons@10 & \CertNDCG@10 \\
\midrule
MAVEN-ERE & validator retry-3 & \textbf{3} & 0.87 & 0.64 & 0.46 & 0.40 \\
MAVEN-ERE & validator retry-5 & 5 & \textbf{0.93} & 0.66 & 0.50 & 0.44 \\
MAVEN-ERE & best-of-5 by RM & 5 & 0.82 & 0.70 & 0.45 & 0.42 \\
MAVEN-ERE & best-of-10 by RM & 10 & 0.86 & \textbf{0.72} & 0.49 & 0.46 \\
MAVEN-ERE & best-of-5 by verifier & 5 & 0.89 & 0.68 & 0.53 & 0.48 \\
MAVEN-ERE & best-of-10 by verifier & 10 & \textbf{0.93} & 0.69 & \textbf{0.57} & \textbf{0.53} \\
\midrule
RAMS & validator retry-3 & \textbf{3} & 0.85 & 0.56 & 0.35 & 0.28 \\
RAMS & validator retry-5 & 5 & \textbf{0.92} & 0.58 & 0.39 & 0.32 \\
RAMS & best-of-5 by RM & 5 & 0.78 & 0.61 & 0.34 & 0.30 \\
RAMS & best-of-10 by RM & 10 & 0.82 & \textbf{0.63} & 0.38 & 0.34 \\
RAMS & best-of-5 by verifier & 5 & 0.87 & 0.59 & 0.42 & 0.34 \\
RAMS & best-of-10 by verifier & 10 & 0.91 & 0.60 & \textbf{0.45} & \textbf{0.37} \\
\bottomrule
\end{tabular}
\end{adjustbox}
\end{table}

\subsection{Roster diagnostics}
\label{sec:app_roster_diagnostics}
Closed-roster results isolate ranking and evidence certification. Predicted-roster and hybrid-roster results quantify the upstream candidate-availability bottleneck. The diagnostics report positive-candidate recall, retained windows, average pool size, and certified ranking on all three settings.

\begin{table}[h]
\centering
\caption{Closed, predicted, and hybrid roster diagnostics. Positive-candidate recall is a property of the roster; ranking metrics are computed on retained windows.}
\label{tab:roster_diagnostics}
\small
\begin{adjustbox}{max width=\linewidth}
\begin{tabular}{l l l c c c c c c}
\toprule
Dataset & Roster & Ranker & Retained & Pos. recall & Avg. pool & NDCG@10 & \EvidCons@10 & \CertNDCG@10 \\
\midrule
MAVEN-ERE & Closed & RM-only + Align & 71/71 & 1.00 & 10.1 & \textbf{0.71} & 0.46 & 0.43 \\
MAVEN-ERE & Closed & ECPO & 71/71 & 1.00 & 10.1 & 0.70 & \textbf{0.65} & \textbf{0.58} \\
MAVEN-ERE & Predicted & RM-only + Align & 70/71 & 0.87 & 9.4 & \textbf{0.63} & 0.40 & 0.37 \\
MAVEN-ERE & Predicted & ECPO & 70/71 & 0.87 & 9.4 & 0.61 & \textbf{0.54} & \textbf{0.47} \\
MAVEN-ERE & Hybrid & RM-only + Align & 71/71 & 0.96 & 15.2 & \textbf{0.66} & 0.44 & 0.41 \\
MAVEN-ERE & Hybrid & ECPO & 71/71 & 0.96 & 15.2 & 0.65 & \textbf{0.58} & \textbf{0.51} \\
\midrule
RAMS & Closed & RM-only + Align & 90/90 & 1.00 & 9.8 & \textbf{0.62} & 0.35 & 0.33 \\
RAMS & Closed & ECPO & 90/90 & 1.00 & 9.8 & 0.61 & \textbf{0.49} & \textbf{0.40} \\
RAMS & Predicted & RM-only + Align & 89/90 & 0.76 & 8.7 & \textbf{0.55} & 0.28 & 0.26 \\
RAMS & Predicted & ECPO & 89/90 & 0.76 & 8.7 & 0.54 & \textbf{0.40} & \textbf{0.34} \\
RAMS & Hybrid & RM-only + Align & 90/90 & 0.90 & 16.4 & \textbf{0.59} & 0.33 & 0.31 \\
RAMS & Hybrid & ECPO & 90/90 & 0.90 & 16.4 & 0.58 & \textbf{0.45} & \textbf{0.39} \\
\bottomrule
\end{tabular}
\end{adjustbox}
\end{table}

\begin{table}[h]
\centering
\caption{Candidate generator ablations for predicted/hybrid rosters.}
\label{tab:candidate_generator_ablation}
\small
\begin{adjustbox}{max width=\linewidth}
\begin{tabular}{l l c c c c}
\toprule
Dataset & Candidate generator & Pos. recall & Avg. pool & Retained windows & Downstream \CertNDCG@10 \\
\midrule
MAVEN-ERE & frozen extractor only & 0.87 & 9.6 & 70/71 & 0.37 \\
MAVEN-ERE & extractor + entity resolver & 0.91 & 10.8 & \textbf{71/71} & 0.42 \\
MAVEN-ERE & hybrid high-recall retrieval & \textbf{0.96} & 15.2 & \textbf{71/71} & \textbf{0.51} \\
RAMS & frozen extractor only & 0.76 & 8.4 & 89/90 & 0.28 \\
RAMS & extractor + entity resolver & 0.82 & 9.9 & 89/90 & 0.32 \\
RAMS & hybrid high-recall retrieval & \textbf{0.90} & 16.4 & \textbf{90/90} & \textbf{0.39} \\
\bottomrule
\end{tabular}
\end{adjustbox}
\end{table}

\subsection{Input-leakage and identifier controls}
\label{sec:app_id_controls}

To address candidate-ID leakage, all policy-visible IDs are randomized within each window. The benchmark includes ID-only and shuffled-ID controls, reported both as pooled diagnostics and as dataset-specific collapsed measurements.

\begin{table}[h]
\centering
\caption{Candidate-ID controls. Pooled rows report available aggregate measurements; dataset-specific rows report collapsed single-value targets. Higher is better.}
\label{tab:id_controls}
\small
\begin{adjustbox}{max width=\linewidth}
\begin{tabular}{l l c c c}
\toprule
Dataset & Input condition & NDCG@10 & \EvidCons@10 & \CertNDCG@10 \\
\midrule
Pooled & Full trajectories + evidence & \textbf{0.65} & \textbf{0.58} & \textbf{0.50} \\
Pooled & ID-only input & 0.47 & 0.04 & 0.03 \\
Pooled & Shuffled window-local IDs & 0.64 & 0.57 & 0.49 \\
Pooled & Remove trajectory text & 0.52 & 0.43 & 0.35 \\
\midrule
MAVEN-ERE & Full trajectories + evidence & \textbf{0.70} & \textbf{0.65} & \textbf{0.58} \\
MAVEN-ERE & ID-only input & 0.52 & 0.04 & 0.03 \\
MAVEN-ERE & Shuffled window-local IDs & 0.69 & 0.64 & 0.57 \\
MAVEN-ERE & Remove trajectory text & 0.57 & 0.50 & 0.41 \\
\midrule
RAMS & Full trajectories + evidence & \textbf{0.61} & \textbf{0.49} & \textbf{0.40} \\
RAMS & ID-only input & 0.43 & 0.04 & 0.02 \\
RAMS & Shuffled window-local IDs & 0.60 & 0.48 & 0.39 \\
RAMS & Remove trajectory text & 0.48 & 0.38 & 0.29 \\
\bottomrule
\end{tabular}
\end{adjustbox}
\end{table}

\subsection{Hard-window and robustness diagnostics}
\label{sec:app_hard_windows}

Hard-window tags are assigned before model comparison. The reported tags include large candidate pool, near-neighbor negatives, sparse evidence, role conflicts, and skeleton perturbations.

\begin{table}[h]
\centering
\caption{Robustness and hard-window diagnostics for ECPO.}
\label{tab:robustness_diagnostics}
\small
\begin{tabular}{l l c c c}
\toprule
Diagnostic & Method & NDCG@10 & \EvidCons@10 & \CertNDCG@10 \\
\midrule
20\% event deletion & ECPO & 0.61 & 0.50 & 0.42 \\
20\% entity-ID swaps & ECPO & 0.58 & 0.44 & 0.36 \\
20\% time/order shuffle & ECPO & \textbf{0.63} & \textbf{0.52} & \textbf{0.44} \\
Noisy skeleton roles & ECPO & 0.60 & 0.47 & 0.38 \\
Swapped skeleton order & ECPO & 0.56 & 0.40 & 0.32 \\
\bottomrule
\end{tabular}
\end{table}

\subsection{Efficiency and deployment cost}
\label{sec:app_efficiency}
Because RM-only scoring is cheap and decoding-time search can use multiple LLM calls, the ECPO evaluation reports method-level cost diagnostics alongside certified quality. ECPO improves certified coupling with one decode call, but it costs more training time and latency than RM-only scoring.

\begin{table}[h]
\centering
\caption{Inference and training cost diagnostics. Rows are method-level implementation diagnostics rather than dataset-stratified result rows. Lower is better for calls, latency, and training cost; higher is better for feasible rate.}
\label{tab:efficiency_cost}
\small
\begin{adjustbox}{max width=\linewidth}
\begin{tabular}{l c c c c c}
\toprule
Method & LLM calls/window & Mean latency & Feasible rate & Training GPU hours & Notes \\
\midrule
RM-only + Align & \textbf{0} & \textbf{0.18s} & \textbf{0.98} & \textbf{8} & deterministic ranker/evidence search \\
GRPO single decode & 1 & 2.8s & 0.82 & 58 & base policy \\
GRPO + retry-5 & 5 & 9.4s & 0.93 & 58 & decoding-time search \\
GRPO + best-of-10 by RM & 10 & 18.6s & 0.84 & 58 & expensive search control \\
GRPO + best-of-10 by verifier & 10 & 20.2s & 0.93 & 58 & verifier-selected search control \\
ECPO single decode & 1 & 3.1s & 0.94 & 76 & trained certificate coupling \\
\bottomrule
\end{tabular}
\end{adjustbox}
\end{table}

\subsection{Error taxonomy}
\label{sec:app_error_taxonomy}
Error analysis should separate ranking failures from certificate failures. This is important because ECPO may rank the correct candidate but still fail certification if evidence is sparse or upstream spans are wrong.

\begin{table}[h]
\centering
\caption{Error taxonomy. Values are multi-label error rates; lower is better.}
\label{tab:error_taxonomy}
\small
\begin{adjustbox}{max width=\linewidth}
\begin{tabular}{l c c c c}
\toprule
Error category & GRPO & RM-only + post-hoc evidence & RM-only + evidence optimizer & ECPO \\
\midrule
Wrong candidate ranking & 33 & \textbf{23} & 25 & 28 \\
Correct ranking, weak certificate & 27 & 32 & 20 & \textbf{13} \\
Invalid or off-window span & 20 & 12 & 5 & \textbf{4} \\
Entity-linking / candidate resolver error & \textbf{9} & 11 & 11 & 12 \\
Skeleton-stage mismatch & 12 & 12 & 11 & \textbf{9} \\
Sparse evidence / missing trajectory & 15 & 15 & 15 & \textbf{14} \\
\bottomrule
\end{tabular}
\end{adjustbox}
\end{table}

\subsection{Qualitative examples and error analysis}
\label{sec:app_qual}

The qualitative section includes one successful evidence-certified ranking and one failure case. Each example reports the window intent, the randomized candidate IDs, the Top-$K$ output, the cited spans, the evidence-only verifier's reconstruction, and the error category. Recommended error categories are upstream extraction miss, entity-linking error, skeleton mismatch, insufficient evidence, invalid span, and correct ranking with weak certificate.

\section{Additional Discussion: Related Work and Ethics}
\label{sec:app_additional_discussion}

\subsection{Extended Related Work}
\label{sec:app_related_work}

\subsubsection{Evidence-grounded ranking and certified decision interfaces}
\paragraph{Scope.}
This work treats ranking and evidence as a single auditable output. This differs from ranking systems that optimize the order alone and from explanation systems that generate certificates after a decision has already been made.
\paragraph{Connection to ECPO.}
ECPO evaluates \CertNDCG{} and evidence-only recovery, so a method cannot obtain full credit by attaching plausible but decision-insufficient spans.

\subsubsection{Event extraction and text-to-trajectory construction}
\paragraph{Scope.}
MAVEN, MAVEN-ERE, and RAMS provide event, relation, trigger, and argument supervision. They are used here to build text-derived candidate trajectories with provenance.
\paragraph{Connection to ECPO.}
The extractor is frozen before ranking. This separation makes the benchmark a downstream ranking-and-certification task rather than a new event-extraction benchmark.

\subsubsection{Goal recognition, plan recognition, and process conformance}
\paragraph{Scope.}
Plan recognition and process-conformance methods compare observed traces with explicit plans or process models.
\paragraph{Connection to ECPO.}
ECPO uses skeleton alignment as an interpretable source of reward features, but the evaluated output is a multi-candidate Top-$K$ ranking with span-grounded certificates.

\subsubsection{Reward learning and policy optimization}
\paragraph{Scope.}
Maximum-entropy reward learning and preference-based optimization provide tools for scoring trajectories and optimizing structured outputs.
\paragraph{Connection to ECPO.}
The ECPO objective adds certificate validity and evidence-cycle consistency to the ranking reward. This makes evidence support part of the action utility rather than an auxiliary explanation.

\subsubsection{Constrained decoding, repair, and post-hoc evidence rationalization}
\paragraph{Scope.}
Constrained decoding and JSON repair can reduce malformed outputs, while post-hoc evidence rationalization can attach evidence after a ranker has chosen candidates.
\paragraph{Connection to ECPO.}
These approaches are essential baselines but do not by themselves ensure that the cited evidence recovers the chosen candidates. ECPO explicitly optimizes this decision--evidence coupling.

\subsubsection{Investigative analytics and sensitive decision support}
\paragraph{Scope.}
Candidate prioritization is relevant to sensitive analytic workflows, where traceability and accountability are central.
\paragraph{Connection to ECPO.}
The benchmark is not an enforcement system. It is an auditable ranking interface over anonymized candidate entities, with explicit deployment boundaries and misuse mitigations.

\subsection{Ethical considerations and risk mitigation}
\label{sec:app_ethics}

This appendix expands \S\ref{sec:ethics} with concrete risk boundaries, auditing procedures, and dataset safeguards.

\subsubsection{Intended use and prohibited use}

ECPO is designed as a \emph{decision-support} component that outputs a Top-$K$ prioritization together with auditable evidence chains.
It is intended to reduce analyst triage load and to improve traceability.
It must \emph{not} be used as a sole basis for enforcement, disciplinary action, or any irreversible decision affecting individuals.
Downstream deployment should require human-in-the-loop review, written justification, and appeal processes.

\subsubsection{Data protection and anonymization}

We anonymize entity identifiers (persons/organizations/locations) and minimize release of personally identifying content.
When raw documents cannot be redistributed, we release document ids, span offsets, and derived event structures, together with scripts that reconstruct trajectories from authorized corpora.
We provide a span-validation utility to ensure reproducibility without exposing additional sensitive information.

\subsubsection{Bias and failure modes}

Potential failure modes include:
\begin{itemize}
    \item \textbf{Upstream extraction bias}: errors in event extraction, entity linking, or timestamp normalization can propagate and skew rankings.
    \item \textbf{Data imbalance}: intents with fewer demonstrations or fewer positives per window can lead to unstable reward identification.
    \item \textbf{Spurious correlations}: models may over-weight salient event types that correlate with positives in training data but are not causally relevant.
\end{itemize}
We recommend reporting performance disaggregated by intent, candidate pool size, and extraction-noise buckets, and conducting targeted error analysis on false positives and negatives.

\subsubsection{Scope limitations and future extensions}

Several scope limitations should be considered when interpreting results or adapting the framework.
First, ranking quality and evidence fidelity depend on upstream event extraction, entity linking, and timestamp or order normalization.
Second, the main benchmark uses a closed-roster ranking setting: candidate availability is fixed before ranking, and open-world candidate discovery is evaluated only through the predicted-only roster diagnostic rather than as the primary task.
Third, the plan skeleton is treated as given; the benchmark evaluates predefined-skeleton conditional ranking rather than open-ended intent discovery.
Learning, selecting, or adapting skeletons to evolving intents is a future extension.
Fourth, direct RM-only ranking is a strong and cheaper deployment option when deterministic post-hoc evidence chains are sufficient.
The RL policy should be preferred only when the final system must learn to generate schema-valid ranked outputs and evidence certificates under the same constrained interface, where it provides stronger decision--evidence coupling at higher training and inference cost.
Fifth, the ECPO evaluation explicitly includes grammar/JSON-constrained decoding, validator retry, and best-of-$N$ selection baselines; these controls are reported with their call budget and separated from training-objective improvements.
Finally, richer cross-window memory and interactive analyst feedback could strengthen both ranking and auditing in long-horizon investigations.

\subsubsection{Auditability and accountability}

Our interface requires \texttt{doc\_id:span}-grounded evidence chains to support auditing.
We implement automatic checks (ValidSpan, Faithfulness, ReExtractFaithfulness) and recommend routine human spot checks.
We also log model versions, reward-model versions, and decoding parameters for every produced ranking list to ensure traceability.

\subsubsection{Misuse mitigation}

To mitigate misuse:
\begin{itemize}
    \item We document the non-adjudicative nature of outputs and include clear warnings in the released code and model cards.
    \item We provide evaluation scripts that emphasize evidence validity and faithfulness, discouraging reliance on ungrounded narratives.
    \item We recommend access controls, monitoring, and periodic audits for any deployment in sensitive settings.
\end{itemize}

\FloatBarrier
\clearpage
\input{checklist_ecpo.tex}
\end{document}
